\documentclass[12pt,a4paper]{article}

\usepackage[utf8]{inputenc}
\usepackage[T2A]{fontenc}
\usepackage[bulgarian]{babel}
\usepackage{graphicx}
\usepackage{geometry}
\usepackage{enumitem}
\usepackage[colorlinks=true,linkcolor=blue]{hyperref}

\geometry{margin=2cm}

\usepackage{cmap}
\geometry{top=2cm, bottom=2cm, left=2cm, right=2cm}

\usepackage{array}
\usepackage{tabularx}
\usepackage{setspace}
\usepackage{tikz}
\usepackage{color,soul}
\usepackage{xcolor}

\begin{document}
\begin{titlepage}
\centering

\vspace*{2cm}
{\large МИНИСТЕРСТВО НА ОБРАЗОВАНИЕТО И НАУКАТА\par}
\vspace{-0.2cm}
\noindent\rule{0.8\textwidth}{0.5pt}
\vspace{0.3cm}
{\large ПОДГОТОВКА ЗА ДЪРЖАВЕН ЗРЕЛОСТЕН ИЗПИТ\par}
\vspace{-0.4cm}
\noindent\rule{0.8\textwidth}{0.5pt}
\vspace{0.8cm}

\vspace{1.5cm}
{\Huge\bfseries БЪЛГАРСКИ ЕЗИК И ЛИТЕРАТУРА\par}
\vspace{0.6cm}
{\Large Подготовка за ДЗИ – 12. клас\par}
\vspace{0.6cm}
{\Large Сборник с тестове и примерни изпити\par}
\vspace{2.4cm}

\noindent\rule{1.0\textwidth}{0.5pt}
% \begin{center}
% \Huge $\ast\ \ast\ \ast$
% \end{center}
\vspace{0.4cm}
{\Large\bfseries Учебна година\par}
{\Large\bfseries 2026/2027 \par}
\vspace{-0.2cm}

\noindent\rule{1.0\textwidth}{0.5pt}
{\Large\bfseries Автор:\par}
{\large Катерина Димитрова\par}
\noindent\rule{1.0\textwidth}{0.5pt}
\vfill
{\large • София • \par}
{\large 2026 \par}
\end{titlepage}

\tableofcontents
\newpage

\section{11. клас}

\subsection{Родното и чуждото}
\subsubsection{\colorbox{green!40}{Димитър Талев – „Железният светилник“}}

\textbf{1. Социокултурен и литературен контекст}
\vspace{0.1cm}
\newline
\indent Творчеството на Димитър Талев e насочено към миналото на народа, като източник на гордост и самочувствие. Романите му отразяват различни етапи от родната история, като авторът се фокусира върху отношението между личността и обществото. Интересът е провокиран от това как пътят на индивида се преплита с този на колектива, за да бъдат обяснени понятия като дълг и чест.
\vspace{0.1cm}
\newline
\indent Романът “Железният светилник” e първата част от тетралогията на Талев, в която е проследен пътят на народа към свободата. Димитър Талев продължава романовата традиция на Иван Вазов и също като него разкрива промените в бита на българина под напора на историческите събития. И двамата автори се вглеждат в народопсихологията. На фона на голямо историческо платно е проследен и процесът на вграждането на индивидуалността в общия дом на народа.
\vspace{0.1cm}
\newline
\indent Романът отразява авторовите виждания за ролята на добродетелите и за патриархалната  устойчивост като гаранция за съхранението на народа. В текста е изобразен моментът, в който личността излиза от света на родното, себеопознава се и се устремява към свободата си.
\vspace{0.1cm}
\newline
\textbf{2. Жанрово определение}
\vspace{0.1cm}
\newline
\indent “Железният светилник” носи всички характеристики на романа като жанр. Това НЕ е исторически роман, тъй като в него има и моменти, разкриващи народопсихологията на възрожденския човек.
\vspace{0.1cm}
\newline
\textbf{3. Тълкуване на заглавието}
\vspace{0.1cm}
\newline
\indent Паратекстът е разгърната метафора на надеждата, която крепи човека, възпитан в патриархалните традиции на своя род. Светилникът въвежда образа на светлината като прогрес и прозрение, но и като алюзия за изгаряне. Определението „железният“ акцентира върху вещественото непреходното символизирал устойчивостта на традициите.
\vspace{0.1cm}
\newline
\textbf{4. Сюжет и композиция}
\vspace{-0.70cm}
\newline
\begin{itemize}
    \item Композиция
    \newline
    \hspace*{0.5cm} Четири части, всяка със свое заглавие и епиграф, част от народна песен, смислово насочваща към съдържанието.
    \begin{itemize}
        \item I част -"Хаджи Серафимовата внука"
        \item II част -"В тъмни времена"
        \item III част -"Народът се пробужда"
        \item IV част -"Корени и гранки"
    \end{itemize}
    \newline
    \hspace*{0.5cm} Всички части проследяват историята на рода Глаушеви и романът завършва със сватбата на Лазар и Ния.
\end{itemize}
\vspace{0.1cm}
\newline
\textbf{5. Персонажна система}
\vspace{-0.70cm}
\newline
\begin{itemize}
    \item Султана Глаушева
    \newline
    \hspace*{0.5cm} В нейният образ Димитър Талев представя родното в неговото интимно измерение. За Султана домът е въплъщение на най-ценното и непреходното в живота на човека. Неслучайно името ѝ насочва към владетелска титла, тъй като тя е истинският стожер на дома. Тя предопределя съдбата на близките си и очертава техният житейски път. Тя е първият герой нарушител в романа, защото престъпва установените норми, изправена пред заплахата да види загиването Хаджи Серафимовия род. Родното в представите на султана е познатото, консервативното, провереното и доказано в миналото като гаранция за оцеляване. Тя изповядва житейската философия за „прибирането“ в себе си. На следване на утъпкан път от другите. Чуждото за героинята е всичко, което не познава, възприема го като заплаха. Султана категорично отрича правото на индивида за себеизява като в патриархалното общество той може да оцелее само чрез връзката си с колектива.
    \vspace{-0.2cm}
    \item Лазар Глаушев
    \newline
   \hspace*{0.5cm} Името на героя препратка към библейският персонаж от Новия завет на Библията. Подобно на него, героят на Талев „възкръсва“ и се превръща във водач, апостол на поробените си събратя. Самият герой е причината съгражданите му да прогледнат, да преодолеят робското у себе си и да излязат от времето на тъмните времена на пътя на промяната. Той също е герой нарушител, тъй като престъпва майчината воля и се опълчва на духовните и политическите врагове на народа. Образът му е изграден чрез преплитане на мотивите за: СЛОВО, ЗНАНИЕ, СВЕТЛИНА (сравняваме с Левски). 
    \newline
    \hspace*{0.5cm} Повратният момент в живота на Лазар е срещата със Средногорския монах. Димитър Талев въплъщава в неговия образ идеята за народен будител. Лазар Глаушев оглавява борбата за независима българска църква, българско училище и читалище.
    \vspace{-0.2cm}
    \item Стоян Глаушев
    \newline
   \hspace*{0.5cm} Бягството му от село Гранче е спонтанен отпор срещу робството. Името му означава устойчивост, здравина, сила.
   \vspace{-0.2cm}
    \item Катерина Глаушева
    \newline
   \hspace*{0.5cm} Най-малката дъщеря на Глаушеви е въплъщение на самия живот (буйна непокорна, красива, истинска). Тя става герой нарушител, защото превръща любовта в своя религия: изпреварва времето си, следвайки поривите на сърцето, а не на разума. Прави свой собствен свободен избор, пренебрегвайки неписаните правила на поведение. Катерина става жертва на консервативното мислене.
   \vspace{-0.2cm}
    \item Рафе Клинче
    \newline
   \hspace*{0.5cm} Герой, въплъщение на ренесансовия човек, за което говори и името му. Той не се вписва в затворената гръцка общност, защото талантът му определя друг светоглед. Завинаги остава чужд на преспанци, въпреки че оставя след себе си иконостаса – Мъртво дърво, в което той въплъщава душата си. Остава като знак на вярата и любовта, защото в него е вграден образът на Катерина.
\end{itemize}
\vspace{0.1cm}
\newline
\textbf{6. Родното и чуждото в романа}
\vspace{0.1cm}
\newline
\indent “Железният светилник” се гради върху конфликтът родното и чуждото. Своето е представено най напред чрез образа на рода - затворена общност от близки, свои хора, познати лица, свързани с кръвни или родствени връзки.
\vspace{0.1cm}
\newline
\indent Приказно-легендарен образ на родното изгражда Рилският монах – възрожденски модел на връзки, основани върху ЕЗИК, ВЯРАТА, ПРОИЗХОД, ТЕРИТОРИЯ. Родното за него осмислено като свръхценност, затова трябва да бъде опазено и съхранено.
\vspace{0.1cm}
\newline
\indent Рафе Клинче дава друга представа за родното и чуждото. За него чуждото не е това, което НЕ познава, а това, което го застрашава.
\vspace{0.1cm}
\newline
\indent Чужди в романа не са само турците и гръцкото духовенство. Чужди за патриархалната общност са всички различни, които не се вписват в общоприетите представи.
\vspace{0.1cm}
\newline
\textbf{7. Железният светилник като символичен за общността предмет}
\vspace{0.1cm}
\newline
\indent Той символизира светлината най-напред във родовата общност. Стопля и огрява най-интимното българско пространство. Неговият пламък обаче и възможност труда да се вгледа в миналото си, да осъзнае силата на корените си, за да предопредели бъдещето си. Светлината присъства при всяко изпитание на общността - всяка малка победа е причина от просветление. В същото време светилникът е силата, устойчивостта, но и белег на отминалите времена. Затова в рамките на тетралогията, той постепенно е заменен от газената лампа. Родът не само трябва да се съхрани, но трябва и да се развива. Народът излиза от тъмните времена и поема по пътя на прогреса, но трябва да оцелеят корените на българския род, за да потекат живителните сокове на новото към младите гранки.
\newpage

\subsubsection{\colorbox{green!40}{Алеко Константинов – „Бай Ганьо журналист“}}
\textbf{1. Социокултурен и литературен контекст}
\vspace{0.1cm}
\newline
\indent Творчеството на Алеко Константинов отразява преломното време на 90-те години на XIX век, Когато в България се установяват новите икономически отношения и се формират партийни  фракции (партии или формирования), които се борят за електорат (привърженици). Възрожденските ценности остават назад в миналото и отстъпват място на личностни интереси, егоизъм, финансово облагодетелстване. Срещу моралната деградация на следосвобожденското общество се обявява европеецът по дух и образование Алеко Константинов (Пояснение: Завършва право в Одеса), в чиито творби като в криво огледало се отразяват пороците в новоосвободена България.
\vspace{0.1cm}
\newline
\indent Неуспял юрист, пътешественик и гражданин, Щастливеца не подминава нищо в обществения живот, което противоречи на представите му за дълг, чест и морал.
\vspace{0.1cm}
\newline
\indent Фейлетонът е част от книгата „Бай Ганьо“, където Алековият търговец на розово масло, завърнал се от пътуването си в Европа, се впуска в обществения живот в преследване на келепира. Текстът е изобличаване на продажната журналистика и политическото хамелеонство.
\vspace{0.1cm}
\newline
\textbf{2. Тълкуване на заглавието}
\vspace{0.1cm}
\newline 
\indent Героят е директно назован като основен обект на повествованието и носител на проблема. Посочената професия показва способността му да се адаптира и подготвя читателя за нарасналото му самочувствие да търси друго, по-широко, поле за изява.
\vspace{0.1cm}
\newline
\textbf{3. Жанр и композиция}
\vspace{0.1cm}
\newline 
\indent Творбата най често се определя като фейлетон, тъй като засяга злободневни и актуални за обществото проблеми. Носи характеристиките и на разказ и но очерк, тъй като героите са плод на авторовото въображение, но в творбата присъстват и правдоподобни личности и събития.
\vspace{0.1cm}
\newline 
\indent Героите на Алеко са ОБРАЗИ ТИПОВЕ, които са събирателни образи на конкретна социална група. Чрез тях авторът иронизира проявата на байганьовщината в следосвобожденската ни действителност. Иронията му е насочена към продажната журналистика, пошлостта, изопачаването на истината.
\vspace{0.1cm}
\newline 
\indent Използван е похватът РАЗКАЗ В РАЗКАЗА в композицията, като действията на главния герой са представени чрез друг разказвач.
\vspace{0.1cm}
\newline
\textbf{4. Образна система}
\vspace{0.1cm}
\newline 
\indent Героите в очерка са разделени на две групи, които представят двете лица на българската действителност. По-специфична и интересна за разглеждане това е групата на интелигентите (кръгът „Весела България“). Групата на разказвачите представя онова разбиране за родното и чуждото, което изповядва и самият писател. Според тях трябва да се заимства от постиженията на другите народи, НО да се съхраняват непреходните култура, обичаи, душевност, за да се запази онова, което ни прави различни от останалите.
\vspace{-0.7cm}
\newline 
\begin{itemize}
    \item Бай Ганьо
    \newline
    \hspace*{0.5cm} Героят не се е променил особено от първата част на книгата. Той е в разцвета на силите си, запазил високото си самочувствие, но си остава все така ограничен и прост. Пътуването е внесло промени във визията му, но това не значи, че Бай Ганьо се е култивирал (правим сравнение с „Вълкът козината си мени, но нрава – не.“). Той вече е придобил самочувствие на всезнаещ човек, което обяснява и амбициите му. Променена е визията на героя (костюм, вратовръзка, риза), но тези детайли не го превръщат в европеец – модерната визия не съответства на неговата духовна ограниченост, пошлост, егоизъм. Алековият персонаж въплъщава всички фалшиви стойности на своето време. Вестникът за него е начин НЕ да работи за народното благо, а за себе си. Продажбата на съвестта се оправдава с финансовите облаги. Разбирането му за журналистика се свежда до политическо хамелеонство. Словото е превърнато в подлога/слуга на политическите интереси. Името на вестника е параван, зад който са скрити истинските цели на издателя. Изданието НЕ е трибуна на истината и обективността. НЕ работи за народното благо. В „превод“ заглавието би трябвало да звучи по следният начин: „Всички в наша изгода.“.
    \vspace{-0.2cm}
    \item Дочоолу и Гочоолу
    \newline
    \hspace*{0.5cm} Не са реални личности, но чрез имената им е показана все още непреодоляната културна изостаналост на народа ни. Двамата се поддават на манипулациите на Бай Ганьо и техния проблем е, че се поддават и на ширещият се масов материализъм.
    \vspace{-0.2cm}
    \item Гуньо Адвокатина
    \newline
    \hspace*{0.5cm} Прозвището му назовава юридическа професия, което изисква чест и морал, но героят служи само на своите интереси. Изповядва устойчива житейска позиция – да се преструваме, че сме европейци, за да извлечем изгода, НО да си запазим нашенското.
    \vspace{-0.2cm}
    \item Данко Харсъзина
    \newline
    \hspace*{0.5cm} Това е образ, обединяващ всичко първично, примитивно и пошло. За него дори парите НЕ са важни, а възможността и удоволствието да упражнява насилие над другите (бий и псувай). За него чуждият е винаги враг, който трябва да бъде наказан с тояга. 
\end{itemize}
\vspace{0.1cm}
\newline
\textbf{5. Конфликтът родното и чуждото в творбата}
\vspace{0.1cm}
\newline 
\indent Основният конфликт в книгата е сблъсъкът между лицето и опакото на демокрацията. След освобождението институциите, законите, политиката, пресата гарантират равноправието между/на всички граждани. Байганьовщината обаче съзнателно изкривява тази роля и я подчинява на егоистичните си интереси. Европейското е само привидно. Позитивното се побългарява и се опорочава.
\newpage

\subsubsection{\colorbox{green!40}{Станислав Стратиев – „Балкански синдром“}}
\textbf{1. Социокултурен и литературен контекст}
\vspace{0.1cm}
\newline
\indent Станислав Стратиев, (Станко Стратиев Миладинов) е роден е София на 9.09.1941 г. Завършва гимназия в София (1959) и българска филология в Софийския университет (1968). Редактор е във вестник „Народна младеж“ (1964-1968) работи и за вестник „Стършел“ (1968–1975). От 1975 е драматург, директор (1983—1986) и художествен ръководител (1986—1989) в Държавен сатиричен театър, София. Оставя след себе си богато литературно наследство: драми, есета, комедии, миниатюри, повести, разкази, сатира, сценарии, фейлетони, фрагменти.
\vspace{0.1cm}
\newline
\indent Творчеството на Станислав Стратиев е създадено в условията на особена историческа ситуация, при която само литературата има смелостта да засегне болезнени, но неоспорими истини за изгубените идеали, за кошмара на бюрокрацията, за властващите сред хората меркантилизъм и конформизъм. Именно затова творецът поставя като централен проблем в своята пиеса „Балкански синдром" сблъсъкът между материалното и духовното, между нравствено извисените личности и безличната маса. В центъра на изображението са поставени незнайните хора от улицата, типизираните лица на човешката безпомощност, на оцеляващия инстинкт, на широко затворените очи за силата на задкулисието. Персонажите са назовани не с имена, а със социалните им роли или според биологичния им пол: Първи мъж, Зрител от Бели Искър, Кум, Бяла врана, Директор на театъра, Жената символ, Жената с проблем и др. Обезличаването е всъщност изискване на обществения строй за превъзходство на масовостта пред индивидуалността.
\vspace{0.1cm}
\newline
\textbf{2. Жанрово определение}
\vspace{0.1cm}
\newline
\indent Пиесата на Станислав Стратиев е комедия, чрез която са представени различни човешки слабости, пречупени през призмата на смешното. Авторът използва похватите на фарса и водевила, а представянето на злободневни проблеми сближават текста с този на фейлетона.
\vspace{0.1cm}
\newline
\textbf{3. Тълкуване на заглавието}
\vspace{0.1cm}
\newline
\indent Паратекстът е нестандартен, защото действието на пиесата излиза извън рамките на родното пространство. Заглавието е код за разбиране смисъла на произведението. Синдром означава проява, признак на нещо; външен белег на болест. Заглавието „казва“, че балканското е нещо нестандартно, то е отклонение от общоприетото, от нормалното т.е. европейското. Иначе казано паратекстът може да се разчете така:“ Българският характер като част от балканската нестандартност“.
\vspace{0.1cm}
\newline
\textbf{4. Композиция}
\vspace{0.1cm}
\newline
\indent Представени са различни сцени, като на финала е разрешено недоразумението. Времето на действието обхваща 80-те години на 20 век, времето на „преустройството“. Време на лицемерна загриженост на властта у нас към проблемите на обикновения човек. Пиесата е изградена от две действия, но не съдържа обичайната за драматургията структура.
\vspace{0.1cm}
\newline
\textbf{5. Основни проблеми в комедията}
\vspace{0.1cm}
\newline
\indent Началото въвежда във водещия проблем – фалша на обществото, преобърнатите ценности, за неистинското, абсурдно съществуване. Проблемът за липсата на връзка между науката и живота, за фалшивото изкуство. Проблемът на тоталитаризма – нагаждачеството на обикновения човек към силните на деня. Според Станислав Стратиев големият проблем на обществото от втората половина на 90-те г. на 20 век е у самите хора. Това е невъзможността да се откъснат от байганьовското, свързано с бездуховността, с примиренчеството, с нищонеправенето.
\vspace{0.1cm}
\newline
\textbf{6. Светът на родното и чуждото в пиесата}
\vspace{0.1cm}
\newline
\indent Изобразявайки абсурдни епизоди от нашата действителност, Стратиев изгражда иронично-пародиен образ на „родното“, което официално не признава изостаналостта си от „чуждото“, но все пак се стреми да го достигне. Чрез персонажа на Извънземното се представя друг, различен поглед към „родното“. За своите то, макар да има доста недостатъци, си остава незаменима ценност. И за да могат именно те - своите, да видят истинския му облик и да възприемат промяната като нещо неизбежно, е нужен неутралният поглед на Извънземното. Оценката, дадена чрез думите на Кума, за родното е, че през 80-те години на 20 век у нас всеки е добре, особено когато не плаща. И все пак в „родното „ има нещо хубаво – жените. Извънземното го разпознава в Цонка и се влюбва в нея. Така Стратиев се опитва да оневини своето чрез наличието на нещо, с което действително се славим – хубостта на българките.
\vspace{0.1cm}
\newline
\indent В света на „родното“ властва статиката, никой не е в състояние да излезе от застоя, защото пред всеки опит за промяна се изпречват препятствия, които обезсърчават човека и той се примирява със статуквото. Темата за единението е представена пародийно и иронично. „Чуждото“ е представено като контрастно на своето. То е символ на движението, на напредъка, на цивилизацията. Там човешкият дух е разкрепостен и прави свободен полет напред и нагоре.
\newpage


\subsection{Миналото и паметта}
\subsubsection{\colorbox{yellow!40}{Иван Вазов – „Паисий“}}
\textbf{1. Социокултурен и литературен контекст}
\vspace{0.1cm}
\newline
\indent Иван Вазов е феномен в българската литература, тъй като творчеството му показва най-богато жанрово многообразие, а в творбите му се оглежда българската същност. Авторът живее и твори както преди, така и след освобождението до началото на XX век. Това е време на преход, когато страната отново е самостоятелна държава, нуждаеща се от своя самобитен, културен, политически и икономически живот.
\vspace{0.1cm}
\newline
\indent Вазовото творчество оформя представата за българина, за самият себе си. Насочва погледът му към красивата природа и утвърждава родния език като свръхценност. Ако Ботев воюва за България с перо и меч, то Вазов служи на родината със словото си. Често творбите му се превръщат в щит срещу злостните нападки към родното. Приел мисията си на учител на нацията, Вазов утвърждава миналото като време на висши идеали, но и критикува следосвобожденското общество. Преобърнатата ценностна система, апатичността на младото население и късата национална памет провокират написването на цикъла „Епопея на забравените“. Дванадесетте оди изграждат националния пантеон, чрез който поетът превръща „забравените“ в „незабравимите“. Сред великите личности се откроява и монах Паисий сложил началото на Възраждането в България, осмислил службата си към Бог като служба към ближния. Неговата книжица се превръща в убежище на националната идентичност, защото миналото се противопоставя на забравата, затова одата е идейният център на цикъла.
\vspace{0.1cm}
\newline
\textbf{2. Тълкуване на заглавието и епиграфа}
\vspace{0.1cm}
\newline
\indent Епиграфът е цитат от Паисиевата творба, който съдържа идеологията на възраждането –  познаването на рода, езика и паметта като средство да се съхрани националното самосъзнание.
\vspace{0.1cm}
\newline
\indent Заглавието извежда главния герой като въплъщение на апостолската мисия на духовника. Задава очакване за автобиографичност на одата, но историческата недостоверност превръща текста в мит, в предание, в легенда за една изключителна личност.
\vspace{0.1cm}
\newline
\textbf{3. Жанрово определение}
\vspace{0.1cm}
\newline
\indent Вазовото произведение се определя в жанрово отношение като ода. Лирическият говорител изразява преклонението си пред митологизирания светогорски монах, представен като част от пантеона на легендарните национални герои. На преден план е прославата на обикновения духовен подвиг, който е разкрит като модел за подражание. Патетично, с ораторски стил, отличаващ се с експресивни езикови средства - тропи и фигури, - е осъществена апологията на персонажа и са утвърдени националните ценности, чието олицетворение е Паисий.
\vspace{0.1cm}
\newline
\textbf{4. Сюжет и композиция}
\vspace{0.1cm}
\newline
\indent Условния сюжет представя момента, в който се написва „История славянобългарска“.  Одата се фокусира върху съдържанието на Паисиевия труд. Монологът на героя е по-голямата част от текста, тъй като така Хилендарският монах дава обективна оценка за своето дело (на труда си). Представя възгледите си за връзката между историята и националното самоосъзнаване.
\vspace{0.1cm}
\newline
\indent Творбата на Вазов е съчетание между ода и поема, тъй като има събитие, което провокира емоции. Вазов утвърждава и възхвалява делото на монаха. Патетичният тон допринася за превръщане на образа в пример за подражание.
\vspace{0.1cm}
\newline
\indent Основната опозиция в одата е мрак – светлина. 
\vspace{0.1cm}
\newline
\indent Мракът е метафора на робското време, на една епоха, в която цари хаос и сякаш времето е спряло. Самият герой носи същата характеристика (тъмен), което се асоциира с изостаналост, невзрачност и дивота. Монахът пространствено е отдалечен от своите събратя, ограден е от пустошт, глухота и застой.
\vspace{0.1cm}
\newline
\indent Втората част на текста извежда образа на светлината – книгата се превръща в искрата, която осветява верния път на българите към бъдещето им. Книгата разказва мрака на примирението и възвръща българското себеуважение. От тъмен монах, Паисий се превръща в библейски пророк с ореол на святост. Подобно на глава първа от Библията, Паисий сътворява светлината на знанието, която побеждава хаоса и робството.
\vspace{0.1cm}
\newline
\textbf{5. Образът на Паисий като въплъщение на народен будител}
\vspace{0.1cm}
\newline
\indent Началото на одата чрез времевия топос 120 години връща времето назад преди написването на История славянобългарска – време на робство, несигурност, хаос, в който надеждата и вярата са изчезнали от народното битие. Тъмнината маркира цяла една епоха, в която българинът е подложен на не само на физическа, но и на духовна асимилация. Погледът на лирическия говорител се стеснява до пространството на килийката. Умалителното внася иронично пренебрежително отношение към божия служител и неговите занимания. Монологът на героя представлява трансформирането на мрака в светлина. Създаденото творение не е поредната църковна творба, а книга, нова, революционна за времето си, чрез която Паисий пропагандира вижданията си за миналото като източник на самочувствие и национална гордост. Това е така, защото Паисий е първият българин възприел паметта като живот. Мракът на робското търпение е изместен от светлината на пробуждането. Историята отваря глава първа от новата Библия на самосъзнание се българин. Повелителните глаголни форми „четете“ и „знайте“ въплъщават идеята на светогорския монах за словото като чудо. Знанието е светлината в началото на пътя на бъдещата нация (бъдещето на българската нация). Героят пренебрегва всички страдания и изпитания, осъзнавайки, че историята на народа си заслужава лишенията. Сътворява новото мислене, което ще изведе сънародниците му от хаоса на робството. Превръща обречения народ в нация. Самочувствието на автора оправдано и затова и Вазов го опоетизира и възвеличава, давайки му легендарно наименование – „Светогорец“.
\vspace{0.1cm}
\newline
\indent Одата „Паисий“ улавя момента, в който една книга гарантира победата на духовното умиране на един народ и осъществява чудото на възкресението – на роба в човек.
\vspace{0.1cm}
\newline
\indent С написването на историята се променя и самия образ на главния герой, който от прегърбен непознат монах се изправя в целия си ръст и се превръща във фигура, подобна на библейски пророк. Това е делото, съзнанието, че е променил целия български свят. В своя монолог Паисий излага виждането си за връзката между историята и националното себеутвърждаване. Посочените величави моменти са средство, чрез което авторът „прогонва мрака на примирението“ и показва, че културният и държавен падеж в миналото предопределят възхода на роба и в бъдещето. Пространство на тясната килийка се разширява до размерите на национален космос, в който е поставено началото на превръщането на народа в нация.
\vspace{0.1cm}
\newline
\indent Миналото, за което говори Паисий, е устояла на времето чрез паметта. За него родното ще оцелее. Идеята, че българското е вечно. Той е единственият герой в епопеята, който остава жив. Така той се превръща в метафора не само на Възраждането, но и на започналото духовно, българско пробуждане, което продължава и до днес. Преписите на  "История славянобългарска" говорят за факта, че тогава книгата е била нужна, но и днес, в наши дни, тя е в основата на българското самочувствие.
\newpage

\subsubsection{\colorbox{red!40}{Никола Вапцаров – „История“}}
\textbf{1. Социокултурен и литературен контекст}
\vspace{0.1cm}
\newline
\indent 
\subsubsection{\colorbox{red!40}{Йордан Радичков – „Ноев ковчег“}}

\subsection{Обществото и властта}
\subsubsection{\colorbox{red!40}{Христо Ботев – „Борба“}}
\subsubsection{\colorbox{red!40}{Елин Пелин – „Андрешко“}}
\subsubsection{\colorbox{red!40}{Христо Смирненски – „Приказка за стълбата“}}

\subsection{Животът и смъртта}
\subsubsection{\colorbox{red!40}{Христо Ботев – „До моето първо либе“}}
\subsubsection{\colorbox{red!40}{Иван Вазов – „Новото гробище над Сливница“}}
\subsubsection{\colorbox{red!40}{Емилиян Станев – „Крадецът на праскови“}}

\subsection{Природата}
\subsubsection{\colorbox{red!40}{Иван Вазов – „При Рилския манастир“}}
\subsubsection{\colorbox{red!40}{Пейо Яворов – „Градушка“}}
\subsubsection{\colorbox{red!40}{Пенчо Славейков – „Спи езерото“}}

\newpage

\section{12. клас}

\subsection{Любовта}

\subsubsection{\colorbox{green!40}{Димчо Дебелянов – „Аз искам да те помня все така“}}

\textbf{1. Социокултурен и литературен контекст}
\vspace{0.1cm}
\newline
\indent Символистичната лирика на Димчо Дебелянов представя характерните за направлението търсения на смисъла на живота, преследваното, но недостижимо щастие и разочарованието на модерния човек от действителността, лишена от хармония. Стиховете на автора носят и характерните опозиции, превърнати в код за развитието на творчеството му:
\begin{itemize}
    \vspace{-0.2cm}
    \item ред - хаос
    \vspace{-0.2cm}
    \item светлина - мрак
    \vspace{-0.2cm}
    \item живот - смърт
    \vspace{-0.2cm}
\end{itemize}

\begin{itemize}[leftmargin=*]
\item Под влиянието на Шопенхауер (германски философ, който следва и развива философията на Имануел Кант, отнасяща се за начина, по който преживяваме нещата от живота.) лиричните текстове на Дебелянов съдържат основните смислови ядра на символизма. Преобладаващият елегичен жанр в творчеството му не е случаен, тъй като битието е представено като страдание, а щастието е измамно. Осъзнал невъзможността за постигане на хармония, Дебеляновият лирически човек е страдащият и самотен индивид.
\vspace{-0.85cm}
\end{itemize}

\begin{itemize}[leftmargin=*]
\item Споменът се оказва единственият начин да се съхрани това, което е било и което в настоящето се оказва невъзможно.
\vspace{-0.3cm}
\end{itemize}

\begin{itemize}[leftmargin=*]
\item Дисхармоничната реалност налага усещане за безнадеждност, за невъзможност мечтата да се превърне в реалност. Любовта е неосъществима. Елегията „Аз искам да те помня все така“ е поредният Дебелянов текст, в който лирическият Аз изплаква своята болка, провокирана от факта, че красивото, истинско и дълбоко чувство е обречено. Съдбата се превръща във фактора, отнемащ правото на щастие на героите.
\vspace{-0.2cm}
\end{itemize}

\vspace{0.1cm}
\newline
\noindent 
\textbf{2. Тълкуване на заглавието}
\vspace{0.1cm}
\newline
\indent Както повечето Дебелянови текстове и този използва за паратекст първия стих. Личното местоимение „аз“ акцентира върху първоличната форма и налага звученето на изповед. Повелителното наклонение на глагола „искам“ заявява категоричното желание на Аза красотата на любовното чувство да бъде съхранена и пренесена от миналото през трагичното настояще към неясното бъдеще. 

\begin{itemize}[leftmargin=*]
\vspace{-0.2cm}
\item Глаголът „помня“ въвежда на паратекстово ниво основния за Дебеляновата лирика образ-символ – споменът, който съхранява онова, което не е било, но е могло да бъде.
\end{itemize}

\vspace{0.1cm}
\newline
\noindent 
\textbf{3. Жанрово определение}
\vspace{0.1cm}
\newline
\indent Творбата е изповедно лирическо стихотворение с елегично чувство. Преобладаващият тъжен тон се определя от ясната осъзнатост на героя, че любовта е обречена.

\vspace{0.1cm}
\newline
\noindent 
\textbf{4. Творческа предистория}
\vspace{0.1cm}
\newline 
\indent Произведението е написано по конкретен повод – любовта на поета към Мария Василева -  Звънчето. Първата редакция на текста носи посвещение и съдържа конкретни биографични факти. Втората редакция обаче е изчистена от всякаква конкретика, защото любовта е надличното чувство, универсално такова и така текстът звучи като обобщена представа за красивата, но обречена любов.

\vspace{0.1cm}
\newline
\noindent 
\textbf{5. "Сюжет" и композиция}
\vspace{0.1cm}
\newline 
\indent Условният за лирическия текст сюжет улавя момента на неизбежната раздяла, която определя преобладаващите емоции и размишления. Сцената акцентира върху  двамата влюбени, принудени да се разделят, защото смъртта тегне като прокоба над любимата жена.

\begin{itemize}[leftmargin=*]
\vspace{-0.2cm}
\item Композицията на текста е \colorbox{yellow!40}{кръгова}, като по този начин в края на текста се дава отговор на въпроса как лирическият Аз ще запомни любимата си – чрез спомена.
\end{itemize}

\vspace{0.1cm}
\newline
\noindent 
\textbf{6. Основни теми и мотиви}
\vspace{0.1cm}
\newline 
\indent Елегията възпроизвежда момента на неизбежната раздяла на двамата влюбени. Те са изведени извън неуютното, враждебно пространство на града и там, на хълма, те са заедно в своята самота и болка. Влюбените осъзнават, че любовта им е обречена и само споменът е в състояние да съхрани отминалото щастие. Споменът като символ НЕ е откритие на Дебелянов, но при него се превръща във феноменално средство, хвърлящо мост между миналото и бъдещето. Това се постига чрез характерните за поета „да“ конструкции. Те изразяват желанието на лирическия Аз мечтата да се превърне в бъдещ спомен. Затова е използвано и сегашно време, назоваващо и бъдещето („да те помня“).

\begin{itemize}[leftmargin=*]
\vspace{-0.2cm}
\item Трагизмът на текста идва от факта, че външен фактор е причината за неизбежната раздяла. Намесата на съдбата е необратима. Присъдата е произнесена и двамата влюбени са безсилни да променят ситуацията. Болестта взима надмощие, което проличава от описанието на любимата в текста. Епитетите в градация „бездомна“, „безнадеждна“ и „унила“ изграждат представата за отиващ си живота. Метафорите „пламнала ръка“ и „скръбен взор“ засилват усещането за безнадеждност, за обреченост. Смъртта се оказва непобедимия фактор, враг и на двамата, и става ясно, че е невъзможно да се избегне предначертаното. Парадоксалното твърдение „и любовта ни сякаш по е свята, защото трябва да се разделим.“ представя убеждението на лирическия Аз това, че любовта е надличностна, тя свръхценност и определението „свята“ я сакрализира. Именно раздялата е превръща в изключителна ценност, затова и споменът като „капсула на времето“ ще съхрани и пренесе святото чувство от миналото през трагичното настояще към неясното бъдеще.
\end{itemize}

\begin{itemize}[leftmargin=*]
\vspace{-0.2cm}
\item Втората строфа не уточнява причината за раздялата, а се фокусира върху „тръгването“ на любимата. Условното име, с което е назована тя („Морна, морна“ (унила, тъжна, уморена)) засилва представата за силата на смъртта – болестта вече е взела надмощие. Доказателство за това е библейският образ на скършения злак (подобно на зелена трева). По този начин е въведен мотива за безвреме прекършения живот, за отнетата младост. Личното местоимение \colorbox{yellow!40}{тя} е с главна буква, защото назовава съдбата (смъртта). Превръща е в герой. Лирическият герой се опитва да вдъхне увереност както на себе си, така и на своята любима. Глаголните форми в бъдеще време „ще останеш“, „ще се върнеш“ представляват убедеността на влюбения, че нещо толкова истинско и красиво не може да отиде да безвъзвратно.
\end{itemize}

\begin{itemize}[leftmargin=*]
\vspace{-0.2cm}
\item По този начин любовта се превръща във вечност, побеждава тленното. Въпреки дълбочината на чувството, обаче съдбата не може да бъде променена. Смъртта все по осезаемо, заявява присъствието си. Падащата нощ, подобно на погребален саван (плат, с който се увива трупа) обгръща двамата влюбени, а мрежите на прилепите създават усещане за пленници на съдбата. Ясната осъзнатост, че любимата е обречена, провокира оксиморонното твърдение „а в свойта вяра сам не вярвам аз.“. Лирическият Дебелянов човек не намира устои причини да вярва в положителния изход от ситуацията. Ако за Вапцаров вярата е живот и мотивация за промяна, то за символиста вярата е погубена от силата на съдбата.
\end{itemize}

\begin{itemize}[leftmargin=*]
\vspace{-0.2cm}
\item Кръговата композиция на текста представя разбирането на поета за омагьосания кръг на живота, в който човек е обречен на самота и болка. Последният стих дава всъщност отговор на въпроса как лирическият герой ще запази спомена за любимата. \colorbox{yellow!40}{Апосиопезата} в края изразява силните емоции на модерния човек, разбрал, че „тук“ и „сега“ щастието му е отказано и само споменът е в състояние да съхрани красотата на онова, което е било, но никога няма да се повтори. 
\end{itemize}
\newpage

\subsubsection{\colorbox{green!40}{Христо Фотев – „Колко си хубава!“}}

\textbf{1. Социокултурен и литературен контекст}
\vspace{0.1cm}
\newline
\indent Роденият в Цариград (Истанбул) творец завинаги свързва живота си и делото си с Бургас. Лириката на Христо Фотев представя един по-необичаен, нетрадиционен поглед към живота. Очарователният бохем, неуспелият художник и талантлив творец създава стихове, в които любовта е неизменна част от съществуването на човека, а животът е дар и затова трябва да се изживява истински и пълноценно. Бургаският бард има едновременно наивен и възторжен поглед към живота. Лирическият му герой с почти детско любопитство открива чудесата в прозата, в делничното. Той пренарежда ежедневното, с което и провокира/нарушава общоприетите представи за поведение.

\begin{itemize}[leftmargin=*]
\vspace{-0.2cm}
\item Поезията на Фотев появила се в средата на XX век, поставя две основни теми – любовта и свободата.  Възторженото чувство е представено като изпепеляване и страст, а свободата е отговорността на индивида към любимата, но и морална сила да си отидеш, когато не можеш и не искаш да приемеш правилата. Лириката на Христо Фотев превръща красотата в свръхценност. Тя предизвиква преклонение и възхищение и провокира любовта. Именно емблематичното стихотворение „Колко си хубава“ представя любовта като чувство, което променя мисленето и усещането на индивида. Творбата е от стихосбирката „Сантиментални посвещения“ е създадена по конкретен повод – силните чувства на Фотев към Милка Керемедчиева. Както Дебеляновият текст и този на Бургаския бохем е провокиран от конкретен субект, но биографичността е избегната, за да бъде представена любовта като универсална човешка емоция и красотата като етична категория, провокираща преклонение. Между флота вият текст и старобиблейските (старозаветните) „Песен на песните“ на Соломон съществува сериозна прилика – женската красота провокира възторг у любения и става причина за откритото възхищение и в двете творби.
\end{itemize}

\vspace{0.1cm}
\newline
\noindent 
\textbf{2. Жанрово определение}
\vspace{0.1cm}
\newline
\indent Лирическо изповедно стихотворение (НЕ любовно). Аз-формата дава възможност за постигане на максимална откровеност за разкриване на емоционалния свят на лирическия Аз. Текстът е диалогично насочен към любимата, която обаче е пасивен получател на информацията.

\vspace{0.1cm}
\newline
\noindent 
\textbf{3. Тълкуване на заглавието}
\vspace{0.1cm}
\newline
\indent

\subsubsection{\colorbox{green!40}{Петя Дубарова – „Посвещение“}}

\subsection{Вярата и надеждата}
\subsubsection{\colorbox{yellow!40}{Елин Пелин – „Спасова могила“}}
\subsubsection{\colorbox{red!40}{Атанас Далчев – „Молитва“}}
\subsubsection{\colorbox{red!40}{Никола Вапцаров – „Вяра“}}

\subsection{Трудът и творчеството}
\subsubsection{\colorbox{red!40}{Елин Пелин – „Ветрената мелница“}}
\subsubsection{\colorbox{red!40}{Йордан Йовков – „Песента на колелетата“}}
\subsubsection{\colorbox{red!40}{Виктор Пасков – „Балада за Георг Хених“}}

\subsection{Изборът и раздвоението}
\subsubsection{\colorbox{red!40}{Пейо Яворов – „Две души“}}
\subsubsection{\colorbox{red!40}{Елисавета Багряна – „Потомка“}}
\subsubsection{\colorbox{red!40}{Борис Христов – „Честен кръст“}}

\newpage

\section{Таблици върху раздели}
\subsection{11 клас}
\subsubsection{Родното и чуждото}

\newcommand{\bq}[1]{\quotedblbase #1\textquotedblleft}
\newcolumntype{L}[1]{>{\RaggedRight\hspace{0pt}\arraybackslash}p{#1}}
\setlength{\tabcolsep}{1.5pt}
\renewcommand{\arraystretch}{1.15}
\scriptsize

\begin{tabular}{|L{1.4cm}|L{1.4cm}|L{1.4cm}|L{2.1cm}|L{2.3cm}|L{2.6cm}|L{2.1cm}|L{3.0cm}|L{1.3cm}|}
\hline
\textbf{Произ\-веде\-ние} &
\textbf{Автор} &
\textbf{Жанр} &
\textbf{Особе\-ности на композицията} &
\textbf{Герои} &
\textbf{Теми, образи, мотиви} &
\textbf{Опозиции} &
\textbf{Идейни внушения} &
\textbf{Забе\-лежки} \\
\hline

\bq{Железният светилник} &
Димитър Талев &
Роман &
Четири части, всяка част има епиграф (мото) -- цитат от фолклорна песен; композиционна рамка -- сватбата &
Султана, Стоян, Лазар, Ния, Божана, Катерина, Рафе Клинче, Аврам Немтур &
Родното и чуждото, родното чуждо; родът и народът; светлината и желязото (железният светилник), словото &
Родно -- чуждо, традиционно (старо) -- ново, светлина -- мрак, материално -- духовно, морал -- любов &
Родът и домът се градят и отстояват с цената на жертви; словото и изкуството са дар; любовта е изпитание &
I част от тетралогия \\
\hline

\bq{Бай Ганьо журналист} &
Алеко Константинов &
Очерк (разказ, фейлетон) &
Разказ в разказа, композиционна рамка &
Гедрос (герой разказвач), Бай Ганьо, Дочоолу, Гочоолу, Гуньо Адвокатина, Данко Харсъзина &
Родното и чуждото не като етнически, а като етически категории, властта, журналистическата етика, демагогията &
Морал -- деградация на ценностната система (подлост, хамелеонщина, опортюнизъм, лична изгода) &
Интелигенцията като носител на морала е изблъскана на заден план, в обществото господстват Бай Ганьо и сподвижниците му; новоизградeният български социален свят е абсурден &
Последен фейлетон от книгата \bq{Бай Ганьо} \\
\hline

\bq{Балкански синдром} &
Станислав Стратиев &
Комедия на абсурда &
Две действия, фрагментарност &
Жена с проблем, Бяла Врана, Мъж от Бели Искър, Жена-символ, Директор, Кум, Бай Цончо, Йовчо, Печо, баба Кера, Извънземно &
Родното (балканското също е родно) и чуждото (извънземното -- цивилизация, която не ни разбира); разминаването между думи и дела, малкият човек, изгубен в алогичния свят &
Думи -- дела, профанно -- възвишено, безразличието към своето -- и устремът на модерната цивилизация &
Драмата на почтения човек, живеещ в непочтен свят. Свят, в който дела и думи много рядко се припокриват. Изостаналост, лична изгода, властта като сила. &
Заглавието отвежда към представата за болест \\
\hline
\end{tabular}
\normalsize

\subsubsection{Миналото и паметта}
\subsubsection{Обществото и властта}
\subsubsection{Животът и смъртта}
\subsubsection{Природата}
\subsection{12 клас}
\subsubsection{Любовта}

\newcommand{\bq}[1]{\quotedblbase #1\textquotedblleft}
\newcolumntype{L}[1]{>{\RaggedRight\hspace{0pt}\arraybackslash}p{#1}}
\setlength{\tabcolsep}{1.5pt}
\renewcommand{\arraystretch}{1.15}
\scriptsize

\begin{tabular}{|L{1.4cm}|L{1.4cm}|L{1.4cm}|L{2.1cm}|L{2.3cm}|L{2.6cm}|L{2.1cm}|L{3.0cm}|L{1.3cm}|}
\hline
\textbf{Произ\-веде\-ние} &
\textbf{Автор} &
\textbf{Жанр} &
\textbf{Особе\-ности на композицията} &
\textbf{Герои} &
\textbf{Теми, образи, мотиви} &
\textbf{Опозиции} &
\textbf{Идейни внушения} &
\textbf{Забе\-лежки} \\
\hline

\bq{Аз искам да те помня все така} &
Димчо Дебелянов &
Елегия &
Кръгова композиция, известна огледалност свързва първото и последното четиристишие &
Лирическият Аз, любимата (Морна) &
Обречената любов, макар и свръхценна не е възможно да бъде опазена. &
Надежда -- обреченост, смърт -- живот след смъртта &
Споменът е убежище за двама. Съзнателно извайван в настоящето, за да съхрани частица от него за бъдещето. Стаено е желание за вечност отвъд тленното. &
Първоначалното заглавие на творбата е \bq{Елегия} и е посветена \bq{На Зв.} \\
\hline
 
\bq{Колко си хубава!} &
Христо Фотев &
Стихотворение &
Рефрен, композиционна рамка &
Лирическият Аз и любимата &
Възхищението, желанието, възторга на влюбването, &
Загуба -- опазване (на любовта) &
Усещане за всеобхватност на любовта -- чувство, което прониква във всички човешки дейности и осмисля цялостното съществуване. &
Фотев избира за заглавие всекидневна фраза, която отказва да надгражда сложности и прилича на обикновен израз на възхита. \\
\hline
 
\bq{Посвещение} &
Петя Дубарова &
Стихотворение &
Стихотворението не е разделено на строфи. Ритъмът му е единен. Организирано е в двойки съседни рими. Така творбата се излива като ритмична, римувана, песенна цялост. &
Лирически Аз и любимият &
Темите за любовта и страха &
Страх -- увереност, мечта -- реалност, мрак -- светлина &
Свободата на личния избор е в противовес на обществените условности и общоприетото. Любовта е разностранна трансформация. Важно е не само да бъдат захвърлени оковите на страха и предразсъдъците, трябва мисията на облагородяването да бъде основна мотивация в любовта. &
Чрез заглавието лирическият субект изписва знаците на своята посветеност и ги насочва към любимия, като едновременно с това дарява и себе си. \\
\hline
\end{tabular}
\normalsize

\subsubsection{Вярата и надеждата}
\subsubsection{Трудът и творчеството}
\subsubsection{Изборът и раздвоението}

\newpage

\section{Модул трети теми - ДЗИ}
\subsection{Теми за есета}

\subsubsection{Властта на думите и думите на властта }
\vspace{0.5cm}
\begin{center}
\textbf{ЧАСТ 3 (Време за работа: 120 минути)}
\vspace{0.2cm}
{\textit{\underline{Аргументативния текст напишете в листа за отговори на предвиденото място.}}}
\vspace{0.1cm}
\end{center}
\textbf{41.} Напишете аргументативен текст в обем до 4 страници по ЕДНА от следните две теми:
\vspace{0.2cm}
\textbf{Тема за есе:}
\vspace{0.1cm}
\begin{center}
\textbf{„Властта на думите и думите на властта“}
\\
\textit{(Есе)}
\end{center}

\subsubsection{В търсене на очертанията на реалността и абсурда}
\vspace{0.5cm}
\begin{center}
\textbf{ЧАСТ 3 (Време за работа: 120 минути)}
\vspace{0.2cm}
{\textit{\underline{Аргументативния текст напишете в листа за отговори на предвиденото място.}}}
\vspace{0.1cm}
\end{center}
\textbf{41.} Напишете аргументативен текст в обем до 4 страници по ЕДНА от следните две теми:
\vspace{0.2cm}
\textbf{Тема за есе:}
\vspace{0.1cm}
\begin{center}
\textbf{„В търсене на очертанията на реалността и абсурда“}
\\
\textit{(Есе)}
\end{center}

\subsubsection{Всеки един строи кораба на своя живот.}
\vspace{0.5cm}
\begin{center}
\textbf{ЧАСТ 3 (Време за работа: 120 минути)}
\vspace{0.2cm}
{\textit{\underline{Аргументативния текст напишете в листа за отговори на предвиденото място.}}}
\vspace{0.1cm}
\end{center}
\textbf{41.} Напишете аргументативен текст в обем до 4 страници по ЕДНА от следните две теми:
\vspace{0.2cm}
\textbf{Тема за есе:}
\vspace{0.1cm}
\begin{center}
\textbf{„Всеки един строи кораба на своя живот.“}
\\
\textit{(Есе върху цитат от "Ноев ковчег")}
\end{center}

\subsubsection{Традициите - окови или опора}
\vspace{0.5cm}
\begin{center}
\textbf{ЧАСТ 3 (Време за работа: 120 минути)}
\vspace{0.2cm}
{\textit{\underline{Аргументативния текст напишете в листа за отговори на предвиденото място.}}}
\vspace{0.1cm}
\end{center}
\textbf{41.} Напишете аргументативен текст в обем до 4 страници по ЕДНА от следните две теми:
\vspace{0.2cm}
\textbf{Тема за есе:}
\vspace{0.1cm}
\begin{center}
\textbf{„Традициите - окови или опора“}
\\
\textit{(Есе)}
\end{center}

\subsubsection{Човекът – слуга или владетел на природата}
\vspace{0.5cm}
\begin{center}
\textbf{ЧАСТ 3 (Време за работа: 120 минути)}
\vspace{0.2cm}
{\textit{\underline{Аргументативния текст напишете в листа за отговори на предвиденото място.}}}
\vspace{0.1cm}
\end{center}
\textbf{41.} Напишете аргументативен текст в обем до 4 страници по ЕДНА от следните две теми:
\vspace{0.2cm}
\textbf{Тема за есе:}
\vspace{0.1cm}
\begin{center}
\textbf{„Човекът – слуга или владетел на природата“}
\\
\textit{(Есе)}
\end{center}

\subsubsection{Паметта срещу унищожителната сила на времето}
\vspace{0.5cm}
\begin{center}
\textbf{ЧАСТ 3 (Време за работа: 120 минути)}
\vspace{0.2cm}
{\textit{\underline{Аргументативния текст напишете в листа за отговори на предвиденото място.}}}
\vspace{0.1cm}
\end{center}
\textbf{41.} Напишете аргументативен текст в обем до 4 страници по ЕДНА от следните две теми:
\vspace{0.2cm}
\textbf{Тема за есе:}
\vspace{0.1cm}
\begin{center}
\textbf{„Паметта срещу унищожителната сила на времето“}
\\
\textit{(Есе)}
\end{center}

\subsubsection{Красотата в живота на човека}
\vspace{0.5cm}
\begin{center}
\textbf{ЧАСТ 3 (Време за работа: 120 минути)}
\vspace{0.2cm}
{\textit{\underline{Аргументативния текст напишете в листа за отговори на предвиденото място.}}}
\vspace{0.1cm}
\end{center}
\textbf{41.} Напишете аргументативен текст в обем до 4 страници по ЕДНА от следните две теми:
\vspace{0.2cm}
\textbf{Тема за есе:}
\vspace{0.1cm}
\begin{center}
\textbf{„Красотата в живота на човека“}
\\
\textit{(Есе)}
\end{center}

\subsubsection{Любовта е като лабиринт, в който лесно можеш да се изгубиш.}
\vspace{0.5cm}
\begin{center}
\textbf{ЧАСТ 3 (Време за работа: 120 минути)}
\vspace{0.2cm}
{\textit{\underline{Аргументативния текст напишете в листа за отговори на предвиденото място.}}}
\vspace{0.1cm}
\end{center}
\textbf{41.} Напишете аргументативен текст в обем до 4 страници по ЕДНА от следните две теми:
\vspace{0.2cm}
\textbf{Тема за есе:}
\vspace{0.1cm}
\begin{center}
\textbf{„Любовта е като лабиринт, в който лесно можеш да се изгубиш.“}
\\
\textit{(Есе върху цитат от Харуки Мураками)}
\end{center}

\subsubsection{Изкушенията на властта. Посветено на всички, които ще кажат: Това не се отнася до мене.}
\vspace{0.5cm}
\begin{center}
\textbf{ЧАСТ 3 (Време за работа: 120 минути)}
\vspace{0.2cm}
{\textit{\underline{Аргументативния текст напишете в листа за отговори на предвиденото място.}}}
\vspace{0.1cm}
\end{center}
\textbf{41.} Напишете аргументативен текст в обем до 4 страници по ЕДНА от следните две теми:
\vspace{0.2cm}
\textbf{Тема за есе:}
\vspace{0.1cm}
\begin{center}
\textbf{„Изкушенията на властта. "Посветено на всички, които ще кажат: "Това не се отнася до мене."}\textbf{"}\textbf{“}
\\
\textit{(Есе върху цитат от "Приказка за стълбата")}
\end{center}

\subsubsection{Самотата - благослов и прокоба}
\vspace{0.5cm}
\begin{center}
\textbf{ЧАСТ 3 (Време за работа: 120 минути)}
\vspace{0.2cm}
{\textit{\underline{Аргументативния текст напишете в листа за отговори на предвиденото място.}}}
\vspace{0.1cm}
\end{center}
\textbf{41.} Напишете аргументативен текст в обем до 4 страници по ЕДНА от следните две теми:
\vspace{0.2cm}
\textbf{Тема за есе:}
\vspace{0.1cm}
\begin{center}
\textbf{„Самотата - благослов и прокоба“}
\\
\textit{(Есе)}
\end{center}

\subsubsection{Нещастието създава човека, а човекът - щастието}
\vspace{0.5cm}
\begin{center}
\textbf{ЧАСТ 3 (Време за работа: 120 минути)}
\vspace{0.2cm}
{\textit{\underline{Аргументативния текст напишете в листа за отговори на предвиденото място.}}}
\vspace{0.1cm}
\end{center}
\textbf{41.} Напишете аргументативен текст в обем до 4 страници по ЕДНА от следните две теми:
\vspace{0.2cm}
\textbf{Тема за есе:}
\vspace{0.1cm}
\begin{center}
\textbf{„Нещастието създава човека, а човекът - щастието“}
\\
\textit{(Есе)}
\end{center}

\subsubsection{Добро, добрина и добродушие}
\vspace{0.5cm}
\begin{center}
\textbf{ЧАСТ 3 (Време за работа: 120 минути)}
\vspace{0.2cm}
{\textit{\underline{Аргументативния текст напишете в листа за отговори на предвиденото място.}}}
\vspace{0.1cm}
\end{center}
\textbf{41.} Напишете аргументативен текст в обем до 4 страници по ЕДНА от следните две теми:
\vspace{0.2cm}
\textbf{Тема за есе:}
\vspace{0.1cm}
\begin{center}
\textbf{„Добро, добрина и добродушие“}
\\
\textit{(Есе)}
\end{center}

\subsubsection{Каква полза за човек ако придобие цял свят, а погуби душата си...}
\vspace{0.5cm}
\begin{center}
\textbf{ЧАСТ 3 (Време за работа: 120 минути)}
\vspace{0.2cm}
{\textit{\underline{Аргументативния текст напишете в листа за отговори на предвиденото място.}}}
\vspace{0.1cm}
\end{center}
\textbf{41.} Напишете аргументативен текст в обем до 4 страници по ЕДНА от следните две теми:
\vspace{0.2cm}
\textbf{Тема за есе:}
\vspace{0.1cm}
\begin{center}
\textbf{„Каква полза за човек ако придобие цял свят, а погуби душата си...“}
\\
\textit{(Есе върху цитат от Евангелие от Марко)}
\end{center}

\subsubsection{Бай Ганьо – Безсмъртния Алеков герой}
\vspace{0.5cm}
\begin{center}
\textbf{ЧАСТ 3 (Време за работа: 120 минути)}
\vspace{0.2cm}
{\textit{\underline{Аргументативния текст напишете в листа за отговори на предвиденото място.}}}
\vspace{0.1cm}
\end{center}
\textbf{41.} Напишете аргументативен текст в обем до 4 страници по ЕДНА от следните две теми:
\vspace{0.2cm}
\textbf{Тема за есе:}
\vspace{0.1cm}
\begin{center}
\textbf{„Бай Ганьо – Безсмъртния Алеков герой“}
\\
\textit{(Есе)}
\end{center}

\subsubsection{Унизените и оскърбените в 21-ви век}
\vspace{0.5cm}
\begin{center}
\textbf{ЧАСТ 3 (Време за работа: 120 минути)}
\vspace{0.2cm}
{\textit{\underline{Аргументативния текст напишете в листа за отговори на предвиденото място.}}}
\vspace{0.1cm}
\end{center}
\textbf{41.} Напишете аргументативен текст в обем до 4 страници по ЕДНА от следните две теми:
\vspace{0.2cm}
\textbf{Тема за есе:}
\vspace{0.1cm}
\begin{center}
\textbf{„Унизените и оскърбените в 21-ви век“}
\\
\textit{(Есе)}
\end{center}

\subsubsection{Човекът между хаоса и светлината}
\vspace{0.5cm}
\begin{center}
\textbf{ЧАСТ 3 (Време за работа: 120 минути)}
\vspace{0.2cm}
{\textit{\underline{Аргументативния текст напишете в листа за отговори на предвиденото място.}}}
\vspace{0.1cm}
\end{center}
\textbf{41.} Напишете аргументативен текст в обем до 4 страници по ЕДНА от следните две теми:
\vspace{0.2cm}
\textbf{Тема за есе:}
\vspace{0.1cm}
\begin{center}
\textbf{„Човекът между хаоса и светлината“}
\\
\textit{(Есе)}
\end{center}

\subsubsection{Най-изящните цветя цъфтят измежду хищни плевели.}
\vspace{0.5cm}
\begin{center}
\textbf{ЧАСТ 3 (Време за работа: 120 минути)}
\vspace{0.2cm}
{\textit{\underline{Аргументативния текст напишете в листа за отговори на предвиденото място.}}}
\vspace{0.1cm}
\end{center}
\textbf{41.} Напишете аргументативен текст в обем до 4 страници по ЕДНА от следните две теми:
\vspace{0.2cm}
\textbf{Тема за есе:}
\vspace{0.1cm}
\begin{center}
\textbf{„Най-изящните цветя цъфтят измежду хищни плевели.“}
\\
\textit{(Есе)}
\end{center}

\subsubsection{Път}
\vspace{0.5cm}
\begin{center}
\textbf{ЧАСТ 3 (Време за работа: 120 минути)}
\vspace{0.2cm}
{\textit{\underline{Аргументативния текст напишете в листа за отговори на предвиденото място.}}}
\vspace{0.1cm}
\end{center}
\textbf{41.} Напишете аргументативен текст в обем до 4 страници по ЕДНА от следните две теми:
\vspace{0.2cm}
\textbf{Тема за есе:}
\vspace{0.1cm}
\begin{center}
\textbf{„Път.“}
\\
\textit{(Есе)}
\end{center}

\subsubsection{Самочувствие и чувство за малоценност}
\vspace{0.5cm}
\begin{center}
\textbf{ЧАСТ 3 (Време за работа: 120 минути)}
\vspace{0.2cm}
{\textit{\underline{Аргументативния текст напишете в листа за отговори на предвиденото място.}}}
\vspace{0.1cm}
\end{center}
\textbf{41.} Напишете аргументативен текст в обем до 4 страници по ЕДНА от следните две теми:
\vspace{0.2cm}
\textbf{Тема за есе:}
\vspace{0.1cm}
\begin{center}
\textbf{„Самочувствие и чувство за малоценност.“}
\\
\textit{(Есе)}
\end{center}

\subsubsection{Тежестта на думите}
\vspace{0.5cm}
\begin{center}
\textbf{ЧАСТ 3 (Време за работа: 120 минути)}
\vspace{0.2cm}
{\textit{\underline{Аргументативния текст напишете в листа за отговори на предвиденото място.}}}
\vspace{0.1cm}
\end{center}
\textbf{41.} Напишете аргументативен текст в обем до 4 страници по ЕДНА от следните две теми:
\vspace{0.2cm}
\textbf{Тема за есе:}
\vspace{0.1cm}
\begin{center}
\textbf{„Тежестта на думите.“}
\\
\textit{(Есе)}
\end{center}

\subsubsection{Свободата като илюзия}
\vspace{0.5cm}
\begin{center}
\textbf{ЧАСТ 3 (Време за работа: 120 минути)}
\vspace{0.2cm}
{\textit{\underline{Аргументативния текст напишете в листа за отговори на предвиденото място.}}}
\vspace{0.1cm}
\end{center}
\textbf{41.} Напишете аргументативен текст в обем до 4 страници по ЕДНА от следните две теми:
\vspace{0.2cm}
\textbf{Тема за есе:}
\vspace{0.1cm}
\begin{center}
\textbf{„Свободата като илюзия.“}
\\
\textit{(Есе)}
\end{center}

\subsubsection{Изкуството - абстрактна лудост или духовно величие}
\vspace{0.5cm}
\begin{center}
\textbf{ЧАСТ 3 (Време за работа: 120 минути)}
\vspace{0.2cm}
{\textit{\underline{Аргументативния текст напишете в листа за отговори на предвиденото място.}}}
\vspace{0.1cm}
\end{center}
\textbf{41.} Напишете аргументативен текст в обем до 4 страници по ЕДНА от следните две теми:
\vspace{0.2cm}
\textbf{Тема за есе:}
\vspace{0.1cm}
\begin{center}
\textbf{„Изкуството - абстрактна лудост или духовно величие.“}
\\
\textit{(Есе)}
\end{center}

\subsubsection{Свестните у нас считат за луди}
\vspace{0.5cm}
\begin{center}
\textbf{ЧАСТ 3 (Време за работа: 120 минути)}
\vspace{0.2cm}
{\textit{\underline{Аргументативния текст напишете в листа за отговори на предвиденото място.}}}
\vspace{0.1cm}
\end{center}
\textbf{41.} Напишете аргументативен текст в обем до 4 страници по ЕДНА от следните две теми:
\vspace{0.2cm}
\textbf{Тема за есе:}
\vspace{0.1cm}
\begin{center}
\textbf{„Свестните у нас считат за луди.“}
\\
\textit{(Есе върху цитат на Христо Ботев)}
\end{center}

\subsubsection{Изборът, който не принадлежи на едноликата ни душевност}
\vspace{0.5cm}
\begin{center}
\textbf{ЧАСТ 3 (Време за работа: 120 минути)}
\vspace{0.2cm}
{\textit{\underline{Аргументативния текст напишете в листа за отговори на предвиденото място.}}}
\vspace{0.1cm}
\end{center}
\textbf{41.} Напишете аргументативен текст в обем до 4 страници по ЕДНА от следните две теми:
\vspace{0.2cm}
\textbf{Тема за есе:}
\vspace{0.1cm}
\begin{center}
\textbf{„Изборът, който не принадлежи на едноликата ни душевност.“}
\\
\textit{(Есе)}
\end{center}

\subsubsection{Кривото огледало, в което се оглежда чуждия образ на традиционалността.}
\vspace{0.5cm}
\begin{center}
\textbf{ЧАСТ 3 (Време за работа: 120 минути)}
\vspace{0.2cm}
{\textit{\underline{Аргументативния текст напишете в листа за отговори на предвиденото място.}}}
\vspace{0.1cm}
\end{center}
\textbf{41.} Напишете аргументативен текст в обем до 4 страници по ЕДНА от следните две теми:
\vspace{0.2cm}
\textbf{Тема за есе:}
\vspace{0.1cm}
\begin{center}
\textbf{„Кривото огледало, в което се оглежда чуждия образ на традиционалността.“}
\\
\textit{(Есе)}
\end{center}

\subsubsection{Надеждата като психологическа необходимост.}
\vspace{0.5cm}
\begin{center}
\textbf{ЧАСТ 3 (Време за работа: 120 минути)}
\vspace{0.2cm}
{\textit{\underline{Аргументативния текст напишете в листа за отговори на предвиденото място.}}}
\vspace{0.1cm}
\end{center}
\textbf{41.} Напишете аргументативен текст в обем до 4 страници по ЕДНА от следните две теми:
\vspace{0.2cm}
\textbf{Тема за есе:}
\vspace{0.1cm}
\begin{center}
\textbf{„Надеждата като психологическа необходимост.“}
\\
\textit{(Есе)}
\end{center}

\subsubsection{Вярата като отказ от рациото или като негов щит.}
\vspace{0.5cm}
\begin{center}
\textbf{ЧАСТ 3 (Време за работа: 120 минути)}
\vspace{0.2cm}
{\textit{\underline{Аргументативния текст напишете в листа за отговори на предвиденото място.}}}
\vspace{0.1cm}
\end{center}
\textbf{41.} Напишете аргументативен текст в обем до 4 страници по ЕДНА от следните две теми:
\vspace{0.2cm}
\textbf{Тема за есе:}
\vspace{0.1cm}
\begin{center}
\textbf{„Вярата като отказ от рациото или като негов щит.“}
\\
\textit{(Есе)}
\end{center}

\subsubsection{Хармонията в природата - перфектната човешка мимикрия}
\vspace{0.5cm}
\begin{center}
\textbf{ЧАСТ 3 (Време за работа: 120 минути)}
\vspace{0.2cm}
{\textit{\underline{Аргументативния текст напишете в листа за отговори на предвиденото място.}}}
\vspace{0.1cm}
\end{center}
\textbf{41.} Напишете аргументативен текст в обем до 4 страници по ЕДНА от следните две теми:
\vspace{0.2cm}
\textbf{Тема за есе:}
\vspace{0.1cm}
\begin{center}
\textbf{„Хармонията в природата - перфектната човешка мимикрия.“}
\\
\textit{(Есе)}
\end{center}

\subsubsection{Ониоманиятa – проблем на нашето}
\vspace{0.5cm}
\begin{center}
\textbf{ЧАСТ 3 (Време за работа: 120 минути)}
\vspace{0.2cm}
{\textit{\underline{Аргументативния текст напишете в листа за отговори на предвиденото място.}}}
\vspace{0.1cm}
\end{center}
\textbf{41.} Напишете аргументативен текст в обем до 4 страници по ЕДНА от следните две теми:
\vspace{0.2cm}
\textbf{Тема за есе:}
\vspace{0.1cm}
\begin{center}
\textbf{„Ониоманиятa – проблем на нашето“}
\\
\textit{(Есе)}
\end{center}

\subsubsection{Автентичност и имитация в съвременния свят}
\vspace{0.5cm}
\begin{center}
\textbf{ЧАСТ 3 (Време за работа: 120 минути)}
\vspace{0.2cm}
{\textit{\underline{Аргументативния текст напишете в листа за отговори на предвиденото място.}}}
\vspace{0.1cm}
\end{center}
\textbf{41.} Напишете аргументативен текст в обем до 4 страници по ЕДНА от следните две теми:
\vspace{0.2cm}
\textbf{Тема за есе:}
\vspace{0.1cm}
\begin{center}
\textbf{„Автентичност и имитация в съвременния свят“}
\\
\textit{(Есе)}
\end{center}

\newpage

\subsection{Теми за интерпретативни съчинения}
\subsubsection{Свободата като резултат от съзнателно извършен грях (Елисавета Багряна)}
\vspace{0.5cm}
\begin{center}
\textbf{ЧАСТ 3 (Време за работа: 120 минути)}
\vspace{0.2cm}
{\textit{\underline{Аргументативния текст напишете в листа за отговори на предвиденото място.}}}
\vspace{0.1cm}
\end{center}
\textbf{41.} Напишете аргументативен текст в обем до 4 страници по ЕДНА от следните две теми:
\vspace{0.2cm}
\textbf{Тема за интерпретативно съчинение:}
\vspace{0.1cm}
\begin{center}
\textbf{„Свободата като резултат от съзнателно извършен грях“}
\\
\textit{(Интерпретативно съчинение върху „Потомка“ на Елисавета Багряна)}
\end{center}
\begin{flushleft}
Няма прародителски портрети,
\\
ни фамилна книга в моя род
\\
и не знам аз техните завети,
\\
техните лица, души, живот.
\vspace{0.5cm}
\\
Но усещам, в мене бие древна,
\\
скитническа, непокорна кръв.
\\
Тя от сън ме буди нощем гневно,
\\
тя ме води към греха ни пръв.
\vspace{0.5cm}
\\
Може би прабаба тъмноока,
\\
в свилени шалвари и тюрбан,
\\
е избягала в среднощ дълбока
\\
с някой чуждестранен, светъл хан.
\vspace{0.5cm}
\\
Конски тропот може би кънтял е
\\
из крайдунавските равнини
\\
и спасил е двама от кинжала
\\
вятърът, следите изравнил.
\vspace{0.5cm}
\\
Затова аз може би обичам
\\
необхватните с око поля,
\\
конски бяг под плясъка на бича,
\\
волен глас, по вятъра разлян.
\vspace{0.5cm}
\\
Може би съм грешна и коварна,
\\
може би средпът ще се сломя —
\\
аз съм само щерка твоя вярна,
\\
моя кръвна майчице-земя.
\end{flushleft}

\subsubsection{Място и роля на поезията в живота на обществото и човека (Борис Христов)}
\vspace{0.5cm}
\begin{center}
\textbf{ЧАСТ 3 (Време за работа: 120 минути)}
\vspace{0.2cm}
{\textit{\underline{Аргументативния текст напишете в листа за отговори на предвиденото място.}}}
\vspace{0.1cm}
\end{center}
\textbf{41.} Напишете аргументативен текст в обем до 4 страници по ЕДНА от следните две теми:
\vspace{0.2cm}
\textbf{Тема за интерпретативно съчинение:}
\vspace{0.1cm}
\begin{center}
\textbf{„Място и роля на поезията в живота на обществото и човека“}
\\
\textit{(Интерпретативно съчинение върху „Честен кръст“ на Борис Христов)}
\end{center}

\begin{table}[h]
\centering
\renewcommand{\arraystretch}{1.2}

\begin{tabular}{p{0.51\textwidth} p{0.47\textwidth}}

\textbf{I.} & \\

\begin{minipage}[t]{\linewidth}
Затварям ви за думите, уши - резето \\спускам.\\
Не искам в моя дом да се говори за \\изкуство -\\
\\
не искам в хленчене и мъдрословия угодни\\
напразно да се споменава името господне.\\

Когато имаше какво да каже на сърцето,\\
устата ми не криеше езика под небцето\\

и в стихове - спасителни - макар и тихо,\\
звънях отвътре, докато отвън ме биха...\\

Ала сега, докато в разума е още светло,\\
прибирам стълбата, с която стигах до небето,\\

и бързам да напусна на поезията храма,\\
преди да запълзят змии по младото ми рамо.\\

Каквото съм изпял дотук - ще го изтрия.\\
Ще вържа двете си ръце - устата ще зашия.\\

И ще се хвърля лудо на живота по вълните,\\
това, което някога измислих - да изпитам.\\

Не е за мойте зъби ореха да бъда гений,\\
щом не дойде смъртта за мен да се ожени,\\
\end{minipage}

&

\begin{minipage}[t]{\linewidth}

щом не можах до днес да изкова от думи кръста,\\
на който да се изкача и да възкръсна...\\

И затова сега, отскубнал се до корен,\\
ще се заема с труд по-благодарен и\\ достоен.\\

Ще мъкна камъни за мой дом и пот ще лея\\
над мисълта до утре как да преживея...\\

Или ще тръгна бос по връхчетата на тревата\\
подир виновницата за изгубения рай - жената.\\

И вързан до последната халка на нейната окова,\\
ще гледам как изсъхва моето дърво\\ духовно...\\

(…)
\end{minipage}
\\
\end{tabular}
\end{table}

\subsubsection{Робство, борба и свобода (Христо Ботев)}
\vspace{0.5cm}
\begin{center}
\textbf{ЧАСТ 3 (Време за работа: 120 минути)}
\vspace{0.2cm}
{\textit{\underline{Аргументативния текст напишете в листа за отговори на предвиденото място.}}}
\vspace{0.1cm}
\end{center}
\textbf{41.} Напишете аргументативен текст в обем до 4 страници по ЕДНА от следните две теми:
\vspace{0.2cm}
\textbf{Тема за интерпретативно съчинение:}
\vspace{0.1cm}
\begin{center}
\textbf{„Робство, борба и свобода“}
\\
\textit{(Интерпретативно съчинение върху „Борба“ на Христо Ботев)}
\end{center}

\begin{table}[h]
\centering
\renewcommand{\arraystretch}{1.2}
\begin{tabular}{p{0.48\textwidth} p{0.48\textwidth}}

\begin{minipage}[t]{\linewidth}

В тъги, в неволи младост минува,\\
кръвта се ядно в жили вълнува,\\
погледът мрачен, умът не види\\
добро ли, зло ли насреща иде…\\
На душа лежат спомени тежки,\\
злобна ги памет често повтаря,\\
в гърди ни любов, ни капка вяра,\\
нито надежда от сън мъртвешки\\
да можеш свестен човек събуди!\\
Свестните у нас считат за луди,\\
глупецът вредом всеки почита:\\
„Богат е“, казва, пък го не пита\\
колко е души изгорил живи,\\
сироти колко той е ограбил\\
и пред олтарят бога изамамил\\
с молитви, с клетви, с думи лъжливи.\\
И на обществен тоя мъчител\\
и поп, и черква с вяра слугуват;\\
нему се кланя дивак учител,\\
и с вестникарин зайдно мъдруват,\\
че страх от бога било начало\\
на сяка мъдрост… Туй е казало\\
стадо от вълци във овчи кожи,\\
камък основен за да положи\\
на лъжи свети, а ум човешки\\
да скове навек в окови тежки!\\
Соломон, тоя тиран развратен,\\
отдавна в раят найде запратен,\\
със свойте притчи между светците,\\


\end{minipage}

&

\begin{minipage}[t]{\linewidth}

казал е глупост между глупците,\\
и нея светът до днес повтаря —\\
„Бой се от бога, почитай царя!“\\

Свещена глупост! Векове цели\\
разум и совест с нея се борят;\\
борци са в мъки, в неволи мрели,\\
но кажи, що са могли да сторят!\\
Светът, привикнал хомот да влачи,\\
тиранство и зло и до днес тачи;\\
тежка желязна ръка целува,\\
лъжливи уста слуша с вяра:\\
млъчи, моли се, кога те биять\\
кожата ти да одере звярът\\
и кръвта да ти змии изпият,\\
на бога само ти се надявай:\\
„Боже, помилуй — грешен съм азе“\\
думай, моли се и твърдо вярвай —\\
бог не наказва, когото мрази…\\
Тъй върви светът! Лъжа и робство\\
на тая пуста земя царува!\\
И като залог из род в потомство\\
ден и нощ — вечно тук преминува.\\
И в това царство кърваво, грешно,\\
царство на подлост, разврат и сълзи,\\
царство на скърби — зло безконечно!\\
кипи борбата и с стъпки бързи\\
върви към своят свещени конец…\\
Ще викнем ние: „Хляб или свинец!“\\


\end{minipage}

\\

\end{tabular}
\end{table}

\subsubsection{Преобърнатите ценности (Алеко Константинов)}
\vspace{0.5cm}
\begin{center}
\textbf{ЧАСТ 3 (Време за работа: 120 минути)}
\vspace{0.2cm}
{\textit{\underline{Аргументативния текст напишете в листа за отговори на предвиденото място.}}}
\vspace{0.1cm}
\end{center}
\textbf{41.} Напишете аргументативен текст в обем до 4 страници по ЕДНА от следните две теми:
\vspace{0.2cm}
\textbf{Тема за интерпретативно съчинение:}
\vspace{0.1cm}
\begin{center}
\textbf{„Преобърнатите ценности“}
\\
\textit{(Интерпретативно съчинение върху „Бай Ганьо журналист“ на Алеко Константинов)}
\end{center}

    — А пък мен най ми харесва, като напопържам някого — изтърси чистосърдечно Данко Харсъзина, — отлеква ми н’ам как на душата!\\
    — Ашколсун бе, Харсъз, да живейш! — провиква се възхитен бай Ганьо и потупва Данка по рамото. — Работата е опечена. Ти, Гуньо, ще напишеш, както рекох, уводната статия, тъй ли? А вий, Гочоолу, и ти, Дочоолу, някои дописки, някои телеграми да изкалъпите.\\
    — Че какви телеграми, какви дописки? — запитват в недоумение двоицата.\\
    — Какви ли? Всякакви. Не виждате ли другите вестници? Пишете там: Ваше Царско Височество, народа ликува коленопреклонно и едногласно моли Всевишния... и прочее; наблъскайте там, каквото ви дойде на ума. Кажете там: Българският народ е дал хиляди доказателства,че когато се коснат до правата му, всички до един стават на крак... и... таквозинака... със сълзи на очите... и тъй нататък. Най-сетне какво ще ми дращете много-много; напсувайте опозицията, па вер селям. Ами! Какво ще седнем сега да се лигавим с разни философии. И кой ще те разбере! Нали е работата да замажем очи — карай колата. Не е ли тъй?\\
    — Е, хайде сбогом, че ме чакат клиенти — измъмра Гуньо Адвокатина и си взе шапката.\\
    — Чакат те тебе дяволите... Нейсе. Хайде, сбогом. Па чувай, Гуньо, утре рано уводната да бъде готова. Много здраве!\\
    — Всичко бошлаф. Да ви кажа ли аз вам? — заявява авторитетно бай Ганьо. — Нашия вестник ще кръстим или   „България за нас“, или „Народно величие“. Изберете си едно от двете.\\
    — „Народно величие“! Съгласен! „Народно величие“! Тъй, тъй, да живей!\\
    — Е сбогом, бай Ганьо.\\
    — Сбогом.\\
Гочоолу, Дочоолу и Гуньо Адвокатина излизат.\\
    — Ти, Данко, остани, ний с тебе ще пишем антрефилета.\\
    — Добре. Заръчай сега да донесат мастика и мезе и да почнем работата. Па да не домъкнат пак кисела бамя, гледай там някое мезе по-като хората — тук вестник има да се пише, не е шега.\\
Донесоха мастика и мезе. Ама ще попита бай Манолчо, какво мезе? То не е важно, важното е, че бай Ганьо и Данко Харсъзина се запретнаха да ръководят общественото мнение.\\
    — Данко бе, нашия комшия страшно ми се пери, учен бил, честен бил и н’ам какви дивотии. Да му теглим ли един калай?\\
    — Калай не, ами със земята барабар го направваме — заявява специалистът на псувните.\\
И почва се писането... „Научаваме се, че...“ — пише бай Ганьо и излага върху един лист бяла хартия такива черни хули срещу своя съсед, за каквито не само че не се е „научавал“ някой път, но нито насън му са минавали. Пише бай Ганьо, пише и зачерква, той е все недоволен от ядовитостта на своите стрели: крадец е за него нежна думица, той я зачерква и пише хайдук,   но тая дума е станала обикновена, бай Ганьо добавя пладнешки и я съединява с едно фатален. Съседът, жена му, децата му, родът му — излизат под перото на бай Ганя феноменални изверги... Той чете своето произведение на Данка Харсъзина. Данко, със светнали от мастиката очи, поощрява с въодушевление майстора на антрефилетата.\\
    — Карай, карай, карай! Блъскай, майка му стара, пред нищо недей се спира, пред нищо! Блъскаай! — гърми Данко Харсъзина, като че командува някое артилерийско сражение...\\

\subsubsection{Словото - знание и памет (Иван Вазов)}
\vspace{0.5cm}
\begin{center}
\textbf{ЧАСТ 3 (Време за работа: 120 минути)}
\vspace{0.2cm}
{\textit{\underline{Аргументативния текст напишете в листа за отговори на предвиденото място.}}}
\vspace{0.1cm}
\end{center}
\textbf{41.} Напишете аргументативен текст в обем до 4 страници по ЕДНА от следните две теми:
\vspace{0.2cm}
\textbf{Тема за интерпретативно съчинение:}
\vspace{0.1cm}
\begin{center}
\textbf{„Словото - знание и памет“}
\\
\textit{(Интерпретативно съчинение върху откъс от  „Паисий“ на Иван Вазов)}
\end{center}

\begin{tabular}{p{0.48\textwidth} p{0.48\textwidth}}
\begin{minipage}[t]{\linewidth}
Най-после отдъхна и рече: „Конец!\\
На житие ново аз турих венец.“\\
И той фърли поглед любовен, приветен\\
към тоз труд довършен, подвиг многолетен,\\
на волята рожба, на бденьето плод,\\
погълнал безшумно полвина живот —\\
житие велико! Заради което\\
той забрави всичко, дори и небето!\\
(...)\\
\\
„От днеска нататък българският род\\
история има и става народ!\\
Нека той познае от мойто писанье,\\
че голям е той бил и пак ще да стане,\\
че от славний Будин до светий Атон\\
бил е припознаван нашият закон.\\
Нека всякой брат наш да чете, да помни,\\
че гърците са люде хитри, вероломни,\\
че сме ги блъскали, и не един път —\\
и затуй не можат нази търпят —\\
и че сме имали царства и столици,\\
и от нашта рода светци и патрици;\\
че и ний сме дали нещо на светът\\
и на вси Словене книга да четът;\\
и кога му викат: «Българину!» бесно,\\
той да се гордее с това име честно.\\
Нека наш брат знае, че бог е велик\\\\
и че той разбира българский язик,\\
че е срам за всякой, който се отрича\\
от своя си рода и при гърци тича\\
и своето име и божия дар\\
зафърля безумно като един твар.\\
(...)\\
\end{minipage}
&
\begin{minipage}[t]{\linewidth}
Четете да знайте, що в стари години\\
по тез земи славни вършили деди ни,\\
как със много кралства имали са бран\\
и от царе силни вземали са дан,\\
и била велика българската държава;\\
как свети Борис се покръстил в Преслава,\\
как е цар Асен тук храмове градил\\
и дарове пращал; кой бе Самуил,\\
дето си изгуби душата във ада,\\
покори Дурацо и влезе в Елада;\\
четете и знайте кой бе цар Шишман\\
и как нашто царство сториха го плян;\\
кой би Иван Рилски, чийто свети мощи\\
чудеса се славят до тоя ден йоще;\\
как се Крум преславний с Никифора би\\
и из черепа му руйно вино пи\\
и как Симеон цар угрите прогони\\
и от Византия приема поклони.\\
А тоя беше учен, философ велик\\
и не се срамеше от своя език\\
и кога нямаше кого да надвива,\\
той пишеше книги, за да си почива.\\
Четете и знайте, що съм аз писал,\\
от много сказанья и книги събрал,\\
четете, о, братя, да ви се не смеят\\
и вам чужденците да не се гордеят…\\
На ви мойта книга, тя е вам завет,\\
нека де преписва и множи безчет\\
и пръска по всички поля и долини,\\
де българин страда, въздиша и гине.\\
Тя е откровенье, божа благодат —\\
младий прави мъдър, а стария — млад,\\
който я прочита няма да се кае,\\
който знае нея, много ще да знае.“\\
(...)\\
\\
Тъй мълвеше преди сто и двайсет годин\\
тоз див Светогорец — за рая негоден,\\
и фърляше тайно през мрака тогаз\\
най-първата искра в народната свяст.\\
\end{minipage}
\end{tabular}
\\

\subsubsection{Възраждащата сила на красотата (Емилиян Станев)}
\vspace{0.5cm}
\begin{center}
\textbf{ЧАСТ 3 (Време за работа: 120 минути)}
\vspace{0.2cm}
{\textit{\underline{Аргументативния текст напишете в листа за отговори на предвиденото място.}}}
\vspace{0.1cm}
\end{center}
\textbf{41.} Напишете аргументативен текст в обем до 4 страници по ЕДНА от следните две теми:
\vspace{0.2cm}
\textbf{Тема за интерпретативно съчинение:}
\vspace{0.1cm}
\begin{center}
\textbf{„Възраждащата сила на красотата“}
\\
\textit{(Интерпретативно съчинение върху „Крадецът на праскови“ на Емилиян Станев)}
\end{center}

\indentТя беше красива жена, преминала първата си младост, и сега излъчваше уморената и презряла хубост на отминаващо лято. В големите й очи се долавяше нещо замислено, твърдо, дори мрачно, което придаваше на погледа й студен израз, какъвто имаха очите на бездетните и незадоволени жени.\\
(...)
\indentСърцето й лудо блъскаше в гърдите, ту замираше, както замираше дишането й, а когато се озова при липата, беше принудена да седне под сянката на дървото, защото й се струваше, че няма да издържи. После видя пленника да се приближава и я обхвана безумно желание да избяга и да се затвори в старата колиба. Но беше късно. Тялото й остана покорно, безволно и натежало от желание…\\
\indentОт тоя ден тя престана да бъде тая Елисавета, която познаваше, като че под сянката на липата, в нейните подрасти, душата й бе разделена на същества: едното — примирената, угнетената жена, чакаща пристъпващата насреща й старост с безразлично отчаяние и тъга, и другото — непознато досега, вярващо, любещо и ликуващо същество, което отхвърляше нейния разум и желаеше да живее свободно и щастливо.\\
\indentВсеки следобед тя чакаше пленника до липата с притаен дъх и отмалели колене. Тогава усещаше тишината плътно да тежи край нея, кръвта й да пулсира в слепите очи и устните й да съхнат. И когато го видеше да се промъква край черниците и да идва насреща й, преставаше да диша…\\
\indentНямаше вече тия мъчителни пристъпи на съвест, които се появяваха в началото на техните срещи, нямаше ги страха, колебанието и нерешителността. Чувствата й ставаха все по-пълни и все по-силни и тя му се отдаваше с готовността на влюбена жена, познала възраждащата магия на любовта. Можеше да се каже, че се отдаваше за пръв път на мъж с цялото си тяло и с цялата си душа, алчна към всеки отлитащ миг от тия два часа, които двамата прекарваха под липата.\\
\indentПостепенно липата бе станала за нея живо същество и тя я гледаше отдалеч с весел поглед, както се гледа мълчалив и верен съучастник, или отиваше при нея в предвечерните часове, когато дългата и мощна сянка на дървото се просваше на изток като черна мантия, метната върху изгорялата от сушата трева. Тогава в главата й шумяха рояк смели мисли, тя мечтаеше, опиянена от спомените, изпълнена с благодарност и обич към света.

\subsubsection{Пътят към себе си (Йордан Йовков)}
\vspace{0.5cm}
\begin{center}
\textbf{ЧАСТ 3 (Време за работа: 120 минути)}
\vspace{0.2cm}
{\textit{\underline{Аргументативния текст напишете в листа за отговори на предвиденото място.}}}
\vspace{0.1cm}
\end{center}
\textbf{41.} Напишете аргументативен текст в обем до 4 страници по ЕДНА от следните две теми:
\vspace{-0.7cm}
\begin{center}
\textit{Паднало се 2026 година (Кати врачка)}
\end{center}
\vspace{-0.4cm}
\\
\vspace{0.2cm}
\textbf{Тема за интерпретативно съчинение:}
\vspace{0.1cm}
\begin{center}
\textbf{„Пътят към себе си“}
\\
\textit{(Интерпретативно съчинение върху откъс от  „Песента на колелетата“ на Йордан Йовков)}
\end{center}

\indentНикакви увещания, никакви молби не можеха да накарат Сали Яшар да остане до покъсно и да довърши някоя работа, колкото и бърза да беше тя. „Ще я направим – говореше
той, – и утре има време, и утре е божи ден.“ И той казваше това кротко, но твърдо – и сякаш
беше вече далеч някъде, откъснат от всичко, забравил за всичко, потънал в тая странна
замисленост, която пълнеше очите му, скрити под гъстите му вежди. И наметнал някоя дреха
на изпотените си плещи, той тръгваше за дома си, прегърбен малко, спокоен и загледан в
земята, а хората, които го срещаха, след като почтително го поздравяваха, гледаха го
учудено и си мислеха, че някаква болка трябва да гложди сърцето на Сали Яшар и че ако той
бърза да си отиде, не е, за да си почине, а за да остане сам със себе си и с мислите си. (…)\\
\indentИдеха часовете, когато Сали Яшар преставаше да гледа навън и обръщаше мислите си
към самия себе си. И прави бяха хората от Али Анифе: имаше нещо, което мъчеше Сали
Яшар, и всяка вечер по това време той си мислеше за него. (…)\\
\indentКаруците пееха по пътищата и сякаш разказваха как един човек може да бъде много
богат, но и много злочест. Защото за какво му беше това богатство на Сали Яшар? То дойде
много късно. На баира, отдясно на царския път, където каруците, които се връщат от града,
разсипват най-мелодичните си звуци, лежат в гробищата двамата му сина, два сокола.
Високи, грубо одялани камъни се издигат над тях, земята е обрасла с трева, като че никога не
е била разравяна, и люляката, която цъфти там всяка пролет, сякаш сам бог беше я посадил
отвеки. „Няма да им занеса нищо от моите жълтици – мислеше си Сали Яшар, – нито ще
взема нещо за себе си, когато и аз отида при тях.“ Тъмнината скриваше очите на Сали Яшар
и не можеше да се види какво има в тях, но той гледаше в земята като изтръпнал и мислеше,
че по-хубаво би било да беше беден, но да бяха живи двамата юнака, които и сега му се
усмихваха със сините си очи, румени, широкоплещести.\\
\indentДостигнал дотук, Сали Яшар започваше да мисли, че е длъжен да направи някакво
голямо благодеяние, истинско безкористно благодеяние, което да засегне повече хора, да
надживее много поколения. (…)\\
\indentЕто какво благодеяние искаше да направи и Сали Яшар. Но чешма не навсякъде
можеше да се направи, нямаше къде. Идеше му на ум да направи някъде кладенец, мост на
някое лошо място, хан, където замръкват пътници. Но всичко това му се виждаше или недостатъчно, или не твърде почтено. И той се биеше с мислите си, не знаеше какво да предприеме, безсилен като пред някоя мъчна задача, която не може да реши. (…)\\
\indentНа другия ден Сали Яшар беше в ковачницата си и – което рядко се случваше с него,
когато работеше – той три пъти подред слага чука си на наковалнята и, като скръстеше ръце
и се подпреше на него, оставаше тъй дълго време замислен, докато духалото работеше.
Първия път той си спомни онова, което Джапар беше му разказвал за връщането на Чауша.
Втория път той си спомни как сам беше чул каруцата, когато идеше Шакире, и, както винаги
в тия случаи, се просълзи. Третия път, когато Сали Яшар се подпря на чука и се замисли, той
сякаш като насън видя хиляди и хиляди каруци, които се връщат, които пеят по пътищата, а
от ниските къщурки, над които се носи вечерен дим и прах, излизат жени с деца на ръце,
други по-големи деца тичат наоколо им и всички вървят срещу каруците, за да ги
посрещнат...\\
\indentНикога Сали Яшар не беше изпитвал такова вълнение и такава радост. „Аллах! –
пошепна той и се улови за челото. – Аз съм бил сляп, аз съм бил глупав! Каква чешма и
какви мостове искам да правя? Себап! Има ли по-голям себап от тоя, който правя? Каруци
трябва да правя аз, каруци!“ (…)\\
\indentС чук в ръка той се спря пред тях и се замисли. Мъдрец наистина беше Сали Яшар,
много нещо беше видял, много нещо беше преживял, но едно беше ясно за него: с мъки, с
нещастия е пълен тоя свят, но все пак има нещо, което е хубаво, което стои над всичко друго
– любовта между хората.

\subsubsection{Защо езерото спи? (Пенчо Славейков)}
\vspace{0.5cm}
\begin{center}
\textbf{ЧАСТ 3 (Време за работа: 120 минути)}
\vspace{0.2cm}
{\textit{\underline{Аргументативния текст напишете в листа за отговори на предвиденото място.}}}
\vspace{0.1cm}
\end{center}
\textbf{41.} Напишете аргументативен текст в обем до 4 страници по ЕДНА от следните две теми:
\\
\vspace{0.2cm}
\textbf{Тема за интерпретативно съчинение:}
\vspace{0.1cm}
\begin{center}
\textbf{„Защо езерото спи?“}
\\
\textit{(Интерпретативно съчинение върху „Спи езерото“ на Пенчо Славейков)}
\end{center}

\newline
\newline
\noindentСпи езерото; белостволи буки\\
над него свождат вити гранки,\\
и в тихите му тъмни глъбини\\
преплитат отразени сянки.\\
\\
Треперят, шептят белостволи буки,\\
а то, замряло, нито трепва…\\
Понякога му сал повърхнини\\
дълга от лист отронен сепва.\\

\subsubsection{Човекът между надеждата и отчаянието (Пейо Яворов)}
\vspace{0.5cm}
\begin{center}
\textbf{ЧАСТ 3 (Време за работа: 120 минути)}
\vspace{0.2cm}
{\textit{\underline{Аргументативния текст напишете в листа за отговори на предвиденото място.}}}
\vspace{0.1cm}
\end{center}
\textbf{41.} Напишете аргументативен текст в обем до 4 страници по ЕДНА от следните две теми:
\\
\vspace{0.2cm}
\textbf{Тема за интерпретативно съчинение:}
\vspace{0.1cm}
\begin{center}
\textbf{„Човекът между надеждата и отчаянието.“}
\\
\textit{(Интерпретативно съчинение върху „Градушка“ на Пейо Яворов)}
\end{center}

\begin{tabular}{p{0.48\textwidth} p{0.48\textwidth}}
\begin{minipage}[t]{\linewidth}
Една, че две, че три усилни\\
и паметни години… Боже,\\
За някой грях ръце всесилни\\
издигна ти и нас наказа.\\
Кой ли може\\
неволя клетнишка изказа,\\
макар и — вчерашна се дума?\\
Да беше мор, да беше чума,\\
че в гроба гърло не гладува,\\
ни жадува!\\
А то градушка ни удари,\\
а то порой ни мътен влече,\\
слана попари, засух беше —\\
в земята зърно се опече…“\\
****\\
Но мина зима снеговита,\\
отиде пролет дъждовита —\\
и знойно лято позлати\\
до вчера злъчни широти.\\
Назряла вече тучна нива,\\
класец натегнал се привежда\\
и утешителка надежда\\
при труженик селяк отива.\\
Затопли радост на сърцето,\\
усмивка свърне на лицето,\\
въздишка кротка, пръст до пръст —\\
ръка неволно прави кръст:\\
„Да бъде тъй неделя още,\\
неделя пек и мирно време,\\
олекне ще и тежко бреме, —\\
на тежки мъки края дойде“.\\
****\\
\end{minipage}
&
\begin{minipage}[t]{\linewidth}
Додето сила има,\\
селяк без отдих труд се труди…\\
Почивка — ей я, би-ще зима.\\
Сега — недремнал и се буди:\\
петлите първи не пропели,\\
дори и куче не залае,\\
на крак е той… — „Катран и върви,\\
бре мъжо, взе ли от пазаря?“ Знае\\
невеста ранобудна — всичко,\\
готово е, но пак ще пита,\\
че утре жетва е; — самичко\\
сърце си знае как се стяга…\\
И сърдита\\
за нещо тя на двора бяга\\
и пак се връща в килеря,\\
тършува… „Днеска да намеря\\
Седмина още — чуеш жено?\\
— Ти ставай, Ваньо, — лентьо стига!\\
Небето ето — го червено,\\
че вече слънцето се дига,\\
добитъкът те гладен чака“.\\
А Ваньо сънен се прозява\\
и зад сайванта в полумрака\\
подал се — Сивчо го задява\\
с носът си влажен по вратлето.\\
„Хе ставай! — вика пак бащата\\
отсреди двора, на колата\\
затракал нещо, — стига…“ Ето\\
в съседен двор се вдига врява —\\
там някой люто се ругае.\\
Наблизко негде чук играе\\
и наковалнята отпява.\\
А сутренник повява леко\\
и звън от хлопки от далеко\\
донася в село; стадо блее…\\
На всякъде живот захваща.\\
И ето вече слънце грее\\
и на земята огън праща.\\
****
\end{minipage}
\end{tabular}

\newpage
\begin{tabular}{p{0.48\textwidth} p{0.48\textwidth}}
\begin{minipage}[t]{\linewidth}
Преваля пладне. Задух страшен.\\
И всеки дигне взор уплашен,\\
с ръкав избрише си челото\\
и дълго гледа към небето,\\
а то сивее мъгловито.\\
И слънцето жълтей сърдито:\\
От юг бухлат се облак дига\\
пълзи и вече го настига;\\
по-доле вирнали главите\\
и други… Знак е — чуй петлите.\\
А гъски около реката\\
защо размахали крилата,\\
и те са глупи закрещели,\\
Какво ли са орали, сели?\\
****\\
Върни се облако неверен, —\\
почакай, пакостнико черен,\\
неделя — две… ела тогази,\\
страшилище! А облак лази,\\
ръсте и вий снага космата,\\
засланя слънце; в небесата\\
тъмней зловещо… Милост няма!\\
Ще стане пак беда голяма. —\\
на завет всичко се прибира,\\
сърцето трепне, в страх премира,\\
че горе — дим и адски тътен.\\
Върхушка, прах… ей свода мътен\\
продран запалва се — и блясък —\\
и още — пак, — О Боже!… Трясък\\
оглася планини, полета —\\
земя трепери… Град! — парчета —\\
яйце и орех… Спри… Недей…\\
Труд кървав, Боже, пожалей!\\
****\\
\end{minipage}
&
\begin{minipage}[t]{\linewidth}
Но свърши. Тихо гръм последен\\
заглъхва нейде надалече\\
и вълк на стадо — вихър леден\\
подгоня облаците вече.\\
А ето слънцето огряло\\
тъжовно гледа върволица\\
от стари, млади и дечица\\
забързали навън от село;\\
в калта подпретнали се боси,\\
глави неволнишки навели,\\
отиват, — черно зло ги носи\\
в нивята грозно опустели.\\
Че там жетварка бясна хала\\
просо, пшеница, ръж, ечмени —\\
безрадно, зрели и зелени,\\
и цвет — надежди е познала…\\
\end{minipage}
\end{tabular}

\subsubsection{Раждането на българския свят (Димитър Талев)}
\vspace{0.5cm}
\begin{center}
\textbf{ЧАСТ 3 (Време за работа: 120 минути)}
\vspace{0.2cm}
{\textit{\underline{Аргументативния текст напишете в листа за отговори на предвиденото място.}}}
\vspace{0.1cm}
\end{center}
\textbf{41.} Напишете аргументативен текст в обем до 4 страници по ЕДНА от следните две теми:
\\
\vspace{0.2cm}
\textbf{Тема за интерпретативно съчинение:}
\vspace{0.1cm}
\begin{center}
\textbf{"Раждането на българския свят"}
\\
\textit{(Интерпретативно съчинение върху „Железният светилник“ на Димитър Талев)}
\end{center}

\vspace{0.5cm}
(Част втора: „Тъмни времена“, глава 12)
\\
\indentПреспанци градяха своя нов храм с необикновено увлечение, с някакво изстъпление. Пролетта беше ранна и топла. Още в първите дни на април работата на църковното място закипя. Идваха доброволни работници от целия град, мъже и жени, млади и стари. Идваха и селяни от близките села, докарваха с колите си ломен камък, вар, дървесен материал, помагаха с десниците си. Цялата Преспа живееше тогава с мисълта за новата църква.
\\
\indentНа църковния двор, ограден с високи дъски, всеки ден се трупаха десетки, стотици хора. Чорбаджии и бедняци работеха рамо до рамо; богати търговци, чиито ръце никога не бяха хващали копач или лопата, сега пренасяха камъни и бъркаха вар, забравили гордост и суета. Жените пренасяха вода в стомни и бакъри, децата метяха и чистеха стърготините около майсторите дюлгери. Никой не искаше пари, никой не жалеше сили.
\\
\indentВсичко се правеше с песен, с шеги, с радостни викове, сякаш хората не бяха на тежка работа, а на някакъв голям, всенароден празник. Старите преспанци, които помнеха миналите времена на страх, на усамотение и унижение, когато всеки се криеше в своята къща като в дупка, гледаха това чудо и плачеха от радост. Раждаше се нещо ново, непознато дотогава в тоя град — преспанци усещаха, че стават едно тяло, един народ, който има една мисъл, една воля и един бог.

\vspace{0.5cm}

(Част трета: „Народ се пробужда“, глава 17)
\\
\indentНаместникът пристъпи пред олтара, облечен в тежките си, златоткани одежди, и почна да чете евангелието на гръцки език. Гласът му беше силен, отривист, но в него се усещаше някаква скрита, трескава тревога. Той бързаше, сякаш искаше по-скоро да свърши, да надвие необикновеното мълчание, което тегнеше в препълнения храм. Още при първите гръцки думи в черквата стана такова движение, такова глухо шумене, сякаш силен вятър изведнъж разлюля гъста, вековна гора. Хората се споглеждаха, притискаха се един в друг, устните им пресъхваха от вълнение.
\\
\indentТогава Лазар Глаушев пристъпи напред от предните редици. Той беше бледен, но очите му горяха с някакъв непознат, устремен пламък. Той излезе пред самия амвон, вдигна високо дясната си ръка и гласът му прозвуча над главите на уплашените богомолци, чист, ясен и властен, та пресече изведнъж думите на гръка:
— Братя българи! Докога ще слушаме чужди думи в Божия дом? Наш е този храм, с наши кървави пари е граден, наш е и Бог, и езикът ни трябва да бъде наш, че да разбираме какво се молим и за какво плачем!
\\
\indentВ храма настана мъртвешка тишина, в която се чуваше само тежкото дишане на стотиците гърди. Гръцкият наместник замръзна на мястото си, изпуснал книгата, с лице, позеленяло от ярост и изненада. И преди той или хората му да се опомнят, Лазар разгърна славянското евангелие и започна да чете същото слово на чист, звънък, разбран за всеки старец и дете български език. В същия миг черковният хор, воден от Климент Бенков, поде песнопението със същата невиждана сила. Гласът на певците отекна в купола, но не остана сам — изведнъж цялата черква, отпред до самия праг, поде песента. Мъже, жени, старци пееха със сълзи на очи. Това вече не беше просто черковна служба, това беше гласът на един събуден народ, който разкъсваше веригите си и заявяваше пред бога и света своето право на живот.

\subsubsection{Казусът "балкански синдром" (Станислав Стратиев)}
\vspace{0.5cm}
\begin{center}
\textbf{ЧАСТ 3 (Време за работа: 120 минути)}
\vspace{0.2cm}
{\textit{\underline{Аргументативния текст напишете в листа за отговори на предвиденото място.}}}
\vspace{0.1cm}
\end{center}
\textbf{41.} Напишете аргументативен текст в обем до 4 страници по ЕДНА от следните две теми:
\\
\vspace{0.2cm}
\textbf{Тема за интерпретативно съчинение:}
\vspace{0.1cm}
\begin{center}
\textbf{"Казусът "балкански синдром""}
\\
\textit{(Интерпретативно съчинение върху „Балкански синдром“ на Станислав Стратиев)}
\end{center}
\indentЖЕНА-СИМВОЛ. Това ли ще гледаме сега?… Тези ли примитиви?… Защо не започва пиесата?… Изхвърлете ги най-после!… Ние сме дошли тук да гледаме театър!…
\\
\indentДИРЕКТОР. Виждате, ситуацията е сложна…
\\
\indentЖЕНА-СИМВОЛ. Не виждам какво му е толкова сложното. Изхвърляте ги тези… и продължавате спектакъла. Стига вече простотии, стига вече балканщини. Къде е тук животът на човешкия дух? Сега сме края на двайсети век!
\\
\indentКУМ (чак сега се ориентира). Ние ли сме „простотии“? Ние ли сме балканщини? Значи народът, дето ви храни и облича, е „простотии“?… Като ви изкараме на пазара чушки и домати, не сме простотии, ама като дойдеме в театъра, сме простотии, така ли?
\\
\indentДИРЕКТОР. Каратанасов, не влизай в пререкания! (Към Жената.) Моля ви, влезте ни в положението!…
\\
\indentЖЕНА-СИМВОЛ. Няма да вляза! Трийсет години вече само това правя – влизам някому в положението!… Защо някой не влезе в моето положение!…
\\
\indentКУМ. Как да влезе, я се погледни каква си!… (Смее се.)
\\
\indentДИРЕКТОР. Каратанасов!…
\\
\indentЖЕНА-СИМВОЛ. Кога ще престанете най-после с тази простащина! Театърът е държавен!…
\\
\indentКУМ. А държавата на кого е? На народа!
\\
\indentЖЕНА-СИМВОЛ. Ще прекратите ли най-после този кич или не?
\\
\indentКУМ. Ти си кич!…
\\
\indentДИРЕКТОР. Каратанасов!…
\\
\indentЖЕНА-СИМВОЛ. И това ми било Сатиричен театър. Правилно ви отнеха званията. Народни артисти, заслужили режисьори, деятели на културата!!!
\\
\indentДИРЕКТОР. Каратанасов, защо влизаш в спор с публиката? Хората са прави, те са дошли на театър. Не можем така да ги разгонваме.
\\
\indentКУМ. Ма тя обижда народа бе, началство, това преглъща ли се? Тебе да те нарекат кич, ти как ще се чувстваш?…
\\
\indentДИРЕКТОР. Ама ти защо се разправяш от сцената? Имаме достатъчно артисти, не сме опрели до тебе!
\\
\indentЖЕНА-СИМВОЛ. Халтураджии! (Излиза.)
\\
\indentКУМ. Началник, не се ядосвай! Тая не е добре. Е, аз се ядосах и огладнях. Ти не си ли гладен?… Минчо, Минчо!… (Минчо се явява.) Я донеси за хапване, че ще повредим гастрита така, гладни!…
\\
\indentМИНЧО. Окей!… (Излиза.)

\subsubsection{...най-първата искра в народната свяст (Иван Вазов)}
\vspace{0.5cm}
\begin{center}
\textbf{ЧАСТ 3 (Време за работа: 120 минути)}
\vspace{0.2cm}
{\textit{\underline{Аргументативния текст напишете в листа за отговори на предвиденото място.}}}
\vspace{0.1cm}
\end{center}
\textbf{41.} Напишете аргументативен текст в обем до 4 страници по ЕДНА от следните две теми:
\\
\vspace{0.2cm}
\textbf{Тема за интерпретативно съчинение:}
\vspace{0.1cm}
\begin{center}
\textbf{"...най-първата искра в народната свяст"}
\\
\textit{(Интерпретативно съчинение върху цитат от „Паисий“ на Иван Вазов)}
\end{center}

\begin{tabular}{p{0.48\textwidth} p{0.48\textwidth}}
\begin{minipage}[t]{\linewidth}
Най-после отдъхна и рече: „Конец!\\
На житие ново аз турих венец.“\\
И той фърли поглед любовен, приветен\\
към тоз труд довършен, подвиг многолетен,\\
на волята рожба, на бденьето плод,\\
погълнал безшумно полвина живот —\\
житие велико! Заради което\\
той забрави всичко, дори и небето!\\
(...)\\
\\
„От днеска нататък българският род\\
история има и става народ!\\
Нека той познае от мойто писанье,\\
че голям е той бил и пак ще да стане,\\
че от славний Будин до светий Атон\\
бил е припознаван нашият закон.\\
Нека всякой брат наш да чете, да помни,\\
че гърците са люде хитри, вероломни,\\
че сме ги блъскали, и не един път —\\
и затуй не можат нази търпят —\\
и че сме имали царства и столици,\\
и от нашта рода светци и патрици;\\
че и ний сме дали нещо на светът\\
и на вси Словене книга да четът;\\
и кога му викат: «Българину!» бесно,\\
той да се гордее с това име честно.\\
Нека наш брат знае, че бог е велик\\\\
и че той разбира българский язик,\\
че е срам за всякой, който се отрича\\
от своя си рода и при гърци тича\\
и своето име и божия дар\\
зафърля безумно като един твар.\\
(...)\\
\end{minipage}
&
\begin{minipage}[t]{\linewidth}
Четете да знайте, що в стари години\\
по тез земи славни вършили деди ни,\\
как със много кралства имали са бран\\
и от царе силни вземали са дан,\\
и била велика българската държава;\\
как свети Борис се покръстил в Преслава,\\
как е цар Асен тук храмове градил\\
и дарове пращал; кой бе Самуил,\\
дето си изгуби душата във ада,\\
покори Дурацо и влезе в Елада;\\
четете и знайте кой бе цар Шишман\\
и как нашто царство сториха го плян;\\
кой би Иван Рилски, чийто свети мощи\\
чудеса се славят до тоя ден йоще;\\
как се Крум преславний с Никифора би\\
и из черепа му руйно вино пи\\
и как Симеон цар угрите прогони\\
и от Византия приема поклони.\\
А тоя беше учен, философ велик\\
и не се срамеше от своя език\\
и кога нямаше кого да надвива,\\
той пишеше книги, за да си почива.\\
Четете и знайте, що съм аз писал,\\
от много сказанья и книги събрал,\\
четете, о, братя, да ви се не смеят\\
и вам чужденците да не се гордеят…\\
На ви мойта книга, тя е вам завет,\\
нека де преписва и множи безчет\\
и пръска по всички поля и долини,\\
де българин страда, въздиша и гине.\\
Тя е откровенье, божа благодат —\\
младий прави мъдър, а стария — млад,\\
който я прочита няма да се кае,\\
който знае нея, много ще да знае.“\\
(...)\\
\\
Тъй мълвеше преди сто и двайсет годин\\
тоз див Светогорец — за рая негоден,\\
и фърляше тайно през мрака тогаз\\
най-първата искра в народната свяст.\\
\end{minipage}
\end{tabular}
\\

\subsubsection{Историята и поезията като памет (Никола Вапцаров)}
\vspace{0.5cm}
\begin{center}
\textbf{ЧАСТ 3 (Време за работа: 120 минути)}
\vspace{0.2cm}
{\textit{\underline{Аргументативния текст напишете в листа за отговори на предвиденото място.}}}
\vspace{0.1cm}
\end{center}
\textbf{41.} Напишете аргументативен текст в обем до 4 страници по ЕДНА от следните две теми:
\\
\vspace{0.2cm}
\textbf{Тема за интерпретативно съчинение:}
\vspace{0.1cm}
\begin{center}
\textbf{"Историята и поезията като памет"}
\\
\textit{(Интерпретативно съчинение върху „История“ на Никола Вапцаров)}
\end{center}

\begin{tabular}{p{0.48\textwidth} p{0.48\textwidth}}
\begin{minipage}[t]{\linewidth}
Какво ще ни дадеш, историйо,\\
от пожълтелите си страници? —\\
Ний бяхме неизвестни хора\\
от фабрики и канцеларии,\\
\\
ний бяхме селяни, които\\
миришеха на лук и вкиснало,\\
и под мустаците увиснали\\
живота псувахме сърдито.\\
\\
Ще бъдеш ли поне признателна,\\
че те нахранихме с събития\\
и те напоихме богато\\
с кръвта на хиляди убити.\\
\\
Ще хванеш контурите само,\\
а вътре, знам, ще бъде празно\\
и няма никой да разказва\\
за простата човешка драма.\\
\\
Поетите ще са улисани\\
във темпове и във агитки\\
и нашта мъка ненаписана\\
сама в пространството ще скита.\\
\\
Живот ли бе — да го опишеш?\\
Живот ли бе — да го разровиш?\\
Разровиш ли го — ще мирише\\
и ще горчи като отрова.\\
\\
По синорите сме се раждали,\\
на завет някъде до тръните,\\
а майките лежали влажни\\
и гризли сухите си бърни.\\
\\
Като мухи сме мрели есен,\\
жените вили по задушница,\\
изкарвали плача на песен,\\
но само бурена ги слушал.\\
\\
Онез, които сме оставали,\\
се потехме и под езика,\\
работехме къде що хванем,\\
работехме като добитък.\\
\end{minipage}
&
\begin{minipage}[t]{\linewidth}
Мъдруваха бащите в къщи:\\
„Така било е и ще бъде…“\\
А ние плюехме намръщено\\
на оглупялата им мъдрост.\\
\\
Зарязвахме софрите троснато\\
и търтвахме навън, където\\
една надежда ни докосваше\\
със нещо хубаво и светло.\\
\\
О, как сме чакали напрегнато\\
в задръстените кафенета!\\
И късно през нощта си легахме\\
с последните комюникета.\\
\\
О, как се люшкахме в надеждите!…\\
А тегнеше небето ниско,\\
свистеше въздуха нажежен…\\
Не мога повече! Не искам!…\\
\\
Но в многотомните писания,\\
под буквите и редовете\\
ще вика нашето страдание\\
и ще се зъби неприветно —\\
\\
защото би ни безпощадно\\
живота с тежките си лапи\\
направо по устата гладни,\\
затуй езика ни е грапав.\\
\\
И стиховете, дето пишем,\\
когато краднем от съня си,\\
парфюмен аромат не дишат,\\
а са навъсени и къси.\\
\\
За мъката — не щем награди,\\
не ще дотегнем и с клишета\\
на томовете ти грамади,\\
натрупани през вековете.\\
\\
Но разкажи със думи прости\\
на тях — на бъдещите хора,\\
които ще поемат поста ни,\\
че ние храбро сме се борили.\\
\end{minipage}
\end{tabular}

\subsubsection{Свирепата философия за света на Йордан Радичков (Йордан Радичков)}
\vspace{0.5cm}
\begin{center}
\textbf{ЧАСТ 3 (Време за работа: 120 минути)}
\vspace{0.2cm}
{\textit{\underline{Аргументативния текст напишете в листа за отговори на предвиденото място.}}}
\vspace{0.1cm}
\end{center}
\textbf{41.} Напишете аргументативен текст в обем до 4 страници по ЕДНА от следните две теми:
\\
\vspace{0.2cm}
\textbf{Тема за интерпретативно съчинение:}
\vspace{0.1cm}
\begin{center}
\textbf{"Свирепата философия за света на Йордан Радичков"}
\\
\textit{(Интерпретативно съчинение върху \textcolor{red}{„Ноев ковчег“} на Никола Вапцаров)}
\end{center}
\indentКолко дни и седмици валя, никой не може да каже, защото дните и седмиците се сляха в едно цяло. После, за да внесе малко бъркотия в еднообразието, иззад хълмовете се втурна виелица. Тя насмете земята, тъй както вълци насмитат стадо овце. Грозен вой нададоха осаждените комини на къщите, селищата притихнаха, стенания, скърцания и пукот се изтръгнаха от горите и като накара всяко сърце да замре, виелицата отмина, кикотейки се. Ден или два подир нея валя много разсеяно, снежинките блуждаеха колебливо из въздуха, въртяха се около себе си, полюляваха се и падаха замислено върху земята, изтърсили се върху нея едва ли не по невнимание. Сякаш в машинарията на небесната мелница някакви механизми се бяха разстроили. Но постепенно, макар и бавно, механизмите се оправиха и снегът отново продължи да вали както и по-рано и да затрупва всичко под себе си. Валеше всеки ден и всяка нощ, ханчето постепенно се изгуби, засипано със сняг, само гуглата му все още продължаваше да стърчи и озъбените срещу небето колове на онзи ужасен плет. Изгуби се също тъй и снежният отпечатък на българския апостол, както и широката пъртина, оставена от заптиетата, а снегът продължаваше да вали.
\\
\indentИ валя той още много дни и нощи, затрупвайки всичко под себе си.
\\
\indentТози сняг не бе почнал да се сипе върху угнетената българска земя от Къкринското ханче. Той валеше от по-рано, засипваше селищата и пътищата и натрупваше големи преспи. Отдавна валеше този сняг, отпреди петстотин години… ту монотонен и тих, ту лепкав и мокър, ту усукващ се като някой дервиш от студения вятър, ту ситен и синкав, носен на юруш от злите виелици. И като валя непрекъснато през всичките тези петстотин години, натрупаха се величествени преспи в душата на угнетения българин.
\\
\indentКак можеше през тези преспи робът да види, че при Къкринското ханче пада неговият апостол, отпечатвайки се върху снега! А и самото ханче как можеше да види роба, когато то лежеше дълбоко под преспите на душата му, където всичко е вцепенено в своята зимна летаргия! Само сокерицата има чудната способност да вижда през преспите, това го ние отбелязахме! Но угнетеният народ не е сокерица и как тогава той може да види всичко това, още повече че ханчето не припълзя там до пътя за изстрадалия народ, а за неговия апостол пропълзя през снеговете на времето, тъй както го отбелязахме още в началото, за да може съдбата, побутвайки леко Апостола, да го отбие в ханчето, да хлопне тихо вратата зад гърба му, та по този начин човек и ханче да свържат завинаги съдбите си, не, ами да ги вплетат като зъл възел завинаги!
\\
\indentЕто, така е било!… Пусто и безнадеждно!
\\
\indentДнес, когато надниквам в душата си, виждам в нейния здрач да се белеят, или по-точно, да се сивеят стари преспи. Зная, че това са останки от оня сняг, натрупван в душата на българина през робството. Под тях лежи засипан и снежният отпечатък на българския апостол, и онези бели меджидиета, разхвърляни из снега през глухата нощ. Като разравям студените останки в душата си, аз понякога се взирам напрегнато и търся дали там, под снега, няма да лъсне някое бяло меджидие. За какво ми е това меджидие, не зная и не мога да обясня, а се ровя и търся в посивелите преспи, разядени от времето, станали вече уродливи. Грозни са тия преспи, истина е, уродливият им вид потиска, но няма как другояче да се освободя от тях, освен като ги топя. Затова непрекъснато се ровя из тях и се ослушвам, за да чувам как всяка снежна капка се отцежда мъчително и как тежко и глухо пада на дъното на душата ми.

\subsubsection{Човекът между социалната несправедливост и вечните добродетели (Елин Пелин)}
\vspace{0.5cm}
\begin{center}
\textbf{ЧАСТ 3 (Време за работа: 120 минути)}
\vspace{0.2cm}
{\textit{\underline{Аргументативния текст напишете в листа за отговори на предвиденото място.}}}
\vspace{0.1cm}
\end{center}
\textbf{41.} Напишете аргументативен текст в обем до 4 страници по ЕДНА от следните две теми:
\\
\vspace{0.2cm}
\textbf{Тема за интерпретативно съчинение:}
\vspace{0.1cm}
\begin{center}
\textbf{"Човекът между социалната несправедливост и вечните добродетели"}
\\
\textit{(Интерпретативно съчинение върху „Андрешко“ на Елин Пелин)}
\end{center}
\indent— Какво е това ужасно клепало, дявол да го вземе! Не ме оставя ня мира…
\\
— Дъската и хлопа на моята каруца, господине. Плеще като учен човек: ни сама разбира, ни другите й разбират.
\\
— Хитрец си ти, Андрешко, хитрец си! Трябва да умееш да лъжеш момите, ако не си се оженил. Вие се жените млади и хубави женички имате.
\\
\indentГосподинът свали яката на кожуха си.
\\
— Думай, каквото искаш, господине, ама госпожите по си ги бива… Аз знам това много добре! Ти какъв си, господине, по каква работа към наше село?
\\
— Аз съм съдия-изпълнител.
\\
\indentАндрешко се обърна и с внимание погледна своя кираджия.
\\
— И по изпълнение, а?
\\
— По изпълнение, разбира се. Играе ме един вашенец там, ама тоя път ще му дам да разбере. Няколко пъти го гоня, а все се изплъзва от ръцете ми… Разбрах, че дяволува, и ще го издебна довечера, та ще ме помни… Житцето му ще секвестирам! Хем него на ум ще науча, хем да ви покажа един пример, та да не кръшкате друг път. Лъжете търговците, лъжете гражданите, та им продавате развалени яйца и граниво масло. Е, ама чакай, жено селяшка — властта се лесно не лъже! Пипа тя — здраво пипа! За вас камшик, руски камшик — така ще се оправите… Станахте пияници, пропаднахте, развалихте се — ще станете неспособни данъкоплатци и ще съсипете държавата. Ах, защо нямам повечко власт! Ангели ще ви направя.
\\
\indentСъдията-изпълнител разкопча кожуха си и тялото му зашава в него като пиле, което се излюпва.
\\
— Е, господин съдия, господ направил света и пресметнал, че на жените брада не трябва, и не им дал… Сметнал, че на магарето му трябват дълги уши, и му ги дал — отговори Андрешко с престорена наивност.
\\
— Хем ти не дрънкай, ами карай, че вече се мръкна… Много ми взе ти, дяволе! Толкоз пари за двадесет километра! Умеете да ни скубете вие… Карай де, карай, че крантите заспаха!
\\
— Дий, ди-и, господа! — извика Андрешко и завъртя камшика.
\\
— Ти господа ги наричаш, а? По-добре наричай ги братя — забележи сърдито съдията.
\\
— Ще се разсърдят, господин съдия! Ще ги докача на чест, ако не им кажа господа… Службата им е чиновнишка — на часове. На час стават, на час лягат, на час ги поим, на час ги храним. После ги впрягаме, влизат, значи, в канцеларията, и толкоз. А понякогаж в яслите и вестник прочитат.
\\
— Я кажи де сръбна, приятелю, па не дрънкай, ами карай, че закъсняхме. Лукави очи имаш ти, лукави!
\\
— Няма вълци, господин съдия, не бойте се — рече каруцарят с такъв тон, че почтеният чиновник се озърна плахо наоколо.
\\
— От вълци аз не се боя, приятелю, ами стана студено. Нямам време да настивам.
\\
— Увийте се с чула, господине, Моите коне никога не се оплакват от настинка. Чулът много топли.
\\
\indent„Ама лапацало, а!“ — помисли съдията и се обърна строго:
\\
— Карай, карай, говедо!
\\
\indentИ като се наду сърдито, потъна в кожуха си и млъкна.
\\
\indent„А-а, попаднал си в ръце, приятелю“ — помисли си Андрешко И като се обърна, попита със сериозен тон:
\\
— Та по изпълнение, а? Кому ше изгорите душицата?
\\
\indentСъдията дълго мълча, после сърдито отговори:
\\
— Има един — Станойчо го казват — нисък такъв… с дебел врат…
\\
— Знам го. Нему ще секвестирате житото, а? Сиромах е той, господин съдия, не закачайте го.
\\
— Сиромах човек — жив дявол!
\\
\indentСъдията пак млъкна. Вече се смрачаваше. Конете едвам лазеха въз баирчинката, зад която трябваше ла се намира селото. Андрешко нито им вече подвикваше, нито махаше камшика над тях. Той престана да говори и да си тананика и се размисли.
\\
\indentКогато прехвърлиха зад баира и слязоха отвъд в полето, нощта беше настъпила, а селото още се не виждаше. Тънък и студен ветрец повея над потъналата във влага земя. Облаците, разпокъсани, се оттегляха към планините. Синият купол на студеното и помръзнало небе се изясняваше, разширяваше и възвишаваше. Скоро по него затрепкаха звезди, ледени светли звезди. Въздухът изстина страшно. А конете вървяха бавно. Съдията непрестанно се ядосваше.
\\
— Карай, хей, шоп! Ще помръзнем!
\\
\indentАндрешко безучастно подвикна на конете и лениво размаха камшика над главите им. Те нехайно и омърлушено влачеха каруцата, сякаш не чуваха. Андрешко мислеше за бедния Станоя, комуто съдията утре е секвестира житото, съдията, когото той сега караше.
\\
\indent„Ти ми докара тая беля, Андрешко.“ — ще му каже Станоя, като се научи, и ще го изпсува. После ще се разтъжи, ще го почерпи, ще се напие и ще плаче.
\\
\indent„Трябва да му се помогне на човека, трябва да му се помогне… — мисли Андрешко. — Трябва да му се каже да скрие житцето през нощта и да омете хамбара, инак цяла година ще опъва уши от глад… Трябва да му се помогне — не може инак!“
\\
\indentБеше тъмно и по земята не личеше нищо освен кал, дълбока и гъста кал. Пътят се губеше в тая кал и не водеше никъде, освен пак в нея.
\\
\indentНа едно място Андрешко дръпна юздите и спря конете:
\\
— Чакай, като че съм сбъркал пътя!
\\
\indentИ момъкът почна да се взира в тъмнината. Съдията гледаше строгото му лице, по което нямаше ни следа от прежната шеговитост, и му думаше:
\\
— Момче, опичай си ума, че не отговарям… Бой ще ядеш!
\\
\indentАндрешко дръпна юздите, завъртя камшика и викна:
\\
— Дръжте се здраво, господин съдия!

\subsubsection{Житейските избори на човека (Христо Смирненски)}
\vspace{0.5cm}
\begin{center}
\textbf{ЧАСТ 3 (Време за работа: 120 минути)}
\vspace{0.2cm}
{\textit{\underline{Аргументативния текст напишете в листа за отговори на предвиденото място.}}}
\vspace{0.1cm}
\end{center}
\textbf{41.} Напишете аргументативен текст в обем до 4 страници по ЕДНА от следните две теми:
\\
\vspace{0.2cm}
\textbf{Тема за интерпретативно съчинение:}
\vspace{0.1cm}
\begin{center}
\textbf{"Житейските избори на човека"}
\\
\textit{(Интерпретативно съчинение върху „Приказка за стълбата“ на Христо Смирненски)}
\end{center}
\begin{flushright}
    \textit{Посветено на всички, които ще кажат: 
    \\ "Това не се отнася до мене!".}
\end{flushright}
— Кой си ти? — попита го Дяволът…
\\
— Аз съм плебей по рождение и всички дрипльовци са мои братя. О, колко грозна е земята и колко са нещастни хората!
\\
\indentТова говореше млад момък, с изправено чело и стиснати юмруци. Той стоеше пред стълбата — висока стълба от бял мрамор с розови жилки. Погледът му бе стрелнат в далечината, дето като мътни вълни на придошла река шумяха сивите тълпи на мизерията. Те се вълнуваха, кипваха мигом, вдигаха гора от сухи черни ръце, гръм от негодувание и яростни викове разлюляваха въздуха и ехото замираше бавно, тържествено като далечни топовни гърмежи. Тълпите растяха, идеха в облаци жълт прах, отделни силуети все по-ясно и по-ясно се изрязваха на общия сив фон. Идеше някакъв старец, приведен ниско доземи, сякаш търсеше изгубената си младост. За дрипавата му дреха се държеше босоного момиченце и гледаше високата стълба с кротки, сини като метличина очи. Гледаше и се усмихваше. А след тях идеха все одрипели, сиви, сухи фигури и в хор пееха протегната, погребална песен. Някой остро свиреше с уста, друг, пъхнал ръце в джобовете, се смееше високо, дрезгаво, а в очите му гореше безумие.
\\
— Аз съм плебей по рождение и всички дрипльовци ми са братя! О, колко грозна е земята и колко са нещастни хората! О, вие там горе, вие…
\\
\indentТова говореше млад момък, с изправено чело и стиснати в закана юмруци.
\\
— Вие мразите ония горе? — попита дяволът и лукаво се приведе към момъка.
\\
— О, аз ще отмъстя на тия принцове и князе. Жестоко ще им отмъстя зарад братята ми, зарад моите братя, които имат лица, жълти като пясък, които стенат по-зловещо от декемврийските виелици! Виж голите им кървави меса, чуй стоновете им! Аз ще отмъстя за тях! Пусни ме!
\\
\indentДяволът се усмихна:
\\
— Аз съм страж на ония горе и без подкуп няма да ги предам.
\\
— Аз нямам злато, аз нямам нищо, с което да те подкупя… Аз съм беден, дрипав юноша… Но аз съм готов да сложа главата си.
\\
\indentДяволът пак се усмихна:
\\
— О, не, аз не искам толкоз много! Дай ми ти само слуха си!
\\
— Слуха си? С удоволствие… Нека никога нищо не чуя, нека…
\\
— Ти пак ще чуваш! — успокои го Дяволът и му стори път. — Мини!
\\
\indentМомъкът се затече, наведнъж прекрачи три стъпала, но косматата ръка на Дявола го дръпна:
\\
— Стига! Спри да чуеш как стенат там доле твоите братя!
\\
\indentМомъкът спря и се вслуша:
\\
— Странно, защо те започнаха изведнъж да пеят весело и тъй безгрижно да се смеят!… — И той пак се затече.
\\
\indentДяволът пак го спря:
\\
— За да минеш още три стъпала, аз искам очите ти!
\\
\indentМомъкът отчаяно махна ръка.
\\
— Но тогава аз няма да мога да виждам нито своите братя, нито тия, на които отивам да отмъстя!
\\
\indentДяволът:
\\
— Ти пак ще виждаш… Аз ще ти дам други, много по-хубави очи!
\\
\indentМомъкът мина още три стъпала и се вгледа надоле.
\\
\indentДяволът му напомни:
\\
— Виж голите им кървави меса.
\\
— Боже мой! Та това е тъй странно; кога успяха да се облекат толкоз хубаво! А вместо кървави рани те са обкичени с чудно алени рози!
\\
\indentПрез всеки три стъпала Дяволът взимаше своя малък откуп. Но момъкът вървеше, той даваше с готовност всичко, стига да стигне там и да отмъсти на тия тлъсти князе и принцове! Ето едно стъпало, само още едно стъпало, и той ще бъде горе! Той ще отмъсти зарад братята си!
\\
— Аз съм плебей по рождение и всички дрипльовци…
\\
— Млади момко, едно стъпало още! Само още едно стъпало, и ти ще отмъстиш. Но аз винаги за това стъпало вземам двоен откуп: дай ми сърцето и паметта си.
\\
\indentМомъкът махна ръка:
\\
— Сърцето ли? Не! Това е много жестоко!
\\
\indentДяволът се засмя гърлесто, авторитетно:
\\
— Аз не съм толкова жесток. Аз ще ти дам в замяна златно сърце и нова памет! Ако не приемеш, ти никога няма да минеш туй стъпало, никога няма да отмъстиш за братята си — тия, които имат лица като пясък и стенат по-зловещо от декемврийските виелици. Юношата погледна зелените иронични очи на Дявола:
\\
— Но аз ще бъда най-нещастният. Ти ми взимаш всичко човешко.
\\
— Напротив — най-щастливият! … Но? Съгласен ли си: само сърцето и паметта си.
\\
\indentМомъкът се замисли, черна сянка легна на лицето му, по сбръчканото чело се отрониха мътни капки пот, той гневно сви юмруци и процеди през зъби:
\\
— Да бъде! Вземи ги!
\\
\indent...И като лятна буря, гневен и сърдит, разветрил черни коси, той мина последното стъпало. Той беше вече най-горе. И изведнъж в лицето му грейна усмивка, очите му заблестяха с тиха радост и юмруците му се отпуснаха. Той погледна пируващите князе, погледна доле, дето ревеше и проклинаше сивата дрипава тълпа. Погледна, но нито мускул не трепна по лицето му: то бе светло, весело, доволно. Той виждаше доле празнично облечени тълпи, стоновете бяха вече химни.
\\
— Кой си ти? — дрезгаво и лукаво го попита Дяволът.
\\
— Аз съм принц по рождение и боговете ми са братя! О, колко красива е земята и колко са щастливи хората!

\subsubsection{Животът, любовта и смъртта (Христо Ботев)}
\vspace{0.5cm}
\begin{center}
\textbf{ЧАСТ 3 (Време за работа: 120 минути)}
\vspace{0.2cm}
{\textit{\underline{Аргументативния текст напишете в листа за отговори на предвиденото място.}}}
\vspace{0.1cm}
\end{center}
\textbf{41.} Напишете аргументативен текст в обем до 4 страници по ЕДНА от следните две теми:
\\
\vspace{0.2cm}
\textbf{Тема за интерпретативно съчинение:}
\vspace{0.1cm}
\begin{center}
\textbf{"Животът, любовта и смъртта"}
\\
\textit{(Интерпретативно съчинение върху „До моето първо либе“ на Христо Ботев)}
\end{center}

\begin{tabular}{p{0.48\textwidth} p{0.48\textwidth}}
\begin{minipage}[t]{\linewidth}
Остави таз песен любовна,\\
не вливай ми в сърце отрова, –\\
млад съм аз, но младост не помня,\\
пък и да помня, не ровя\\
туй, що съм ази намразил\\
и пред тебе с крака погазил.\\
\\
Забрави туй време, га плачех\\
за поглед мил и за въздишка:\\
роб бях тогаз – вериги влачех,\\
та за една твоя усмивка,\\
безумен аз светът презирах\\
и чувства си в калта увирах!\\
\\
Забрави ти онез полуди,\\
в тез гърди веч любов не грее\\
и не можеш я ти събуди\\
там, де скръб дълбока владее,\\
де всичко е с рани покрито\\
и сърце зло в злоба обвито!\\
\\
Ти имаш глас чуден – млада си,\\
но чуйш ли как пее гората?\\
Чуйш ли как плачат сиромаси?\\
За тоз глас ми копней душата,\\
и там тегли сърце ранено,\\
там, де е се с кърви облено!\\
\\
О, махни тез думи отровни!\\
Чуй как стене гора и шума,\\
чуй как ечат бури вековни,\\
как нареждат дума по дума –\\
приказки за стари времена\\
и песни за нови теглила!\\
\end{minipage}
&
\begin{minipage}[t]{\linewidth}
Запей и ти песен такава,\\
запей ми, девойко, на жалост,\\
запей как брат брата продава,\\
как гинат сили и младост,\\
как плаче сиротна вдовица\\
и как теглят без дом дечица!\\
\\
Запей или млъкни, махни се!\\
Сърце ми веч трепти – ще хвръкне,\\
ще хвръкне, изгоро, – свести се!\\
Там, де земя гърми и тътне\\
от викове страшни и злобни\\
и предсмъртни песни надгробни...\\
\\
Там... там буря кърши клонове,\\
а сабя ги свива на венец;\\
зинали са страшни долове\\
и пищи в тях зърно от свинец,\\
и смъртта ѝ там мила усмивка,\\
а хладен гроб сладка почивка!\\
\\
Ах, тез песни и таз усмивка\\
кой глас ще ми викне, запее? –\\
Кървава да вдигна напивка,\\
от коя и любов немее,\\
пък тогаз и сам ще запея\\
що любя и за що милея!..\\
\end{minipage}
\end{tabular}

\subsubsection{Геройство и памет (Иван Вазов)}
\vspace{0.5cm}
\begin{center}
\textbf{ЧАСТ 3 (Време за работа: 120 минути)}
\vspace{0.2cm}
{\textit{\underline{Аргументативния текст напишете в листа за отговори на предвиденото място.}}}
\vspace{0.1cm}
\end{center}
\textbf{41.} Напишете аргументативен текст в обем до 4 страници по ЕДНА от следните две теми:
\\
\vspace{0.2cm}
\textbf{Тема за интерпретативно съчинение:}
\vspace{0.1cm}
\begin{center}
\textbf{"Геройство и памет"}
\\
\textit{(Интерпретативно съчинение върху „Новото гробище над Сливница“ на Иван Вазов)}
\end{center}
\begin{tabular}{p{0.48\textwidth} p{0.48\textwidth}}
\begin{minipage}[t]{\linewidth}
Покойници, вий в други полк минахте,\\
де няма отпуск, ни зов за борба,\\
вий братски се прегърнахте, легнахте\\
и „Лека нощ“ навеки си казахте —\\
\hspace*{1.5cm} до втората тръба.\\
\\
Но що паднахте тук, деца бурливи?\\
За трон ли злат, за някой ли кумир?\\
Да беше то — остали бихте живи,\\
не бихте срещали тъй горделиви\\
\hspace*{1.5cm} куршума… Спете в мир.\\
\\
Българио, за тебе те умряха,\\
една бе ти достойна зарад тях,\\
и те за теб достойни, майко, бяха\\
И твойто име само кат мълвяха,\\
\hspace*{1.5cm} умираха без страх.\\
\\
Но кой ви знай, че спите в тез полета?\\
Над ваший гроб забвеньето цъфти.\\
Кои сте вий? Над сянката ви клета\\
не мисли никой днес освен поета\\
\hspace*{1.5cm} и майките свети.\\
\end{minipage}
&
\begin{minipage}[t]{\linewidth}
Борци, венец ви свих от песен жива,\\
от звукове, що никой не сбира:\\
от дивий рев на битката гръмлива,\\
от екота на Витоша бурлива,\\
\hspace*{1.5cm} от вашето ура.\\
\\
И тоз венец — той няма да завене,\\
и тая песен вечно ще гърми\\
из българските планини зелени,\\
и славата ще вечно пей и стене\\
\hspace*{1.5cm} над гробни ви хълми.\\
\\
Почивайте под тез могили ледни:\\
не ще да чуйте веч тръба, ни вожд,\\
ни славний гръм на битките победни,\\
към вечността е маршът ви последни.\\
\hspace*{1.5cm} Юнаци, лека нощ!\\
\end{minipage}
\end{tabular}

\subsubsection{Природата като храм на душата (Иван Вазов)}
\vspace{0.5cm}
\begin{center}
\textbf{ЧАСТ 3 (Време за работа: 120 минути)}
\vspace{0.2cm}
{\textit{\underline{Аргументативния текст напишете в листа за отговори на предвиденото място.}}}
\vspace{0.1cm}
\end{center}
\textbf{41.} Напишете аргументативен текст в обем до 4 страници по ЕДНА от следните две теми:
\\
\vspace{0.2cm}
\textbf{Тема за интерпретативно съчинение:}
\vspace{0.1cm}
\begin{center}
\textbf{"Природата като храм на душата"}
\\
\textit{(Интерпретативно съчинение върху „При Рилския манастир“ на Иван Вазов)}
\end{center}
\begin{tabular}{p{0.48\textwidth} p{0.48\textwidth}}
\begin{minipage}[t]{\linewidth}
Сега съм у дома. Наокол планини\\
и върхове стърчат; гори високи, диви\\
шумят; потоците, кристални и пенливи,\\
бучат — живот кипи на всичките страни.\\
Природата отвред, кат майка нежна съща,\\
напява ми песнѝ, любовно ме прегръща.\\
\\
Сега съм у дома. Над мен Еленин връх\\
боде лазурний свод и мен при себе кани;\\
отсреща Бричебор ми праща здравий дъх\\
на своите ели и бори великани;\\
а Царев връх от юг издига се огромен\\
с плешивия си лоб и царския си спомен.\\
\\
Сега съм у дома, сега съм в моя мир —\\
мир въжделен и драг. Тук волно дишам ази,\\
по-светло чувствувам; свещен, отраден мир\\
изпълни ми духа, от нов живот талази\\
нахлуват в мен, трептя от нови ощущенья,\\
от прясна сила, мощ и тайни песнопенья…\\
\\
Сега съм у дома, сега съм пак поет —\\
във лоното на таз пустиня горска, свята;\\
разбирам на леса любовний тих привет,\\
на струите шума, на бездната мълвата.\\
Разменям тайна реч с земя и синий свод\\
и сливам се честит във техния живот.\\
\end{minipage}
&
\begin{minipage}[t]{\linewidth}
Сега съм у дома, в сърцето съм на Рила.\\
Световните злини, тревоги са далеч —\\
за тях е тя стена до небеса турила —\\
усещам се добър, почти невинен веч.\\
Духът ми се цери след жизнената битва,\\
вкушавам сладък мир във песни и молитва.\\
\\
Сега съм у дома. По часове, благат,\\
край бистрата река, при звучната й песен,\\
мечтая ил чета… Ил кат орел надвесен\\
над бездните стоя и моят ум хвъркат\\
блуждае в хаоса, до господа отива,\\
на мирозданьето във тайните се впива.\\
\\
Сега съм у дома — не съм тук странен гост.\\
Природата всегда, но буйната природа,\\
що пълни я живот, шум, песен и свобода,\\
бе моят идеал величествен и прост.\\
Поклон, скали, води! Поклон, ели гигантски!\\
Вам, бездни, висоти! Вам, гледки великански!\\
\\
Сега съм у дома — участник в рилский хор.\\
Аз тук не се родих — тук бих желал да тлея —\\
под горский вечен шум — дълбока епопея —\\
и на Еленин връх под вечно будний взор;\\
да имам гроб, подир живот-синджир теглила,\\
\hspace*{1.0cm} в величествените обятия на Рила.\\
\end{minipage}
\end{tabular}

\subsubsection{Човешкия стремеж към хармония (Пенчо Славейков)}
\vspace{0.5cm}
\begin{center}
\textbf{ЧАСТ 3 (Време за работа: 120 минути)}
\vspace{0.2cm}
{\textit{\underline{Аргументативния текст напишете в листа за отговори на предвиденото място.}}}
\vspace{0.1cm}
\end{center}
\textbf{41.} Напишете аргументативен текст в обем до 4 страници по ЕДНА от следните две теми:
\\
\vspace{0.2cm}
\textbf{Тема за интерпретативно съчинение:}
\vspace{0.1cm}
\begin{center}
\textbf{"Човешкия стремеж към хармония"}
\\
\textit{(Интерпретативно съчинение върху „Спи езерото“ на Пенчо Славейков)}
\end{center}
\noindentСпи езерото; белостволи буки\\
над него свождат вити гранки,\\
и в тихите му тъмни глъбини\\
преплитат отразени сянки.\\
\\
Треперят, шептят белостволи буки,\\
а то, замряло, нито трепва…\\
Понякога му сал повърхнини\\
дълга от лист отронен сепва.\\

\subsubsection{Животът е кратък, изкуството е вечно (Виктор Пасков)}
\vspace{0.5cm}
\begin{center}
\textbf{ЧАСТ 3 (Време за работа: 120 минути)}
\vspace{0.2cm}
{\textit{\underline{Аргументативния текст напишете в листа за отговори на предвиденото място.}}}
\vspace{0.1cm}
\end{center}
\textbf{41.} Напишете аргументативен текст в обем до 4 страници по ЕДНА от следните две теми:
\\
\vspace{0.2cm}
\textbf{Тема за интерпретативно съчинение:}
\vspace{0.1cm}
\begin{center}
\textbf{"Животът е кратък, изкуството е вечно"}
\\
\textit{(Интерпретативно съчинение върху „Балада за Георг Хених“ на Виктор Пасков)}
\end{center}
\indentДокато баща ми тичаше по репетиции, представления, полунощни чалги, докато обикаляхме комисии и инстанции, докато някой се бесеше, друг умираше, трети се готвеше да бяга в Америка и докато майка ми пестеше всяка стотинка, за да пълни бюфета си със сервизи, докато на нашето семейство започваха да гледат с все по-лошо око, защото бяха научили, че се каним да приберем някакъв старец — старецът довършваше своята цигулка, вън от този хаотичен свят, едва люлеещ се на краката си, глух за всичко и всеки.
\\
\indentМного трудно бе да го накараш да влезе в разговор. Не ядеше. Пушеше. Изпиваше огромни количества вода. Стискаше тънки устни и мълчеше. Понякога изпадаше в старото си вцепенение — седеше на миндера, гледаше в ъгъла и мърмореше нещо в брадата си, но това не траеше дълго. Най-много четвърт час; тогава ставаше и започваше отново да работи.
\\
\indentКогато почиваше, аз се примъквах, вдигах кърпата, с която завиваше инструмента, и го разглеждах.
\\
\indentБоже, господи! Каква ти цигулка! Има осминки, четвъртинки, три четвърти и цели; вече бях виждал всякакви, но не и такава. Това бе най-голямата цигулка, която някога е съществувала… сигурно. Може би по-голяма и от виола. Не че бе кой знае колко дълга — дължината й не надвишаваше значително тази на виолата, но бе много по-дълбока, следователно душичката й бе по-голяма. Най-голямата душичка, която някога е поставяна в цигулка.
\\
\indentГрифът се съединяваше с охлювчето направо. Охлювче, каквото сме свикнали да виждаме на цигулките, нямаше. Георг Хених бе изваял от дървото глава на жена. Неговата цигулка завършваше — или започваше? — с женска глава, изработена с най-тънкото длето от комплекта и с изключително майсторство. Главата удивително напомняше изтритата снимка на Боженка, която висеше на стената, забита с пирон. Същевременно — косите, скръбно наведената надолу глава, скулестото лице и мекотата, което излъчваше то, напомняха силно божията майка от иконата, аве Мария, грация плена! — очите на жената бяха привързани с превръзка, чийто възел стоеше отзад на главата й като кок.
\\
— Харесва тебе?
\\
— Много! Прилича на майката на Исус. Също и на жена ти. Само че — малко по-стара. Каква е тази глава?
\\
— Амор. На булгарски каже се — любов. Малък, не разбира. Инфант.
\\
— Защо да не разбирам? Любов, ясно. А защо е с превързани очи?
\\
— Защо любов — въздъхваше майстор Хених — сляпа. Не разбира.
\\
— Много е тъжна главата.
\\
— Не тъжна. Кога свири музикант, гледа право у глава. Свири първо за глава, после за други. Глава чува всичко, разбира всичко. Ако нема любов у гърди, ако не гледа любов у лице — музикант не майстор. Може виртуоз, але не майстор. Разбира? Глава слепа, а гледа музикант право у сърце. Але ако нема сърце?
\\
— Нали е за господ?
\\
— Молим да. За господ. Остави ме работи. Гледа, мисли.
\\
* * * *
\\
\indentНяколко дни след разговора с баща ми майстор Франтишек пристигна при Георг Хених. Беше слънчев следобед. Дори в мазето на стареца няколко лъча бяха успели да намерят пролука и весело се боричкаха на пода. Настроението му тази сутрин бе приповдигнато — дали защото работата отиваше към своя край, или заради слънчевите лъчи, но дядо Георги се шегуваше, работейки последните детайли върху главата и наглеждайки от време на време лепилото в металния съд.
\\
\indentРазказваше пак за Чехия. За буйния си брат Антон, за който нямало по-голямо удоволствие от това да препуска на кон из зелените ливади, облечен като древен рицар, с меч на кръста и валдхорна в ръце, и как препускайки, той можел да изпълни и най-сложните пасажи на своята валдхорна, тези от увертюрата на „Вълшебния стрелец“, та дори и концерта на Моцарт — и дядо Георги отпускаше грифа и се опитваше да тананика този концерт. Гласът му бе ръждив и писклив и с нищо не напомняше бляскавите звуци на славния инструмент, на който някога е свирил брат му. В края на краищата той забелязваше, че едва потискам смеха си и сам започваше да се смее ситно, друсайки се цял. Смехът му преминаваше в кашлица и тогава аз го тупах по гърба, докато се успокои.
\\
\indentНа вратата се почука силно. Някой опита веднага след това да я отвори, но Георг Хених работеше вече само при закрити врати: райберът беше пуснат, за да не го изненада комисията.
\\
— Кой? — попита старецът и се ослуша.
\\
— Старче, отваряй? Аз съм, Франта. Или си ме забрави вече, а? Само че аз не те забравям, нищо че ме изгони последния път. Хайде.
\\
— Отива си! Нема време. Не продава инструмент!
\\
— Дръж си инструментите, не ги искам. Разбрах, че си ги подарил.
\\
— Кой тебе каза?
\\
— И това ли си забравил? Нали ти изпрати онези двамата, бащата и детето, при мен. Те ми казаха.
\\
— Кого изпратил? Никого не пратил! Върви! Стари Хених има работа. Отивай, Франта!
\\
— Може ли така? Толкова години си ме учил, може ли така да ме отпращаш? Хайде, пусни Франтишек да види каква цигулка правиш. Чия е поръчката?
\\
— За цигулка знае?
\\
— Детето ми каза. Хайде, стара чешка скръндзо! Покажи ми да я видя! Казах на Ванко, че пак си почнал работа. Той също иска да намине. Какво пък, нали сега сме конкуренция? — и Франтишек гърлено се изсмя, не преставайки да натиска бравата.
\\
— Отива си! Казал. Цигулка не готова.
\\
Майстор Франтишек изруга зад вратата и престана да напира. Съседната врата хлопна. Чухме гласовете на Рижия и жена му.
\\
Георг Хених отпусна грифа в ръцете си. Женската глава докосваше пода.
\\
— Каже мене защо с баща при Франтишек?
\\
— Ето какво — започнах да обяснявам, — ако те с майстор Ванда се съгласят да дават по двадесет лева издръжка, тогава и ние ще даваме и ще те оставят на мира. Никой няма да те закача.
\\
— Баща ходил при Франтишек? Дваес лева за Хених.
\\
— Ами да! Нали са ти ученици?
\\
Старецът унило сложи цигулката на тезгяха и покри лице с едрите си шепи. Мълча дълго. Усещах, че нещо не е наред. Чувствах се неловко и не знаех защо. Нали всичко е за негово добро? Какво по-хубаво от това — да остане при нас? Защо не разбираше?
\\
— Отиде си сега, уморен, иска малко почине.
\\
\indentГонеше ли ме? Никога не ме бе отпращал. Когато искаше да си почине, просто лягаше на кревата и се обръщаше с гръб. Аз се занимавах, разглеждах инструментите, не пречех. С какво го обидих?
\\
— Дядо Георги, нека да остана! Ти легни, аз ще почакам. Като се събудиш, пак ще си говорим. Нали обеща да ми разкажеш за онази цигулка на Гаспаро да Сало? Дето я украсил Бенвенуто Челини? Моля ти се!
\\
— Не. Отиде. Дойде утре, вдругиден. Иска бъде сам.
\\
\indentНямаше как. Станах и си тръгнах. Той ме изпрати до вратата. Промърмори „довиждане“, без да ме погледне, и заключи след мен.

\subsubsection{Любовта между реалността и спомена (Димчо Дебелянов)}
\vspace{0.5cm}
\begin{center}
\textbf{ЧАСТ 3 (Време за работа: 120 минути)}
\vspace{0.2cm}
{\textit{\underline{Аргументативния текст напишете в листа за отговори на предвиденото място.}}}
\vspace{0.1cm}
\end{center}
\textbf{41.} Напишете аргументативен текст в обем до 4 страници по ЕДНА от следните две теми:
\\
\vspace{0.2cm}
\textbf{Тема за интерпретативно съчинение:}
\vspace{0.1cm}
\begin{center}
\textbf{"Любовта между реалността и спомена"}
\\
\textit{(Интерпретативно съчинение върху „Аз искам да те помня все така“ на Димчо Дебелянов)}
\end{center}
\begin{tabular}{p{0.48\textwidth} p{0.48\textwidth}}
\begin{minipage}[t]{\linewidth}
Аз искам да те помня все така:\\
бездомна, безнадеждна и унила,\\
в ръка ми вплела пламнала ръка\\
и до сърце ми скръбен лик склонила.\\
Градът далече тръпне в мътен дим,\\
край нас, на хълма, тръпнат дървесата\\
и любовта ни сякаш по е свята,\\
защото трябва да се разделим.\\
\\
„В зори ще тръгна, ти в зори дойди\\
и донеси ми своя взор прощален —\\
да го припомня верен и печален\\
в часа, когато Тя ще победи!“\\
— О, Морна, Морна, в буря скършен злак,\\
укрий молбите, вярвай — пролетта ни\\
недосънуван сън не ще остане\\
и ти при мене ще се върнеш пак!\\
\end{minipage}
&
\begin{minipage}[t]{\linewidth}
А все по-страшно пада нощ над нас,\\
чертаят мрежа прилепите в мрака,\\
утеха сетна твойта немощ чака,\\
а в свойта вяра сам не вярвам аз.\\
И ти отпущаш пламнала ръка\\
и тръгваш, поглед в тъмнината впила,\\
изгубила дори за сълзи сила. —\\
Аз искам да те помня все така…\\
\end{minipage}
\end{tabular}

\subsubsection{Любовта и красотата – мимолетност и вечност (Христо Фотев)}
\vspace{0.5cm}
\begin{center}
\textbf{ЧАСТ 3 (Време за работа: 120 минути)}
\vspace{0.2cm}
{\textit{\underline{Аргументативния текст напишете в листа за отговори на предвиденото място.}}}
\vspace{0.1cm}
\end{center}
\textbf{41.} Напишете аргументативен текст в обем до 4 страници по ЕДНА от следните две теми:
\\
\vspace{0.2cm}
\textbf{Тема за интерпретативно съчинение:}
\vspace{0.1cm}
\begin{center}
\textbf{"Любовта и красотата – мимолетност и вечност"}
\\
\textit{(Интерпретативно съчинение върху „Колко си хубава!“ на Христо Фотев)}
\end{center}
\begin{tabular}{p{0.48\textwidth} p{0.48\textwidth}}
\begin{minipage}[t]{\linewidth}
На М. К.\\
Колко си хубава!\\
Господи,\\
колко си хубава!\\
\\
Колко са хубави ръцете ти.\\
И нозете ти колко са хубави.\\
И очите ти колко са хубави.\\
И косите ти колко са хубави\\
\\
Не се измъчвай повече — обичай ме!\\
Не се щади — обичай ме!\\
Обичай ме\\
със истинската сила на ръцете си,\\
нозете си, очите си — със цялото\\
изящество на техните движения.\\
Повярвай ми завинаги — и никога\\
ти няма да си глупава — обичай ме!\\
И да си зла — обичай ме!\\
Обичай ме!\\
По улиците, след това по стълбите,\\
особено по стълбите си хубава.\\
Със дрехи и без дрехи, непрекъснато\\
си хубава… Най-хубава си в стаята.\\
Във тъмното, когато си със гребена.\\
И гребенът потъва във косите ти.\\
Косите ти са пълни с електричество —\\
докосна ли ги, ще засветя в тъмното.\\
\end{minipage}
&
\begin{minipage}[t]{\linewidth}
Наистина си хубава — повярвай ми.\\
И се старай до края да си хубава.\\
Не толкова за мене, а за себе си,\\
дърветата, прозорците и хората.\\
Не разрушавай бързо красотата си\\
с ревниви подозрения — прощавай ми\\
внезапните пропадания някъде —\\
не прекалявай, моля те, с цигарите.\\
Не ме изгубвай никога — откривай ме,\\
изпълвай ме с детинско изумление.\\
Отново да се уверя в ръцете ти,\\
в нозете ти, в очите ти… Обичай ме!\\
Как искам да те задържа завинаги.\\
Да те обичам винаги —\\
завинаги.\\
И колко ми е невъзможно… Колко си\\
ти пясъчна… И моля те, не казвай ми,\\
че искаш да ме задържиш завинаги,\\
да ме обичаш винаги,\\
завинаги.\\
\\
Колко си хубава!\\
Господи,\\
колко си хубава!\\
\\
Колко са хубави ръцете ти.\\
И нозете ти колко са хубави.\\
И очите ти колко са хубави.\\
И косите ти колко са хубави.\\
\\
Колко си хубава!\\
Господи,\\
колко си истинска.\\
\end{minipage}
\end{tabular}

\subsubsection{Философията на любовта като израстване (Петя Дубарова)}
\vspace{0.5cm}
\begin{center}
\textbf{ЧАСТ 3 (Време за работа: 120 минути)}
\vspace{0.2cm}
{\textit{\underline{Аргументативния текст напишете в листа за отговори на предвиденото място.}}}
\vspace{0.1cm}
\end{center}
\textbf{41.} Напишете аргументативен текст в обем до 4 страници по ЕДНА от следните две теми:
\\
\vspace{0.2cm}
\textbf{Тема за интерпретативно съчинение:}
\vspace{0.1cm}
\begin{center}
\textbf{"Философията на любовта като израстване"}
\\
\textit{(Интерпретативно съчинение върху „Посвещение“ на Петя Дубарова)}
\end{center}
\begin{tabular}{p{0.48\textwidth} p{0.48\textwidth}}
\begin{minipage}[t]{\linewidth}
В студените нощи, когато пиян\\
сънят се търкаля на моя таван,\\
когато луната тъмнее от грях,\\
когато увисва над мен моя страх,\\
обесен на острия ръб на нощта,\\
подавам ти своята бледа ръка —\\
на теб — непознатия — смугло красив,\\
потаен и питомен, жаден и див,\\
едва деветнайсет години живял,\\
а всичко опитал и всичко видял,\\
подвластен на никого, ничий, сам свой,\\
но тръгнал към мене и истински мой\\
\end{minipage}
&
\begin{minipage}[t]{\linewidth}
и падал по пътя си, плакал, грешил,\\
но нежност момчешка за мен съхранил.\\
Ръката ми — властната — жадно поел,\\
единствено с мен ще си толкова смел!\\
Ела! Ще измием луната от грях!\\
Ще хвърлим трупа на умрелия страх,\\
ще пеем с тътнежния корабен глас\\
на морската нощ във добрия Бургас.\\
А после, когато тя тръгне назад\\
и слънцето бликне над нас благодат,\\
мечтата надрастнал, усмихнат, смутен,\\
ще тръгнеш реален до мен в моя ден!\\
\end{minipage}
\end{tabular}

\subsubsection{Страданието и търсеното спасение (Елин Пелин)}
\vspace{0.5cm}
\begin{center}
\textbf{ЧАСТ 3 (Време за работа: 120 минути)}
\vspace{0.2cm}
{\textit{\underline{Аргументативния текст напишете в листа за отговори на предвиденото място.}}}
\vspace{0.1cm}
\end{center}
\textbf{41.} Напишете аргументативен текст в обем до 4 страници по ЕДНА от следните две теми:
\\
\vspace{0.2cm}
\textbf{Тема за интерпретативно съчинение:}
\vspace{0.1cm}
\begin{center}
\textbf{"Страданието и търсеното спасение"}
\\
\textit{(Интерпретативно съчинение върху „Спасова могила“ на Елин Пелин)}
\end{center}
\indentПък хора, хора! От всички страни, по всички пътеки идат хора един през друг, настигат се, заминават се, все към оная висока, стръмна и островърха могила, Спасова могила. Вижда се кичестият стар дъб на върха и малкото бяло параклисче до него. По голата зелена стръмнина плъзнал свят като мравуняк. Дядо Захари гледа и се чуди. Отде се взе тоя народ, това чудо! С кола, с коне, пеши, отблизо и отдалеч, отвсякъде идат хора и бързат нататък. Между тях какви не, какви не! Бедняци, дрипави, с оголени меса. Богати, пременени празнично. Всички носят по един недъг и по една надежда за изцеление. Някои с пречупени кръстове се влачат като змии, други на кокили, трети с мръсни рани по тялото. Слепи, сакати…
\\
— Дядо, къде отива тоя народ? — пита болният Монка.
\\
— Всички там отиват, чедо.
\\
— Всички болни ли са, дядо?
\\
— То целият свят е болен, синко. Едни от това, Други от онова. Няма здрав човек на света. Гледаш, тялото желязно, а душата — гнила.
\\
\indentОт върха се понесе тихо черковният глас на дървеното клепало и се разля като благословия по цялата зелена околност.
\\
\indentДядо Захари седна уморено край пътя и въздъхна:
\\
— Прекръсти се, чедо!
\\
\indentМонката откопчи ръце от шията на дяда си, па и двамата дълго се кръстиха, седнали в зеления бурен край пътя.
\\
— Дядо, искам да походя.
\\
— Слабичко си, момченцето ми, ще се умориш.
\\
— Искам да походя, бе дядо — настоя плачливо Монката.
\\
\indentДядо Захари го повдигна грижливо и като го улови за ръка, тръгнаха пак. Монката, десетгодишно момче, болнаво и хилаво сираче, тръгна край дяда си като сянка. По бледото му сухо лице се прозираха големи сини жили, краката му се преплитаха от слабост, ръцете му висяха като клечки. Големите му сини очи, някак страшно отворени, се въртяха навред и с едно умилено учудване следяха ту птичките, които мълчаливи прехвръкваха към гнездата си, ту гривите гълъби, които чезнеха във вечерното сияние към тъмната кория на запад, ту златните мушички, които жеркаха като водни капки във въздуха.
\\
\indentДядо Захари вървеше съвсем полека, за да не умо ри малкия, но като видя, че слънцето вече заседна и хората по пътя намаляха, той отново го качи на гръб и тръгна по-бързо.
\\
\indentПо върха на могилката се мяркаха човешки сенки безцелно, като заблудени.
\\
— Хората си търсят вече място, закъсняхме, сине — рече старецът.
\\
— Дядо господ рано ли ще слезне там, дядо? — пита Монката.
\\
— Къде ти рано! Той ще чака да дойдат всички, които са тръгнали. Зер има и отдалече, има сакати, слепи, има такива, дето не знаят пътя, може да се заблуди някой… Като се приберат всички и като заспят всички, дядо господ ще повика едно ангелче да му носи патерицата и ще му каже — хайде! Ще се спуснат мълчаливо там отгоре, измежду звездите, и полека-лека ще слязат долу. Ще мине дядо господ между тоя народ, дето е тук, тихо, никой да го не чуе, никой да го не види, и ще каже изцеление на всички. После полека-лека, както си е дошъл, ще се вдигне и ще си отиде на небето. Също както някога се е възнесъл на планината, те на тоя същия, утрешния ден — Спасовден.
\\
— Болните ще оздравеят ли, дядо?
\\
— Всички — всички! — отговори с дълбока увереност дядо Захари. — Може сега, може след година, след пет, след десет. Който вярва, ще се спаси.

\subsubsection{Сложността и простотата (Атанас Далчев)}
\vspace{0.5cm}
\begin{center}
\textbf{ЧАСТ 3 (Време за работа: 120 минути)}
\vspace{0.2cm}
{\textit{\underline{Аргументативния текст напишете в листа за отговори на предвиденото място.}}}
\vspace{0.1cm}
\end{center}
\textbf{41.} Напишете аргументативен текст в обем до 4 страници по ЕДНА от следните две теми:
\\
\vspace{0.2cm}
\textbf{Тема за интерпретативно съчинение:}
\vspace{0.1cm}
\begin{center}
\textbf{"Сложността и простотата"}
\\
\textit{(Интерпретативно съчинение върху "Молитва“ на Атанас Далчев)}
\end{center}
\begin{tabular}{p{0.48\textwidth} p{0.48\textwidth}}
\begin{minipage}[t]{\linewidth}
Аз не помня, аз не съм видял\\
минаха ли моите години?\\
Ти не ме оставяй да загина,\\
господи, преди да съм живял!\\
\\
Изведи ме вън от всяка сложност,\\
научи ме пак на простота:\\
да отдавам сетния петак\\
от сърце на срещнатия просек.\\
\end{minipage}
&
\begin{minipage}[t]{\linewidth}
Да усещам своя радостта\\
на невинното дете, което\\
първите снежинки от небето\\
сбира със отворена уста.\\
\\
И без свян да мога да говоря\\
с простите на прост неук език…\\
Научи ме, господи велик,\\
да живея като всички хора.\\
\end{minipage}
\end{tabular}

\subsubsection{Вярата - антипод на социалния ад (Никола Вапцаров)}
\vspace{0.5cm}
\begin{center}
\textbf{ЧАСТ 3 (Време за работа: 120 минути)}
\vspace{0.2cm}
{\textit{\underline{Аргументативния текст напишете в листа за отговори на предвиденото място.}}}
\vspace{0.1cm}
\end{center}
\textbf{41.} Напишете аргументативен текст в обем до 4 страници по ЕДНА от следните две теми:
\\
\vspace{0.2cm}
\textbf{Тема за интерпретативно съчинение:}
\vspace{0.1cm}
\begin{center}
\textbf{"Вярата - антипод на социалния ад"}
\\
\textit{(Интерпретативно съчинение върху „Вяра“ на Никола Вапцаров)}
\end{center}
\begin{tabular}{p{0.31\textwidth} p{0.31\textwidth} p{0.31\textwidth}}
\begin{minipage}[t]{\linewidth}
Ето — аз дишам,\\
работя,\\
живея\\
и стихове пиша\\
(тъй както умея).\\
С живота под вежди\\
се гледаме строго\\
и боря се с него,\\
доколкото мога.\\
\\
С живота сме в разпра,\\
но ти не разбирай,\\
че мразя живота.\\
Напротив, напротив! —\\
Дори да умирам,\\
живота със грубите\\
лапи челични\\
аз пак ще обичам!\\
Аз пак ще обичам!\\
\\
Да кажем, сега ми окачат\\
въжето\\
и питат:\\
„Как, искаш ли час да живееш?“\\
Веднага ще кресна:\\
„Свалете!\\
Свалете!\\
По-скоро свалете\\
въжето, злодеи!“\\
\end{minipage}
&
\begin{minipage}[t]{\linewidth}
За него — Живота —\\
направил бих всичко. —\\
Летял бих\\
със пробна машина в небето,\\
бих влезнал във взривна\\
ракета, самичък,\\
бих търсил\\
в простора\\
далечна\\
планета.\\
\\
Но все пак ще чувствам\\
приятния гъдел,\\
да гледам как\\
горе\\
небето синее.\\
Все пак ще чувствам\\
приятния гъдел,\\
че още живея,\\
че още ще бъда.\\
\\
Но ето, да кажем,\\
вий вземете, колко? —\\
пшеничено зърно\\
от моята вера,\\
бих ревнал тогава,\\
бих ревнал от болка\\
като ранена\\
в сърцето пантера.\\
\end{minipage}
&
\begin{minipage}[t]{\linewidth}
Какво ще остане\\
от мене тогава? —\\
Миг след грабежа\\
ще бъда разнищен.\\
И още по-ясно,\\
и още по-право —\\
миг след грабежа\\
ще бъда аз нищо.\\
\\
Може би искате\\
да я сразите\\
моята вяра\\
във дните честити,\\
моята вяра,\\
че утре ще бъде\\
живота по-хубав,\\
живота по-мъдър?\\
\\
А как ще щурмувате, моля?\\
С куршуми?\\
Не! Неуместно!\\
Ресто! — Не струва! —\\
Тя е бронирана\\
здраво в гърдите\\
и бронебойни патрони\\
за нея\\
няма открити!\\
Няма открити!\\
\end{minipage}
\end{tabular}

\subsubsection{„Ветрената мелница“ – разказ за порива и жизнелюбието на българина (Елин Пелин)}
\vspace{0.5cm}
\begin{center}
\textbf{ЧАСТ 3 (Време за работа: 120 минути)}
\vspace{0.2cm}
{\textit{\underline{Аргументативния текст напишете в листа за отговори на предвиденото място.}}}
\vspace{0.1cm}
\end{center}
\textbf{41.} Напишете аргументативен текст в обем до 4 страници по ЕДНА от следните две теми:
\\
\vspace{0.2cm}
\textbf{Тема за интерпретативно съчинение:}
\vspace{0.1cm}
\begin{center}
\textbf{"„Ветрената мелница“ – разказ за порива и жизнелюбието на българина"}
\\
\textit{(Интерпретативно съчинение върху „Ветрената мелница“ на Елин Пелин)}
\end{center}
\indentГайдата писна. Навалицата притаи дъх.
\\
\indentИграчите застанаха един срещу други и едновременно почнаха.
\\
\indentДъбакът чевръсто застъпя срещу Христина. Тя се затича леко на пръсти и го замина. Двамата се изгледаха от глава до крака, като че искаха да покажат надмощието си един пред друг, и заиграха един срещу други. Христина размаха с кърпата, изви като лебед бялата си шийка и леко се понесе във вихъра на веселата свирня. Лицето й пламна, унесено се полусклопиха босилковите й клепки, завълнуваха се облите й гърди.
\\
\indentДъбакът пък съвсем беше се забравил. С ръце, нехайно кръстосани отдире, той подскачаше като елен, заставаше с чудни движения срещу противницата си и заситняше, та сякаш се поколваше на едно място. После, като отърсваше глава да окапе едрата пот, която обливаше лицето му, той правеше движение, като че ще скочи назад. Христина, която дебнеше всяка негова стъпка, излеко и на пръсти политаше към него, но той подскачаше напред и тя неволно се намереше толкова близко до него, че усещаше дъха му и топлината от пламналото му мъжко лице. Той полека-лека се отдалечаваше и неусетно я увличаше по себе си, като че се шегуваше. Едно надмощие, едно мъжко надмощие личеше във всяко негово движение и Христина неволно падаше под силата му. Унесена, тя се мъчеше да победи, но след едночасова борба спря отмаляла и запъхтяла. И като отри потта от лицето си, едвам продума от умора:
\\
— Не мога вече… Надигра ме!
\\
\indentИ един страшен срам щеше да я разплаче.
\\
— Надигра ме, признавам — повтори като на себе си тя.
\\
\indentГайдарят престана да свири.
\\
\indentДъбакът направи последен скок и също запъхтян и уморен, спря. В навалицата избухна весел смях. Стотина гърди си отдъхнаха, стотина премрежени от гледане очи се изясниха. Другарките на Христина почнаха да я подкачат.
\\
\indentДядо Корчан стоеше усмихнат.
\\
\indentДъбакът изтриваше потта от лицето си.
\\
— Е, надигра ме — признавам! — рече засрамено противницата му и скри лицето си в ръце.
\\
\indentНавалицата отново весело зашумя.
\\
— Облога, облога да се изпълни! — завикаха някои и почнаха да подкачат Христина.
\\
\indentТя дигна глава и почти през сълзи отговори:
\\
— Що? Аз стоя на думата си!
\\
\indentШумът утихна. Тия думи бяха толкова смели!
\\
— Да, аз стоя на думата си! — повтори твърдо тя, като се обърна към Лазар Дъбака.
\\
— Истина ли? — поиска повторно да чуе той.
\\
— Истина! — рече твърдо тя.
\\
\indentДъбакът пристъпи при нея и пред всички завърза златния наниз на шията и.
\\
— Тогава да те заведа у дома.
\\
— Заведи ме! — рече тя, като го улови за ръката.
\\
\indentЛазар се усмихна и я поведе надолу към селото.
\\
\indentНавалицата зашумя весело по тях и ги последва.
\\
\indentЕй, Дъбак, ами ветрената мелница? — викна дядо Корчан след тях.
\\
— Аз си намерих друга! — отговори весело Дъбакът.
\\
\indentГлухият старец се засмя така сърдечно, че лулата падна от устата му.

\subsubsection{Двойственият характер на човека (Пейо Яворов)}
\vspace{0.5cm}
\begin{center}
\textbf{ЧАСТ 3 (Време за работа: 120 минути)}
\vspace{0.2cm}
{\textit{\underline{Аргументативния текст напишете в листа за отговори на предвиденото място.}}}
\vspace{0.1cm}
\end{center}
\textbf{41.} Напишете аргументативен текст в обем до 4 страници по ЕДНА от следните две теми:
\\
\vspace{0.2cm}
\textbf{Тема за интерпретативно съчинение:}
\vspace{0.1cm}
\begin{center}
\textbf{"Двойственият характер на човека"}
\\
\textit{(Интерпретативно съчинение върху „Две души“ на Пейо Яворов)}
\end{center}
\begin{tabular}{p{0.48\textwidth} p{0.48\textwidth}}
\begin{minipage}[t]{\linewidth}
Аз не живея: аз горя. Непримирими\\
в гърдите ми се борят две души:\\
душата на ангел и демон. В гърди ми\\
те пламъци дишат и плам ме суши.\\
\\
И пламва двоен пламък, дето се докосна\\
и в каменът аз чуя две сърца…\\
Навсякъде сявга раздвоя несносна\\
и чезнещи в пепел враждебни лица.\\
\end{minipage}
&
\begin{minipage}[t]{\linewidth}
И подир мене с пепел вятъра навсъде\\
следите ми засипва: кой ги знай?\\
Аз сам не живея — горя! — и ще бъде\\
следата ми пепел из тъмен безкрай.\\
\end{minipage}
\end{tabular}

\section{Тестове и матури}
\subsection{Входно ниво}
\noindent
\textit{Име и фамилия:} \dotfill \textit{№} \makebox[2cm]{\dotfill}

\vspace{0.6cm}

\begin{center}
    {\large\bfseries ВХОДНО НИВО}
\end{center}

\vspace{0.6cm}

\noindent
\textbf{Част 1 (време за работа: 60 минути)}

\vspace{0.3cm}

\noindent
\underline{\textit{Отговорите на задачите от 1. до 21. включително отбелязвайте в листа за отговори.}}

\vspace{0.5cm}

\noindent
\textbf{1. В кой ред НЕ е допусната правописна грешка?} \\
\noindent
\textbf{А)} спрени, мятана, стреляли \\
\textbf{Б)} улеснение, осъзнавам, куловози \\
\textbf{В)} таинствен, фотьойли, фоайе \\
\textbf{Г)} издръжлив, задължнял, сълзлив

\vspace{0.6cm}

\noindent
\textbf{2. В кое изречение НЕ е допусната правописна грешка?} \\
\noindent
\textbf{А)} От здравословното му състояние зависи продължителноста на диетата. \\
\textbf{Б)} Италиянският отбор ще пропусне предстоящото световно първенство. \\
\textbf{В)} Обявени са стоките с промоционални цени в големия мебелен магазин. \\
\textbf{Г)} Имам стара травма от падане и при влажно време трудно движа лакъта си.

\vspace{0.6cm}

\noindent
\textbf{3. В кое изречение е допусната правописна грешка?} \\
\noindent
\textbf{А)} Днес жълточерният отбор ще се бори за победа в местното дерби. \\
\textbf{Б)} Професорът ще изнесе още една лекция пред четвъртокурсниците. \\
\textbf{В)} Пациентката направи проверка на своя здравноосигурителен статус. \\
\textbf{Г)} Заместникдиректорският кабинет е разположен на последния етаж.

\vspace{0.6cm}

\noindent
\textbf{4. В кое от изреченията НЕ е допусната граматична грешка?} \\
\noindent
\textbf{А)} Гласуваните поправки в новия закон ще засегнат пряко въздушния и сухопътен транспорт в страната. \\
\textbf{Б)} Убедително написаното мотивационно писмо на младия и амбициозен кандидат впечатли комисията. \\
\textbf{В)} Още утре г-н Каменов, адвоката на братовчедка ми, ще те потърси, за да подготвиш необходимите документи. \\
\textbf{Г)} По Искърско дефиле, което е с дължина около 100 км, са разположени множество интересни обекти.

\vspace{0.6cm}

\noindent
\textbf{5. В коя от подчертаните думи е допусната граматична грешка?} \\
\noindent
За участие в международния фестивал пристигнаха двадесет \underline{музиканта} \textbf{(А)} – шестима млади и награждавани в конкурси \underline{китаристи} \textbf{(Б)} и трима утвърдени \underline{саксофонисти} \textbf{(В)}, разпределени в три \underline{екипа} \textbf{(Г)}.

\vspace{0.6cm}

\noindent
\textbf{6. В кое изречение НЕ е допусната пунктуационна грешка?} \\
\noindent
\textbf{А)} В гардероба в коридора откриха не само празни куфари, но и няколко кашона. \\
\textbf{Б)} Маргарита, неговата голяма и единствена любов неочаквано замина за чужбина. \\
\textbf{В)} Разхождайки се по тесните улици на притихналия град тя се озова на площада. \\
\textbf{Г)} По всяка вероятност, той беше разбрал за предстоящото пътуване на сестра си.

\vspace{0.6cm}

\noindent
\textbf{7. В кое от изреченията е допусната пунктуационна грешка?} \\
\noindent
\textbf{А)} Тя се вцепени от ужас, отвори уста, но не успя да каже нищо. \\
\textbf{Б)} Той понечи да отвори вратата, бравата на която беше счупена. \\
\textbf{В)} Учениците, без да изчакат своя учител се качиха в автобуса. \\
\textbf{Г)} Прибрах се вечерта вкъщи и забелязах, че чантата ми я няма.

\vspace{0.6cm}

\noindent
\textbf{8. В листа за отговори запишете правилната форма на думата, поставена в скоби.} \\[0.3em]
\noindent
По време на пандемията (кандидат-студентската) кампания беше различна.

\vspace{0.6cm}

\noindent
\textbf{9. В листа за отговори запишете правилно САМО думата, в която е допусната правописна грешка.} \\[0.3em]
\noindent
Момчето бързо порастна и стана на 18 години.

\vspace{0.6cm}

\noindent
\textbf{10. В листа за отговори запишете правилната граматична форма на думата, поставена в скоби.} \\[0.3em]
\noindent
Срещнах човека, (който) всички познаваха, и го поздравих.

\vspace{0.6cm}

\noindent
\textbf{11. В листа за отговори запишете правилната форма САМО на думата, в която е допусната граматична грешка.} \\[0.3em]
\noindent
Този човек е моя най-добър приятел.

\vspace{0.6cm}

\noindent
\textbf{12. В листа за отговори запишете правилната форма САМО на думата, в която е допусната граматична грешка.} \\[0.3em]
\noindent
Господин Иванов, Вие сте помислил за всичко.

\vspace{0.6cm}

\noindent
\textbf{13. В текста са пропуснати САМО ПЕТ препинателни знака. Препишете ЦЕЛИЯ текст в листа за отговори, като поставите пропуснатите знаци.} \\[0.3em]
\noindent
На рождения си ден, който бе отбелязан в залата на театъра актьорът призна пред публиката Скъпи приятели цял живот съм на сцената. Съвсем различно обаче е усещането да си сред публиката. Развълнувани от думите му неговите колеги започнаха да го аплодират.

\vspace{0.6cm}

\noindent
\textbf{Запознайте се с текста и диаграмата и изпълнете задачите към тях (от 14. до 21. включително).}

\vspace{0.6cm}

\noindent
\textbf{Текст 1}

\vspace{0.3cm}

\indent Добър 4,05 е средната оценка на задължителния държавен зрелостен изпит по български език и литература (БЕЛ). Зрелостниците са получили средно 52,77 точки от общо 100. Това показват публикуваните днес резултати от задължителните държавни зрелостни изпити по български език и литература и по профилиращи предмети на майската сесия тази
година.
\newline
\indent Усредненият результат на втората задължителна матура по профилиращ предмет е 62 т. и отговаря на оценка добър 4,48. По-високият успех на втория изпит е очакван, като се има предвид, че на първата матура по БЕЛ се явяват всички завършили ХII клас (в случая 43 011), а втората е по избор по профилиращ предмет и на нея отиват по-малко на брой и по- мотивирани зрелостници. На втория изпит по предмет от профилираната подготовка са се явили 24 322 младежи. Останалите са положили държавен изпит по професия.
\newline
\indent С най-висок среден брой точки са чуждите езици. Сред високите е и средният резултат на втората задължителна матура по математика 69.19 точки. В другия край се нареждат природните науки, информационни технологии, география и икономика, предприемачество.
Сравнително слаб е и средният резултат по изобразително изкуство, което се дължи по-скоро на нисък брой точки на затворените отговори, отколкото на практическата задача.
\newline
\indent Няма зрелостник с 0 точки и на двата изпита, но няма и такъв със 100 точки на двата изпита. Максимален брой точки -- 100, на някои от изпитите са постигнали 32 зрелостници. С две пълни шестици са 49 зрелостници, а с две отлични оценки - 1651 зрелостници. Две слаби оценки имат 941 зрелостници. Само девет зрелостници имат 0 точки.

\vspace{0.6cm}

\noindent
\textbf{Текст 2 (диаграма)}

\begin{figure}[htbp]
    \centering
    \includegraphics[width=1\textwidth]{diagram.png} 
\end{figure}

\vspace{0.6cm}

\noindent
\textbf{14. Защо според Текст 1 успехът на втория ДЗИ е по-висок? Посочете НЕВЕРНИЯ отговор.} \\
\noindent
\textbf{А)} Изпитът по БЕЛ е задължителен за всички зрелостници. \\
\textbf{Б)} На втория зрелостен изпит се явяват по-малко ученици. \\
\textbf{В)} Учениците, явяващи се на втория ДЗИ, са по-мотивирани. \\
\textbf{Г)} ДЗИ по БЕЛ е по-труден от останалите изпити.

\vspace{0.6cm}

\noindent
\textbf{15. Посочете НЕВЕРНИЯ отговор. Според Текст 1:} \\
\noindent
\textbf{А)} Няма зрелостник със 100 точки на двата изпита. \\
\textbf{Б)} Няма зрелостник с 0 т. на двата изпита. \\
\textbf{В)} Максимален брой точки на двата изпита имат 32 зрелостници. \\
\textbf{Г)} Две шестици имат 49 зрелостници.

\vspace{0.6cm}

\noindent
\textbf{16. Посочете НЕВЕРНИЯ отговор. Според диаграмата:} \\
\noindent
\textbf{А)} На ДЗИ по БЕЛ 33.84 \% са получили оценка Добър (40.75 -- 62.25 т.). \\
\textbf{Б)} Отличните оценки на втория държавен изпит са с 11.22 \% повече от отличните оценки на изпита по БЕЛ. \\
\textbf{В)} Преобладаващата част от зрелостниците имат оценки Среден и Добър на двата изпита. \\
\textbf{Г)} Оценка Отличен 6.00 (95 -- 100 т.) на ДЗИ по БЕЛ са получили 0.38 \% от зрелостниците.

\vspace{0.6cm}

\noindent
\textbf{17. Посочете НЕВЕРНИЯ отговор. Според Текст 1 и Текст 2:} \\
\noindent
\textbf{А)} Две слаби оценки (до 29.5 т.) имат 941 зрелостници. \\
\textbf{Б)} Две отлични оценки (84 -- 94.5 т.) имат 1651 зрелостници. \\
\textbf{В)} От всички зрелостници, явили се на ДЗИ по БЕЛ (43 011), оценка 6.00 (95 -- 100 т.) имат само 0.38 \% \\
\textbf{Г)} Само 1651 зрелостници имат две отлични оценки (84 -- 94.5 т.).

\vspace{0.6cm}

\noindent
\textbf{18. Според Текст 1 средният резултат на зрелостниците на втория задължителен изпит по математика е 69.19 т. Запишете каква оценка имат зрелостниците според диаграмата?}

\vspace{0.6cm}

\noindent
\textbf{19. Като имате предвид информацията от двата текста, запишете колко точки НЕ достигат на зрелостниците, за да имат оценка Мн. добър на втория зрелостен изпит.}

\vspace{0.6cm}

\noindent
\textbf{20. Запишете СЪС СВОИ ДУМИ защо зрелостниците имат по-нисък среден успех на ДЗИ по БЕЛ?}

\vspace{0.6cm}

\noindent
\textbf{21. Като имате предвид информацията от Текст 1, коментирайте в текст до 5 изречения следното:}

\begin{itemize}
    \item[\textendash] \textbf{Защо зрелостниците имат по-висок успех на втория ДЗИ?}
    \item[\textendash] \textbf{Как си обяснявате факта, че средно най-много точки зрелостниците имат на ДЗИ по чужд език?}
    \item[\textendash] \textbf{Защо зрелостниците имат висок среден успех на ДЗИ по математика?}
\end{itemize}

\newpage

\vspace{0.6cm}

\noindent
\textbf{ЧАСТ 2 (време за работа: 60 минути)}

\vspace{0.3cm}

\noindent
\underline{\textit{Отговорите на задачите от 22. до 40. включително отбелязвайте в листа за отговори.}}

\vspace{0.6cm}

\noindent
\textbf{22. В листа за отговори запишете само паронима, с който да поправите лексикалната грешка в изречението.} \\[0.3em]
\noindent
Обращението към публиката бе написано прочувствено.

\vspace{0.6cm}

\noindent
\textbf{23. За всяко празно място изберете УМЕСТНАТА дума и я запишете срещу съответната буква в листа за отговори.} \\[0.5em]
\noindent
Елин Пелин променя из основи жанра на разказа в българската литература. В разказите на Елин Пелин огромна роля за изграждането на художествените внушения играе \dots \textbf{(А)}. Той не е просто \dots \textbf{(Б)}, той е \dots \textbf{(В)} на събитието.

\vspace{0.4cm}

\noindent
\textbf{А)} пейзажът, случката, диалогът \\
\textbf{Б)} среда, фон, заден план \\
\textbf{В)} основата, сърцевината, душата

\vspace{0.6cm}

\noindent
\textbf{24. Какви са измерения на родното и чуждото за героите? Посочете НЕВЯРНОТО твърдение.} \\
\noindent
\textbf{А)} За Султана чуждо е всичко, което може да й попречи да укрепи и да въздигне западналия си род. \\
\textbf{Б)} За Лазар родното се свързва не само със семейно-родовата общност, то включва и понятията \textit{народ} и \textit{родина}. \\
\textbf{В)} За Рафе Клинче своето е творчеството, красотата и любовта. \\
\textbf{Г)} Родният свят за Стоян Глаушев е представен в романтична светлина.

\vspace{0.6cm}

\noindent
\textbf{25. Кой от проблемите е интерпретиран и в двата посочени откъса?} \\[0.5em]
\noindent
\textit{От днеска нататък българският род} \\
\textit{история има и става народ! (Ив. Вазов, "Паисий")} \\[0.5em]
\textit{Но разкажи със думи прости} \\
\textit{на тях -- на бъдещите хора,} \\
\textit{които ще поемат поста ни} \\
\textit{че ние храбро сме се борили. (Н. Вапцаров, "История")} \\[0.6cm]
\noindent
\textbf{А)} за написаната история като средство за самопознанието на нацията \\
\textbf{Б)} за историята като средство за духовното извисяване на българите \\
\textbf{В)} за написаната история като форма на колективната памет \\
\textbf{Г)} за гордостта на човека от историята

\vspace{0.6cm}

\noindent
\textbf{26. Кой от мотивите НЕ е интерпретиран в \textit{"Новото гробище над Сливница"}?} \\[0.3em]
\noindent
\textbf{А)} за героичната смърт, която осмисля човешкото съществуване \\
\textbf{Б)} за песента като прослава, в която се сливат човешкото и природното \\
\textbf{В)} за забравата и паметта \\
\textbf{Г)} за преходността на човешкото битие

\vspace{0.6cm}

\noindent
\textbf{27. Кое твърдение НЕ е вярно?} \\
\noindent
\textbf{А)} Анафоричното повторение на израза \textit{Сега съм у дома} в началото на всяка от осемте строфи на Вазовата ода \textit{"При Рилския манастир"} разкрива величието на природата. \\
\textbf{Б)} Ботев изгражда в стихотворението \textit{"Борба"} образа на един свят на преобърнатите ценности, като специална роля в това отношение играят оксимороните (\textit{лъжи святи, свещена глупост\dots}). \\
\textbf{В)} Двойната употреба на епитета \textit{бурлив} в стихотворението на Вазов \textit{"Новото гробище над Сливница"} внушава идеята, че човекът, извисен в битката за родината, се приближава до величавостта на природата. \\
\textbf{Г)} Инверсираното тавтологично словосъчетание \textit{ниви глухо опустели} показва трагичната самота на човека в Яворовата поема \textit{"Градушка"}.

\vspace{0.6cm}

\noindent
\textbf{28. Посочете значенията на паратекста (заглавието и епиграфа).} \\[0.5em]
\begin{flushright}
    \textit{"Приказка за стълбата"} \\
    \textit{Посветено на всички, които ще кажат:} \\
    \textit{Това не се отнася до мене!}
\end{flushright}
\vspace{0.4cm}
\noindent
\textbf{А)} Думата \textit{приказка} в заглавието насочва към представата за иносказание, фантазия. \\
\textbf{Б)} Думата \textit{стълба} в заглавието подсказва идеята за изкачване, за слизане. \\
\textbf{В)} Епиграфът (мотото) е предупреждение за всички, изкачващи се по стълбата на властта с цената на непростими нравствени компромиси, но отказващи да се разпознаят в героя на текста. \\
\textbf{Г)} В паратекста е зададен проблемът за обществото и властта, който ще се интерпретира в текста.

\vspace{0.6cm}

\noindent
\textbf{29. Кое от твърденията за пейзажа в Елин-Пелиновия разказ \textit{"Андрешко"} НЕ е вярно?} \\[0.5em]
\noindent
\textit{Земята тънеше в кал и влага. Мъртви и покрусени се тъмнееха пръснатите пейзажи на села, речища, далечни гори и планини. По полето лъщяха големи локви, мътни, студени и оцъклени като мъртвешки очи.} \\[0.6cm]
\noindent
\textbf{А)} Пейзажът е мрачен, потискащ, изпълнен с вариации на мъртвото и безизходното. \\
\textbf{Б)} Значенията на водата, която символизира обновлението, са преобърнати. \\
\textbf{В)} Мотивът за калта също носи внушения за безнадеждност и предопределеност. \\
\textbf{Г)} Мъртвата вода и неизбродимата кал подсказват, че героите са затънали в безчовечността.

\vspace{0.6cm}

\noindent
\textbf{30. Кое от твърденията НЕ е вярно?} \\
\noindent
\textbf{А)} Природата в одата \textit{"При Рилския манастир"} е възприета като храм на човешката душа, докато в поемата \textit{"Градушка"} човекът е безпомощен и обречен пред немилостивото лице на природата. \\
\textbf{Б)} Човекът е самотен в съзиданието и страданието, но в миниатюрата \textit{"Спи езерото"} (за разлика от поемата \textit{"Градушка"}) самотата е съзнателен избор, начин да се достигне самовглъбяване. \\
\textbf{В)} Природният кръговрат в одата \textit{"При Рилския манастир"} е символ на вечността, а в миниатюрата \textit{"Спи езерото"} -- на безвремието. \\
\textbf{Г)} Съзерцанието е път към духовната същност на живота и в одата \textit{"При Рилския манастир"}, и в поемата \textit{"Градушка"}.

\vspace{0.6cm}

\noindent
\textbf{31. Сходни моменти в изображението на родното, на балканските нрави има в:} \\
\noindent
\textbf{А)} \textit{"Железният светилник"} и \textit{"Бай Ганьо журналист"} \\
\textbf{Б)} \textit{"Бай Ганьо журналист"} и \textit{"Балкански синдром"} \\
\textbf{В)} \textit{"Балкански синдром"} и \textit{"При Рилския манастир"} \\
\textbf{Г)} \textit{"При Рилския манастир"} и \textit{"Новото гробище над Сливница"}

\vspace{0.6cm}

\noindent
\textbf{32. Кой герой от \textit{"Крадецът на праскови"} претърпява най-сериозно вътрешно развитие?} \\[0.4em]
\noindent
\textbf{А)} полковникът \hspace{1.5cm} \textbf{Б)} Елисавета \hspace{1.5cm} \textbf{В)} младият разказвач \hspace{1.5cm} \textbf{Г)} Иво

\vspace{0.6cm}

\noindent
\textbf{33. Посочете НЕВЯРНОТО твърдение} \\
\noindent
\textbf{А)} Историята на любовта между Елисавета и Иво е видяна през две различни гледни точки. \\
\textbf{Б)} Първият разказвач въвежда читателя в общата атмосфера на действието (наблюдението и спомените от живота на града по време на двете световни войни. \\
\textbf{В)} Вторият разказвач е учителят. \\
\textbf{Г)} Всевиждащият разказвач е необходим, за да може авторът да разграничи свидетелската от всевиждащата гледна точка.

\vspace{0.6cm}

\noindent
\textbf{34. Коя тема НЕ е интерпретирана в откъса от стихотворението на Иван Вазов \textit{"Към природата"}?} \\[0.5em]
\noindent
\textit{Поклон на теб, природо, създанье необятно} \\
\textit{на твоя свод лазурен, на твойто слънце златно,} \\
\textit{на твойта вечна младост и вечна красота,} \\
\textit{на всичко, що е в тебе божествено и тайно,} \\
\textit{невидимо, кат бога, велико и безкрайно} \\
\textit{и равно с вечността;} \\[0.6cm]
\noindent
\textbf{А)} за природата като дом \\
\textbf{Б)} за природната хармония \\
\textbf{В)} за природата като носител на божествения промисъл \\
\textbf{Г)} за природата като коректив на обществото

\vspace{0.6cm}

\noindent
\textbf{35. Съпоставете интерпретациите на темата за природата в \textit{"При Рилския манастир"} (\textit{"Към природата"}) и в \textit{"Градушка"} и запишете ЕДНА разлика в тях в листа за отговори.}

\vspace{0.6cm}

\noindent
\textbf{36. В листа за отговори разтълкувайте заглавието \textit{"Железният светилник"} в контекста на творбата.}

\vspace{0.6cm}

\noindent
\textbf{37. В листа за отговори посочете двете значения на думата \textit{история} в стихотворението \textit{"История"}.}

\vspace{0.6cm}

\noindent
\textbf{38. В листа за отговори запишете ЕДНА характерна черта за жанра на одата, проявена в \textit{"Паисий"}.}

\vspace{0.6cm}

\noindent
\textbf{39. Свържете заглавието на всяка от творбите с нейния автор. В листа за отговори срещу съответната буква запишете цифрата, съответстваща на името на автора.} \\[0.8em]
\begin{minipage}[t]{0.45\textwidth}
    \textbf{А)} \textit{"Ноев ковчег"} \\
    \textbf{Б)} \textit{"Балкански синдром"} \\
    \textbf{В)} \textit{"Спи езерото"}
\end{minipage}
\begin{minipage}[t]{0.45\textwidth}
    1. Пенчо Славейков \\
    2. Йордан Радичков \\
    3. Станислав Стратиев
\end{minipage}

\vspace{0.6cm}

\vspace{0.8cm}

\noindent
\textbf{40. В текст до 5 изречения съпоставете интерпретациите на темата за невъзможния диалог между народ и държава в посочените откъси.} \\[0.5em]

\newline
\indent Земеделецът, който едно време кажи-речи дивееше щастливо и в хамбара му се намираше от ново до ново, сега дояжда до Коледа и как се чеше по врата край кофтора в механата, говори: „Тежко се живее!“ И наистина е тежко: кръчмарят иска, лихварят иска, попът иска, касата иска, бирникът иска, за къща трябва\dots{} А пък тая пуста пара от нийде не иде. Нито работа има, ни печал има. Ако се опита да открадне една кола дърва от гората, за да вземе някой и друг грош за тях, ето че горският му съставя акт, дето се казва, че тоя човек нарушил член еди-кой си от закона за горите, как отсякъл в еди-коя си държавна гора толко и толко дървета с такава и такава дебелина, и затова му се пада такава и такава глоба\dots{} Па защо и какво е сякъл, това никой не пита. (\textit{Елин Пелин, "За едни майка, за други мащеха"}) \\[0.4em]
\newline
\indent По изпълнение, разбира се. Играе ме един вашенец там, ама тоя път ще му дам да разбере. Няколко пъти го гоня, а все се изплъзва от ръцете ми\dots{} Разбрах, че дяволува, и ще го издебна довечера, та ще ме помни\dots{} Житцето му ще секвестирам! Хем него на ум ще науча, хем да ви покажа един пример, та да не кръшкате друг път. Лъжете търговците, лъжете гражданите, та им продавате развалени яйца и граниво масло. Е, ама чакай, жено селяшка -- властта се лесно не лъже! Пипа тя -- здраво пипа! За вас камшик, руски камшик -- така ще се оправите\dots{} Станахте пияници, пропаднахте, развалихте се -- ще станете неспособни данъкоплатци и ще съсипете държавата. Ах, защо нямам повечечко власт! Ангели ще ви направя. (\textit{Елин Пелин, "Андрешко"})

\newpage

\noindent
\textbf{ЧАСТ 3 (Време за работа: 20 минути)}

\vspace{0.3cm}

\noindent
\underline{\textit{Не пречертавайте цялата таблица, ако не принтирате варианта. Номерирайте редовете}}

\vspace{0.6cm}

\noindent
\textbf{41. Попълнете пропуснатото в таблицата.}

\vspace{0.3cm}

\noindent
\renewcommand{\arraystretch}{1.8}
\begin{tabularx}{\textwidth}{|>{\raggedright\arraybackslash}X|>{\raggedright\arraybackslash}X|>{\raggedright\arraybackslash}X|>{\raggedright\arraybackslash}X|}
\hline
\hfil\bfseries Автор & \hfil\bfseries Произведение & \hfil\bfseries Жанр & \hfil\bfseries Герой \\ \hline
 &  & роман & Аврам Немтур \\ \hline
Емилиян Станев &  &  & ординарецът \\ \hline
 & „Андрешко“ &  & съдия-изпълнителят \\ \hline
 & „Бай Ганьо журналист“ &  & Гуньо Адвокатина \\ \hline
Христо Смирненски &  & разказ, притча &  \\ \hline
\end{tabularx}

\vspace{0.3cm}

\noindent
\renewcommand{\arraystretch}{1.8}
\begin{tabularx}{\textwidth}{|>{\raggedright\arraybackslash}X|>{\raggedright\arraybackslash}X|>{\raggedright\arraybackslash}X|>{\raggedright\arraybackslash}X|}
\hline
\hfil\bfseries Автор & \hfil\bfseries Тема & \hfil\bfseries Произведение & \hfil\bfseries Жанр \\ \hline
 & Природата &  & лирическа миниатюра \\ \hline
Иван Вазов &  & „Паисий“ &  \\ \hline
 &  & „Балкански & комедия на абсурда \\ 
 &  & синдром“ &  \\ \hline
Христо Ботев & Обществото и &  &  \\ 
 & властта &  &  \\ \hline
Емилиян &  &  & повест \\ 
Станев &  &  &  \\ \hline
\end{tabularx}

\newpage
\subsection{Тест върху "Железният светилник" и правописна норма}
\noindent
\textit{Име и фамилия:} \dotfill \textit{№} \makebox[2cm]{\dotfill}

\vspace{0.6cm}

\begin{center}
    {\large\bfseries Тест върху "Железният светилник" и правописна норма}
\end{center}

\vspace{0.6cm}

\noindent
\textbf{Част 1 (време за работа: 60 минути)}

\vspace{0.3cm}

\noindent
\underline{\textit{Отговорите на задачите от 1. до 21. включително отбелязвайте в листа за отговори.}}

\vspace{0.5cm}

\noindent
\textbf{1. В коя от думите има правописна грешка?} \\
\noindent
\textbf{А)} инженер \\
\textbf{Б)} акомулатор \\
\textbf{В)} тирбушон \\
\textbf{Г)} пеницилин

\vspace{0.6cm}

\noindent
\textbf{2. В кой ред НЕ е допусната правописна грешка?} \\
\noindent
\textbf{А)} манисто, астма, бастун \\
\textbf{Б)} бодър, разпра, огризение \\
\textbf{В)} влакче, прекупувам, заводски \\
\textbf{Г)} мързелив, неосетно, овощни

\vspace{0.6cm}

\noindent
\textbf{3. В кое от изреченията е допусната правописна грешка?} \\
\noindent
\textbf{А)} По образование и квалификация фелшерът е между лекаря и медицинската сестра. \\
\textbf{Б)} Агенцията по вписванията издаде методически указания за работата на длъжностните лица по регистрацията. \\
\textbf{В)} Уестминстърското абатство е посещавано всяка година от над милион туристи. \\
\textbf{Г)} Повечето известни писатели са били хора със сложен характер, с богата душевност и интересна съдба.

\vspace{0.6cm}

\noindent
\textbf{4. В кое от изреченията е допусната правописна грешка?} \\
\noindent
\textbf{А)} Титулярът на сметките, по които са получавани преводи от чужбина, е български гражданин. \\
\textbf{Б)} Букварът е част от цялостен учебен комплект по български език и литература за 1. клас по новата учебна програма. \\
\textbf{В)} Формулярът за кандидатстване за неимигрантска виза трябва да бъде попълнен изцяло. \\
\textbf{Г)} Съединителят представлява връзката между скоростната кутия и двигателя, преди да е натиснат педалът.

\vspace{0.6cm}

\noindent
\textbf{5. В кое от изреченията НЕ е допусната правописна грешка?} \\
\noindent
\textbf{А)} Статията разглежда спецификите и параметрите на употреба на понятията „имидж“, „репутация“ и „бранд“. \\
\textbf{Б)} Лудата гонидба на върха продължи стремглаво и в 26-ия кръг на английското футболно първенство. \\
\textbf{В)} Рупите е защитена местност, която е разположена на около 10 км от Петрич и на 2 км от с. Рупите. \\
\textbf{Г)}От около две години стоките на куфърните търговии, които редовно пътуват между България и Турция, се обмитяват.

\vspace{0.6cm}

\noindent
\textbf{5. В кое от изреченията НЕ е допусната правописна грешка?} \\
\noindent
\textbf{А)} Статията разглежда спецификите и параметрите на употреба на понятията „имидж“, „репутация“ и „бранд“. \\
\textbf{Б)} Лудата гонидба на върха продължи стремглаво и в 26-ия кръг на английското футболно първенство. \\
\textbf{В)} Рупите е защитена местност, която е разположена на около 10 км от Петрич и на 2 км от с. Рупите. \\
\textbf{Г)}От около две години стоките на куфърните търговии, които редовно пътуват между България и Турция, се обмитяват.

\vspace{0.6cm}

\noindent
\textbf{6. В коя от подчертаните думи е допусната правописна грешка?} \\
\noindent
\underline{\textbf{Информираността}} \textbf{(А)} относно Европа и европейската парламентарна демокрация, която осигурява Програмата за училища \underline{\textbf{посланици}} \textbf{(Б)} на Европейския парламент, цели да превърне младите хора в активни европейски граждани и да ги \underline{\textbf{насарчи}} \textbf{(В)} да участват в различни \underline{\textbf{инициативи}} \textbf{(Г)} на местно и международно ниво.

\vspace{0.6cm}

\noindent
\textbf{7. Открийте и поправете четиринайсетте правописни грешки в текста.} \\
\noindent
През 1879 година в северната часть на Испания любителят архиолог маркис Де Саутуола, придружен от четиригодишната си дъщеря, решил да разгледа пещерата Алтамира. Малката палълница явно не се интересувала от фкаменелостите, вълнуващи бащата. Но в пещерата имало няколко ниски галерии, които никой оважаващ себе си архиолог не би изследвал и където можело приятно да се поиграе. Детето пропалзяло в долната част на пещерата, носейки свещицата си. Но когато пугледнало към тавана, за най-голям негов ужас, одгоре се изблещили очите на угромен нарисуван бик. Оказало се, че стените и таваните на тези ниски галерии били покрити с прекрасни рисунки, дело на парвобитния човек.

\vspace{0.6cm}

\noindent
\textbf{8. В коя от думите е допусната правописна грешка?} \\
\noindent
\textbf{А)} колониален \\
\textbf{Б)} остта \\
\textbf{В)} влекат \\
\textbf{Г)} целостта

\vspace{0.6cm}

\noindent
\textbf{9. В коя от думите НЕ е допусната правописна грешка?} \\
\noindent
\textbf{А)} гръкът \\
\textbf{Б)} глобалистски \\
\textbf{В)} радиятор \\
\textbf{Г)} бактерия

\vspace{0.6cm}

\noindent
\textbf{10. В кое от изреченията е допусната правописна грешка?} \\
\noindent
\textbf{А)} Котките обичат чистотата и непрестанно се ближат, за да поддържат телата си чисти, както ги е научила еволюциата. \\
\textbf{Б)} Смята се, че малката гънка в края на котешкото ухо, образуваща продълговата вдлъбнатинка, изпълнява функцията на резонатор. \\
\textbf{В)} Котаракът Мирон пътешествал с геолози из платото Устюрт и самоотвержено охранявал лагера на експедицията от отровни змии. \\
\textbf{Г)} Днес котката е сред малкото бозайници, които могат да живеят с нас и в тесния, малък апартамент.

\vspace{0.6cm}

\noindent
\textbf{11. В кой ред е допусната правописна грешка?} \\
\noindent
\textbf{А)} министър на туризма \\
\textbf{Б)} цар Борис Трети \\
\textbf{В)} Столична Община \\
\textbf{Г)} хардрок

\vspace{0.6cm}

\noindent
\textbf{12. В кое от изреченията е допусната правописна грешка?} \\
\noindent
\textbf{А)} От древните летописци научаваме, че Александър Велики пренесъл за пръв път папагалите в Европа. \\
\textbf{Б)} На галапагоските острови живеят над двайсет животински вида, намиращи се под закрила, сред които е и морската игуана. \\
\textbf{В)} Котката в Древен Египет била не само свещено животно, но и ловец на гризачи, а също така и воин. \\
\textbf{Г)} От Египет домашната котка попаднала в Индия, Югоизточна Азия и Южна Европа.

\vspace{0.6cm}

\noindent
\textbf{13. В кое от изреченията НЕ е допусната правописна грешка?} \\
\noindent
\textbf{А)} Плиний Стари твърди, че страните на изтока са първите, в които се употребявали ароматни вещества. \\
\textbf{Б)} В древния Рим козмети се наричали робините, които се грижели за красотата на римлянките. \\
\textbf{В)} В средните векове само по време на църковна служба свещеници кадели тамян и помазвали вярващите с миро. \\
\textbf{Г)} През 1608 година монаси от доминиканския орден основали в манастир близо до Флоренция първата „парфюмерийна фабрика“.

\vspace{0.6cm}

\noindent
\textbf{14. В кое от изреченията е допусната правописна грешка?} \\
\noindent
\textbf{А)} Излагането на слънце през лятото не е опасно, ако спазваме съветите на лекарите. \\
\textbf{Б)} В Средна и Южна Европа август е един от най-хубавите месеци и плажовете са пълни с летовници. \\
\textbf{В)} Понякога ледът на повърхността на Северния Ледовит океан достига през зимата до три метра дебелина. \\
\textbf{Г)} Със застудяването на времето се укриват и редица други животни, които прекарват зимата в летаргичен сън.

\vspace{0.6cm}

\noindent
\textbf{15. В коя от подчертаните думи е допусната правописна грешка?} \\
\noindent
Декораторите обръщали особено внимание на камбаната и гребена на шлема. Заедно с изкованите сцени от \underline{\textbf{Античността}} \textbf{(А)}, \underline{\textbf{инкрустираните}} \textbf{(Б)} със злато или сребро \underline{\textbf{Ренесансови}} \textbf{(В)} мотиви гравирането било един от предпочитаните начини за \underline{\textbf{украсяване}} \textbf{(Г)}.

\vspace{0.6cm}

\noindent
\textbf{16. В коя от подчертаните думи НЕ е допусната правописна грешка?} \\
\noindent
Нова Зеландия, която се намира в югозападната част на Тихия \underline{\textbf{океан}} \textbf{(А)}, е островна държава, като по-голямата част от \underline{\textbf{територията}} \textbf{(Б)} й \underline{\textbf{опхваща}} \textbf{(В)} двата острова, наречени Северен и Южен, разделени от \underline{\textbf{Протока}} \textbf{(Г)} Кук.

\vspace{0.7cm}

\noindent
\textbf{18. В коя от думите има правописна грешка?} \\
\noindent
\textbf{А)} автокъща \\
\textbf{Б)} акваспининг \\
\textbf{В)} жълто-кафяв \\
\textbf{Г)} прахосмукачка-робот

\vspace{0.6cm}

\noindent
\textbf{18. В кой ред е допусната правописна грешка?} \\
\noindent
\textbf{А)} аутокю, астроном любител, закононарушение \\
\textbf{Б)} Дон-Кихотов, кънтри-певец, ландшафт \\
\textbf{В)} културтрегер, културно-просветен, родословие \\
\textbf{Г)} токшоу, помощник-готвач, кафе пауза

\vspace{0.6cm}

\noindent
\textbf{19. В кой ред НЕ е допусната правописна грешка?} \\
\noindent
\textbf{А)} кафе-сладкарница, ПИН-код, екопътека \\
\textbf{Б)} младши лейтенант, лироепически, килокалории \\
\textbf{В)} криотерапия, хипер инфлация, геолого-географски \\
\textbf{Г)} тупа-лупа, високоблагородие, жената котка

\vspace{0.6cm}

\noindent
\textbf{20. В коя от подчертаните думи е допусната правописна грешка?} \\
\noindent
Насекомите са били сред първите животни, които „са излезли от морето, отърсили са се от калта и са се добили с криле“, казва американският \underline{\textbf{ентомолог}} \textbf{(А)} Кати Прудич. Те произвеждат смайващо на брой поколение -- например мравките от вида Dorylus gribodoi снасят \underline{\textbf{три-четири}} \textbf{(Б)} милиона яйца на всеки 28 дни. \\
Малките същества разполагат и със завидни защитни оръжия: от \underline{\textbf{супер-твърдата}} \textbf{(В)} обвивка на някои бръмбари до токсичните шипове на южния фланелен молец. \\
Насекомите правят огромна услуга на \underline{\textbf{земеделието}} \textbf{(Г)} и околната среда, като обработват и обогатяват почвата, разпръскват семена и опрашват количество растения, представляващо около една трета от храната, произвеждана в света.

\vspace{0.6cm}

\noindent
\textbf{21. В кое от изреченията е допусната правописна грешка?} \\
\noindent
\textbf{А)} Таджикистан започва изграждането на водноелектрическата централа „Рогун“ -- проект на стойност 3,9 милиарда долара. \\
\textbf{Б)} Какво трябва непременно да знаете, преди да започнете тренировки по кикбокс? \\
\textbf{В)} Наскоро бившият министър и негови съмишленици учредиха център за евро-атлантическа сигурност. \\
\textbf{Г)} Има ли разлика между сухите и полу-сухите вина освен по отношение на захарта, която съдържат?

\vspace{0.6cm}

\noindent
\textbf{22. В кое от изреченията НЕ е допусната правописна грешка?} \\
\noindent
\textbf{А)} Корупционен скандал, в който са замесени и чужди служби, тресе северно-македонската съдебна система. \\
\textbf{Б)} За производството на шперплат най-често се използват топола, бреза, липа и иглолистна дървесина. \\
\textbf{В)} Николай Бердяев е руски философ и публицист, един от основателите на т.нар. Руски философско-религиозен ренесанс. \\
\textbf{Г)} Невиждайки кръжащия орел, кокошката спокойно разхождаше пиленцата си.

\vspace{0.6cm}

\noindent
\textbf{23. В кое от изреченията е допусната правописна грешка?} \\
\noindent
\textbf{А)} Възложителят поставя изисквания за защита на информацията с конфиденциален характер. \\
\textbf{Б)} Изберете демонстрационен автомобил, регистрирайте се на сайта ни и си запишете час за тест драйв. \\
\textbf{В)} Доброто кръвообръщение е от съществено значение за здравето на сърцето и на целия организъм. \\
\textbf{Г)} За Романтизма и Модернизма ключова е идеята за разкриване на скритото, за манифестиране на невидимото.

\newpage

\vspace{0.6cm}

\noindent
\textbf{ЧАСТ 2 (време за работа: 60 минути)}

\vspace{0.3cm}

\noindent
\textbf{24. Как НЕ може да се определят романите на Димитър Талев?} \\
\noindent
\textbf{А)} като социално-психологически \\
\textbf{Б)} като представителни за възрожденската традиция \\
\textbf{В)} като принадлежащи към реалистичната традиция \\
\textbf{Г)} като исторически

\vspace{0.6cm}

\noindent
\textbf{25. Коя от особеностите НЕ е характерна за четирилогията на Димитър Талев?} \\
\noindent
\textbf{А)} Основната й тема е историческата съдба на Македония. \\
\textbf{Б)} Вписва се в „голямата лирическа вълна“ от 50-те години на 20. век. \\
\textbf{В)} Романите са публикувани от 1952 до 1966 г. \\
\textbf{Г)} Те се появяват в условията на политическа диктатура.

\vspace{0.6cm}

\noindent
\textbf{26. Коя от опозициите е водеща за идеята на „Железният светилник“?} \\
\noindent
\textbf{А)} мрак -- светлина \\
\textbf{Б)} робство -- свобода \\
\textbf{В)} село -- град \\
\textbf{Г)} материално -- духовно

\vspace{0.6cm}

\noindent
\textbf{27. Кое е вярното твърдение?} \\
\noindent
\textbf{А)} „Железният светилник“ е първата част от трилогията на Димитър Талев за българите в Македония. \\
\textbf{Б)} „Железният светилник“ е единствената творба на Талев с тема и герои от епохата на Възраждането. \\
\textbf{В)} Вазовият роман „Под игото“ е оказал силно емоционално влияние върху Димитър Талев. \\
\textbf{Г)} Действието в „Железният светилник“ се развива в град Прилеп по време на османското владичество.

\vspace{0.6cm}

\noindent
\textbf{28. Коя от опозициите НЕ се открива в романа „Железният светилник“?} \\
\noindent
\textbf{А)} родното -- чуждото \\
\textbf{Б)} селото -- градът \\
\textbf{В)} образованието -- трудът \\
\textbf{Г)} патриархалната норма -- индивидуалният избор

\vspace{0.6cm}

\noindent
\textbf{29. Кои са трите основни ценности, на които се уповава събуденото национално самосъзнание на българина в Талевия роман? Посочете НЕВЕРНИЯ отговор.} \\
\noindent
\textbf{А)} християнската вяра като сплотяваща етноса \\
\textbf{Б)} българският език като основен народностен маркер \\
\textbf{В)} просветата – събуждането на историческа памет и на чувство за родов дълг \\
\textbf{Г)} единението с братята гърци срещу турските поробители

\vspace{0.6cm}

\noindent
\textbf{30. Как Лазар Глаушев разбира мисията на народния водач? Посочете НЕВЯРНОТО твърдение.} \\
\noindent
\textbf{А)} като бягство от личните си терзания \\
\textbf{Б)} като отказ от следване на личния интерес в името на обществените задачи \\
\textbf{В)} като поемане на рискове в името на народната кауза \\
\textbf{Г)} като смелост в противопоставянето срещу гръцките духовници и фанариоти

\vspace{0.6cm}

\noindent
\textbf{31. Защо Султана избира Стоян за свой съпруг?} \\
\noindent
\textbf{А)} защото е влюбена в него \\
\textbf{Б)} защото вижда в него здрав и надежден стълб на своето бъдещо семейство \\
\textbf{В)} защото иска да направи напук на старата баба хаджийка \\
\textbf{Г)} защото смята старите патриархални норми и традиции за отживелица

\vspace{0.6cm}

\noindent
\textbf{32. Избройте четири промени в обществения живот в Преспа, които се осъществяват в хода на сюжета.} \\
\noindent
\makebox[\textwidth]{\dotfill} \\
\makebox[\textwidth]{\dotfill} \\
\makebox[\textwidth]{\dotfill} \\
\makebox[\textwidth]{\dotfill}

\vspace{0.6cm}

\noindent
\textbf{33. Кой от метафоричните образи, използвани от рилския духовник, е израз на стремежа на Султана да съхрани преди всичко семейството?} \\
\noindent
\textbf{А)} черупката, в която се опитваме да се скрием, без да подозираме, че „ще дойде черният гарван и ще я счупи с клюна си“ \\
\textbf{Б)} орелът, който лети към небето със свободни крила, „а нашите крила са сковани“ \\
\textbf{В)} българите като „широк залив от безбрежно славянско море“ \\
\textbf{Г)} овце, които гръцките фанариоти искат да смесят със своето стадо, за да стане по-силно

\vspace{0.6cm}

\noindent
\textbf{34. Каква НЕ е позицията на героя повествовател в „Железният светилник“?} \\
\noindent
\textbf{А)} Представя детайлно собствения си поток на съзнанието. \\
\textbf{Б)} Вижда героите обективно и построява техните характеристики. \\
\textbf{В)} Представя обществените събития и протичането на историята. \\
\textbf{Г)} Тази на класическия третоличен „вездесъщ“ повествовател.

\vspace{0.6cm}

\noindent
\textbf{35. За какво НЕ служи честата употреба на полупряка реч в романа „Железният светилник“?} \\
\noindent
\textbf{А)} Повествователят се доближава максимално до гледните точки на героите. \\
\textbf{Б)} Така се дава обяснение на скритите мотиви на героите. \\
\textbf{В)} Служи в някои случаи за иронично дистанциране от думите на героите. \\
\textbf{Г)} Дава се обяснение на значението на архаичните или диалектните думи.

\vspace{0.8cm}

\textit{Прочетете откъса от романа „Железният светилник“ и изпълнете задачите от 36. до 38. включително.}

\vspace{0.4cm}

\noindent
Малко по-късно Султана и Лазар останаха сами в голямата стая. \\
– Майко, какво става с Катерина? \\
Султана притвори очи, като да се замисли, или за миг се отдаде на умората си. \\
– Какво й е... Пометна. Кръвта не спира. \\
– Но, мамо... – викна Лазар с \underline{\textbf{пресекнал}} \textbf{(А)} глас, – трябва да направим нещо! Не може тъй да чакаме. Да вземем една кола, да я заведем в Битоля. Там има добри хекими. \\
– В Битоля... \underline{\textbf{Ако изтрае днеска...}} \textbf{(Б)} Като се пусне еднаж кръв, мъчно е да се спре. Сама може да спре, а Битоля е далеко. Ако не спре кръвта сама... \\
Лазар каза, без да погледне майка си: \\
– Защо направи това, майко? Щом било толкова опасно, да беше останало така най-сетне. Ето, Катерина може да загуби живота си. \\
– Да – въздъхна Султана, – може... Мислех си, толкова жени пометат. Знаех също, че и много жени затриват живота си. Ако Катерина умре, нека сени ми каже: ти я уби! Направих го, а пък аз й съм майка. Никого нема да боли повече от мене. Аз си знам колко ме боли и как се страхувам за живота й. Но и сега да трябва да почна – пак ще го направя. \\
Тя дигна към сина си остър, мрачен поглед: – Не искам моята щерка да роди копеле. Кой както иска, така да ме съди. \\
Тя се наведе над огнището, посбута огъня без нужда. Лазар мълчеше. Той не одобреваше постъпката на майка си, но \underline{\textbf{кой може да съди една майка за това, което върши за децата си, дори}} \underline{\textbf{и когато прави грешки}} \textbf{(В)}?

\vspace{0.8cm}

\noindent
\textbf{36. Със свои думи посочете три примера за грешките, които допуска Султана.} \\
\noindent
\makebox[\textwidth]{\dotfill} \\
\makebox[\textwidth]{\dotfill} \\
\makebox[\textwidth]{\dotfill} \\
\makebox[\textwidth]{\dotfill}

\vspace{0.6cm}

\noindent
\textbf{37. С цитати от речта на повествователя посочете четири примера, които доказват, че Талев е майстор на разгърнатите психологически характеристики на своите герои.} \\
\noindent
\makebox[\textwidth]{\dotfill} \\
\makebox[\textwidth]{\dotfill} \\
\makebox[\textwidth]{\dotfill} \\
\makebox[\textwidth]{\dotfill} \\
\makebox[\textwidth]{\dotfill} \\
\makebox[\textwidth]{\dotfill}

\vspace{0.6cm}

\noindent
\textbf{38. Срещу съответната буква запишете какви са подчертаните изразни средства в текста.} \\
\noindent
\textbf{А)} \makebox[0.8\textwidth]{\dotfill} \\
\textbf{Б)} \makebox[0.8\textwidth]{\dotfill} \\
\textbf{В)} \makebox[0.8\textwidth]{\dotfill}
\newpage
\subsection{Тест върху раздела Любовта и лексикална норма}
\noindent
\textit{Име и фамилия:} \dotfill \textit{№} \makebox[2cm]{\dotfill}

\vspace{0.6cm}

\begin{center}
    {\large\bfseries Тест върху раздела Любовта и лексикална норма}
\end{center}

\vspace{0.6cm}

\noindent
\textbf{Част 1 (време за работа: 60 минути)}

\vspace{0.3cm}

\noindent
\underline{\textit{Отговорите на задачите от 1. до 21. включително отбелязвайте в листа за отговори.}}

\vspace{0.5cm}

\noindent
\textbf{1. В кое изречение е допусната лексикална грешка?} \\
\noindent
\textbf{А)} Бях приключил с обидчивата си гордост; бях приключил с всичките си дребни прояви на суетност и високомерие. \\
\textbf{Б)} Отдавна вече мирската суета не го трогваше, беше съсредоточен в намеренията си, които превъзхождаха възможностите на човешката природа. \\
\textbf{В)} Сънят и вцепенението отстъпиха място на някакво необичайно трескаво и безсмислено суетене. \\
\textbf{Г)} Такива хора трудно признават собствените си грешки - не заради перфекционизъм, а от накърнена суетливост и самолюбие.

\vspace{0.6cm}

\noindent
\textbf{2. В кое изречение е допусната лексикална грешка?} \\
\noindent
\textbf{А)} Вилата е без текуща вода и без електричество. \\
\textbf{Б)} Той беше човек, завладян от фикс идеи, потаен и вероятно параноик. \\
\textbf{В)} Изчезнал без предизвестие - това бе професионалният евфемизъм за убит. \\
\textbf{Г)} Всеки от десетте отпечатъка на човешките пръсти е уникален.

\vspace{0.6cm}

\noindent
\textbf{3. В кое изречение е допусната лексикална грешка?} \\
\noindent
\textbf{А)} Свещените писания са истории и легенди за опитите на човек да разбере потребността си от смисъл на собственото си съществуване. \\
\textbf{Б)} В музея имаше малка папирусна книжка, изписана с чудновати писмена, които се извиваха като змии. \\
\textbf{В)} Без писменост историите били предавани от уста на уста и обикновено украсявани. \\
\textbf{Г)} Йероглифната писменост съчетава в себе си изобразителните елементи с абстрактните.

\vspace{0.6cm}

\noindent
\textbf{4. В кое изречение е допусната лексикална грешка?} \\
\noindent
\textbf{А)} Той се приближи още и започна да се взира в неясните букви. \\
\textbf{Б)} Родителите й бяха починали преди години и единственият й стожер сега беше баба й. \\
\textbf{В)} Неделният следобед е най-чудесното време за мързелуване, книги, филми и лежерни разговори. \\
\textbf{Г)} Той е стоял извън литературните разпри на своето време и сам е наричал себе си дилетант.

\vspace{0.6cm}

\noindent
\textbf{5. В кое изречение НЕ е допусната лексикална грешка?} \\
\noindent
\textbf{А)} Аз успях да предизвикам интересна дилема на тази тема, която за съжаление прерасна в спор и завърши с разправия. \\
\textbf{Б)} Поради авария ТЕЦ „Бобов дол“ спира част от мощите си за десетина дни. \\
\textbf{В)} Домакинът се разположи в едно кресло край масичката до прозореца и отправи нехаен жест към госта си. \\
\textbf{Г)} В подобни случаи полицията има право дори да използва параноичен газ за разпръскване на тълпата.

\vspace{0.6cm}

\noindent
\textbf{6. В кое изречение НЕ е допусната лексикална грешка?} \\
\noindent
\textbf{А)} Кариерата му на писател започва през 1865 г. и той приема литературния прякор Марк Твен, с който е известен като автор на „Приключенията на Том Сойер“. \\
\textbf{Б)} Познавам подробно брат си и мога да заявя с увереност, че той не е способен на подобна постъпка. \\
\textbf{В)} Той е изконен сърбин, роден е в Сърбия, майка му и баща му са сърби. \\
\textbf{Г)} Запалил се беше доста късно по стрелбата с лък, но упорстваше и дори се беше явявал на състезания за аматьори.

\vspace{0.6cm}

\noindent
\textbf{7. В кой ред фразеологичните словосъчетания НЕ са синоними?} \\
\noindent
\textbf{А)} дните ми са преброени - пътник съм \\
\textbf{Б)} замайва ми се главата - изгубвам си ума \\
\textbf{В)} на куков ден - на второ пришествие \\
\textbf{Г)} с ченгел му вадят думите - вързан в устата

\vspace{0.6cm}

\noindent
\textbf{8. В кой ред фразеологичните словосъчетания НЕ са синоними?} \\
\noindent
\textbf{А)} вземам ума - свалям звезди от небето \\
\textbf{Б)} ям душата - пека на шиш \\
\textbf{В)} отпускам му края - удрям го на живот \\
\textbf{Г)} имам редки пръсти - пара не свъртам

\vspace{0.6cm}

\noindent
\textbf{9. В кой ред фразеологичните словосъчетания са синоними?} \\
\noindent
\textbf{А)} не подвивам крак - не паса трева \\
\textbf{Б)} като куче и котка - между Сцила и Харибда \\
\textbf{В)} вятър ме вее на бял кон - вятър работа \\
\textbf{Г)} като две капки вода - одрал му е кожата

\vspace{0.6cm}

\noindent
\textbf{10. В кой ред фразеологичните словосъчетания са синоними?} \\
\noindent
\textbf{А)} опитно зайче - заешка душа \\
\textbf{Б)} вдигам ръце - дигам си парцалите \\
\textbf{В)} бабини деветини - врели-некипели \\
\textbf{Г)} чиста ми е работата - света вода ненапита

\vspace{0.6cm}

\noindent
\textbf{11. В кой ред фразеологичните словосъчетания са антоними?} \\
\noindent
\textbf{А)} работи му ченето - слага си катинар на устата \\
\textbf{Б)} вземам ума - изгубвам си ума \\
\textbf{В)} вземам за чиста монета - гледам накриво \\
\textbf{Г)} тупам се в гърдите - преминавам през огън и вода

\vspace{0.6cm}

\noindent
\textbf{12. В кой ред фразеологичните словосъчетания НЕ са антоними?} \\
\noindent
\textbf{А)} чета конско евангелие - пея осанна \\
\textbf{Б)} вързан съм в ръцете - имам лека ръка \\
\textbf{В)} галеник на съдбата - роден на черен петък \\
\textbf{Г)} ум като бръснач - с два пръста чело

\vspace{0.6cm}

\noindent
\textbf{13. В кой ред фразеологичните словосъчетания НЕ са антоними?} \\
\noindent
\textbf{А)} цепя косъма на две - имам широки пръсти \\
\textbf{Б)} като църковна мишка - като крез \\
\textbf{В)} плаша се от сянката си - не ми мигва окото \\
\textbf{Г)} в черна светлина - на бял свят

\vspace{0.8cm}

\noindent
\textbf{14. За всяко празно място изберете най-уместния от предложените изрази.} \\
\noindent
Всяка сериозна дискусия относно ролята на войната в човешките общества, и преди всичко дебатите за или против нейната неизбежност, \makebox[2cm]{\dotfill} (А) ни отправят към въпроса за нейния произход. Знайно е, че организираната война е неизменен спътник на цивилизацията още от началото на писаната история, но това все пак не са повече от пет хиляди години. Ако военни действия съществуват още преди цивилизацията, та дори и преди човека, то тогава \makebox[2cm]{\dotfill} (Б) тревожната вероятност войната да се окаже неизбежна част от нашето генетично наследство. Подобна мисъл би била действително \makebox[2cm]{\dotfill} (В), но въпреки това сме длъжни да й \makebox[2cm]{\dotfill} (Г) достатъчно внимание. \\
\noindent
\textbf{А)} наложително, фатално, неминуемо \\
\textbf{Б)} изниква, се ражда, се показва \\
\textbf{В)} умопомрачителна, обезкуражаваща, предизвикваща \\
\textbf{Г)} посветим, дадем, удостоим

\vspace{0.6cm}

\noindent
\textbf{15. За всяко празно място изберете най-уместния от предложените изрази.} \\
\noindent
Личните взаимоотношения са \makebox[2cm]{\dotfill} (А) за нас, както въздухът, който дишаме. Всички ние се нуждаем от приятели, любими, партньори, от хора, с които да \makebox[2cm]{\dotfill} (Б) радости, мъки, опасения и успехи. Това общуване \makebox[2cm]{\dotfill} (В) най-нежните ни струни и обогатява и най-затънтените кътчета на душата ни. Всички се нуждаем от дружба, любов, добра компания и чувство за принадлежност и въпреки това често оставаме на разстояние един от друг, не желаейки или не можейки да протегнем ръка и да \makebox[2cm]{\dotfill} (Г) пълноценен контакт. \\
\noindent
\textbf{А)} основополагащи, жизненоважни, изконни \\
\textbf{Б)} разменим, доверим, споделим \\
\textbf{В)} докосва, свири на, обтяга \\
\textbf{Г)} осъществим, правим, сътворим

\vspace{0.6cm}

\noindent
\textbf{16. За всяко празно място изберете най-уместния от предложените изрази.} \\
\noindent
Литературата \makebox[2cm]{\dotfill} (А) мисленето, разширява хоризонта, приучава към правенето на връзки и сравнения, набавя ориентация, защото не е тайна, че по литературата може да се направят изводи за обществата, каквито невинаги могат да се \makebox[2cm]{\dotfill} (Б) от документите. Всяко произведение предполага противоречивост, смътност, неопределеност и по този начин винаги е \makebox[2cm]{\dotfill} (В) за ума. Учителите трябва да се стремят да \makebox[2cm]{\dotfill} (Г) познавателно и изследователско любопитство. \\
\noindent
\textbf{А)} форсира, аргументира, стимулира \\
\textbf{Б)} измъкнат, експлоатират, извлекат \\
\textbf{В)} обстоятелство, предизвикателство, свидетелство \\
\textbf{Г)} култивират, активират, импулсират

\vspace{0.6cm}

\noindent
\textbf{17. За всяко празно място изберете най-уместния от предложените изрази.} \\
\noindent
Докато се разхождаме из приветливите коридори, младежи играят спокойно на карти, а момиченца скачат на въже в прекрасния вътрешен двор. Минаваме от стая в стая и навсякъде децата учат: в голямата и светла библиотека, в лъскавата компютърна зала или просто в класната стая. Откриваме същата \makebox[2cm]{\dotfill} (А) обстановка и в учителската стая, където също не липсват компютри. Но най-невероятното нещо е \makebox[2cm]{\dotfill} (Б) на шум. Защо тези деца не крещят? Защо например не използват отсъствието на учителя, за да се замерят с пергели? Все пак те са в тази особено \makebox[2cm]{\dotfill} (В) възраст. Както всички младежи в Европа, и те имат пиърсинг, качулки, които нахлупват чак до носа, и джинси, завлечени до колена. Само че \makebox[2cm]{\dotfill} (Г) спокойни. \\
\noindent
\textbf{А)} уютна, доброкачествена, китна \\
\textbf{Б)} дефицитът, недостигът, липсата \\
\textbf{В)} опасна, фатална, преломна \\
\textbf{Г)} напомнят на, изглеждат, стават

\vspace{0.6cm}

\noindent
\textbf{18. За всяко празно място изберете най-уместния от предложените изрази.} \\
\noindent
Бъдеще, настояще и минало за мен са части от едно безкрайно време, в него онова, което наричаме минало, за миг се превръща в настояще или бъдеще и обратно! Ако светът \makebox[2cm]{\dotfill} (А) по-добре миналото, вярвам, че много от сегашните \makebox[2cm]{\dotfill} (Б) биха изчезнали. Трябва да живеем заедно, \makebox[2cm]{\dotfill} (В) етическите, етническите, историческите, литературните \makebox[2cm]{\dotfill} (Г). Историята показва, че никой не може да живее сам, затова трябва да живеем заедно. \\
\noindent
\textbf{А)} предугаждаше, познаваше, долавяше \\
\textbf{Б)} перипетии, поразии, недоразумения \\
\textbf{В)} съзерцавайки, съзнавайки, съзирайки \\
\textbf{Г)} различия, варианти, несъответствия

\vspace{0.6cm}

\noindent
\textbf{19. За всяко празно място изберете най-уместния от предложените изрази.} \\
\noindent
Аз съм тази, която се усмихва първа \makebox[2cm]{\dotfill} (А) да съм мила и любезна и съответно ми се отговаря със същото. Държа се добре с всички, това работи и постепенно би могло да стане \makebox[2cm]{\dotfill} (Б) и на останалите. Въпрос на лично \makebox[2cm]{\dotfill} (В), което, разбира се, отнема време. Живели сме дълго време намусени и искаме рязко всичко да се смени, но това са сложни \makebox[2cm]{\dotfill} (Г) в психиката на човека – да свикне да съществува в тези нови условия, не е лесно. \\
\noindent
\textbf{А)} старая се, силя се, определям \\
\textbf{Б)} ритуал, традиция, навик \\
\textbf{В)} превъзнасяне, превъзпитание, превъплъщаване \\
\textbf{Г)} похвати, процеси, качества

\vspace{0.6cm}

\noindent
\textbf{20. За всяко празно място изберете най-уместния от предложените изрази.} \\
\noindent
Желанието на учителите да научат на много неща претоварва учениците - това е един от основните проблеми на образованието. По някои предмети учителите \makebox[2cm]{\dotfill} (А) прекалено много детайли дори и от учениците, които нямат намерение да се занимават в бъдеще с дадената наука. Така реално не се получават знания, а нещо се назубря временно, колкото да мине изпитването. Повечето ученици днес искат да виждат \makebox[2cm]{\dotfill} (Б) в това, което учат. \makebox[2cm]{\dotfill} (В) „трябва да го знаеш, защото утре ще те изпитам“ отдавна вече не работи. Днешните деца повече от всякога имат нужда от посока, от цели и \makebox[2cm]{\dotfill} (Г), които да им помогнат да ги постигнат. \\
\noindent
\textbf{А)} претендират за, налагат, изискват \\
\textbf{Б)} резултат, смисъл, поанта \\
\textbf{В)} аргументът, заповедта, доказателството \\
\textbf{Г)} претекст, привилегия, мотивация

\newpage

\vspace{0.6cm}

\noindent
\textbf{ЧАСТ 2 (време за работа: 60 минути)}

\vspace{0.3cm}

\noindent
\textbf{21. Кой от посочените факти НЕ е свързан със стихотворението на Димчо Дебелянов „Аз искам да те помня все така...“?} \\
\noindent
\textbf{А)} Творбата има кръгова композиция. \\
\textbf{Б)} Имало е посвещение: „На Зв.“. \\
\textbf{В)} Заглавието е цитатно. \\
\textbf{Г)} Първата публикация е в сп. „Родна реч“.

\vspace{0.6cm}

\noindent
\textbf{22. Посочете кой е жанрът на всяка от творбите.}  \\
\noindent
\begin{minipage}[t]{0.6\textwidth}
\textbf{А)} „Аз искам да те помня все така...“ на Димчо Дебелянов \\
\textbf{Б)} „Колко си хубава!“ на Христо Фотев \\
\textbf{В)} „Посвещение“ на Петя Дубарова
\end{minipage}
\begin{minipage}[t]{0.35\textwidth}
1. изповед \\
2. елегия \\
3. възхвала
\end{minipage}

\vspace{0.3cm}
\noindent
А -- \makebox[1cm]{\dotfill} \hspace{0.5cm} Б -- \makebox[1cm]{\dotfill} \hspace{0.5cm} В -- \makebox[1cm]{\dotfill}

\vspace{0.6cm}

\noindent
\textbf{23. Посочете основния мотив, чрез който е представена темата за любовта в стихотворението на Димчо Дебелянов „Аз искам да те помня все така...“.} \\
\noindent
\textbf{А)} печалния спомен \\
\textbf{Б)} възвишената красота \\
\textbf{В)} бленуваната хармония \\
\textbf{Г)} споделената обич

\vspace{0.6cm}

\noindent
\textbf{24. Посочете кое НЕ може да се отнесе към Дебеляновата творба „Аз искам да те помня все така...“.} \\
\noindent
\textbf{А)} невъзможното щастие в любовта \\
\textbf{Б)} тъгата по неизживения живот \\
\textbf{В)} тъжните слова на девойката \\
\textbf{Г)} драматичната раздяла с любимата

\vspace{0.6cm}

\noindent
\textbf{25. Какво замества местоимението \textit{Тя} в творбата „Аз искам да те помня все така...“?} \\
\noindent
\textbf{А)} любимата \\
\textbf{Б)} любовта \\
\textbf{В)} самотата \\
\textbf{Г)} смъртта

\vspace{0.6cm}

\noindent
\textbf{26. Посочете мотивите, които са свързани в стихотворението „Колко си хубава!“ на Христо Фотев.} \\
\noindent
\textbf{А)} красотата и любовта \\
\textbf{Б)} красотата и доброто \\
\textbf{В)} споменът и раздялата \\
\textbf{Г)} споменът и обичта

\vspace{0.6cm}

\noindent
\textbf{27. Кое художествено средство НЕ присъства в цитата от стихотворението „Колко си хубава!“ на Христо Фотев?} \\
\textit{Колко са хубави ръцете ти.} \\
\textit{И нозете ти колко са хубави.} \\
\textit{И очите ти колко са хубави.} \\
\textit{И косите ти колко са хубави.} \\
\noindent
\textbf{А)} реторично възклицание \\
\textbf{Б)} анафора \\
\textbf{В)} синтактичен паралелизъм \\
\textbf{Г)} метафора

\vspace{0.6cm}

\noindent
\textbf{28. Стиховете „в ръка ми вплела пламнала ръка“ са от стихотворението:} \\
\noindent
\textbf{А)} „Колко си хубава!“ на Христо Фотев \\
\textbf{Б)} „Посвещение“ на Петя Дубарова \\
\textbf{В)} „Аз искам да те помня все така...“ на Димчо Дебелянов \\
\textbf{Г)} „До моето първо либе“ на Христо Ботев

\vspace{0.6cm}

\noindent
\textbf{29. Кой мотив НЕ е свързан с темата за любовта в стихотворението на Христо Фотев „Колко си хубава!“?} \\
\noindent
\textbf{А)} красотата на любимата \\
\textbf{Б)} очарованието на спомена \\
\textbf{В)} въздействието на любовта \\
\textbf{Г)} възхищението на влюбения

\vspace{0.6cm}

\noindent
\textbf{30. Кое твърдение НЕ е вярно за стихотворението „Колко си хубава!“ на Христо Фотев?} \\
\noindent
\textbf{А)} Творбата е възхвала на женската красота. \\
\textbf{Б)} Образът на любимата е идеализиран. \\
\textbf{В)} Съкровеното желание на Аза е да опази любовта си. \\
\textbf{Г)} Първата публикация е в стихосбирката „Баладично пътуване“.

\vspace{0.6cm}

\noindent
\textbf{31. Посочете НЕВЯРНОТО.} \\
\textbf{Специфичният поетически ритъм в стихотворението на Христо Фотев „Колко си хубава!“ се постига чрез:} \\
\noindent
\textbf{А)} анафори \\
\textbf{Б)} епифори \\
\textbf{В)} синтактичен паралелизъм \\
\textbf{Г)} метонимии

\vspace{0.6cm}

\noindent
\textbf{32. Кой смислов акцент НЕ се съдържа в повторението „обичай ме!“ в стихотворението на Христо Фотев „Колко си хубава!“?} \\
\noindent
\textbf{А)} съкровен призив \\
\textbf{Б)} изповедна молба \\
\textbf{В)} магично заклинание \\
\textbf{Г)} настойчив спомен

\vspace{0.6cm}

\noindent
\textbf{33. Посочете художественото средство, което се открива в стиха „И да си зла – обичай ме!“ от стихотворението на Христо Фотев „Колко си хубава!“.}  \\
\noindent
\textbf{А)} антитеза \\
\textbf{Б)} смислов парадокс \\
\textbf{В)} градация \\
\textbf{Г)} оксиморон

\vspace{0.6cm}

\noindent
\textbf{34. В стихотворението „Посвещение“ на Петя Дубарова драматичните изживявания на лирическата героиня са изразени чрез антитези. Посочете НЕВЯРНАТА.} \\
\noindent
\textbf{А)} между самотата и любовта \\
\textbf{Б)} между страха и смелостта \\
\textbf{В)} между бляна и спомена \\
\textbf{Г)} между греха и светостта

\vspace{0.6cm}

\noindent
\textbf{35. Образът на лирическата героиня в „Посвещение“ на Петя Дубарова е свързан с:} \\
\noindent
\textbf{А)} изживяването на хармонията в силата на любовните чувства \\
\textbf{Б)} примирението с превратностите на съдбата \\
\textbf{В)} раздвоението между чувство за вина и желанието за свобода \\
\textbf{Г)} признанието за необратимия трагичен край на любовта

\vspace{0.6cm}

\noindent
\textbf{36. Какъв е идейният акцент във финалните стихове на „Посвещение“ на Петя Дубарова?}  \\
\noindent
\textbf{А)} духовен порив за свобода \\
\textbf{Б)} мечта за духовна хармония \\
\textbf{В)} копнеж за духовно съзряване \\
\textbf{Г)} романтичен блян за любимия

\vspace{0.6cm}

\noindent
\textbf{37. Посочете НЕВЯРНОТО.} \\
\textbf{Основната идея, изразена чрез оксиморона в стиха „а в свойта вяра сам не вярвам аз“ от стихотворението на Димчо Дебелянов „Аз искам да те помня все така...“ е за:} \\
\noindent
\textbf{А)} раздялата на човека с надеждата \\
\textbf{Б)} липсата на ценностен ориентир \\
\textbf{В)} безсилието на човека пред съдбата \\
\textbf{Г)} вечното човешко страдание

\vspace{0.6cm}

\noindent
\textbf{38. Коментирайте в рамките на 4 – 5 изречения внушенията на стиховете „и любовта ни сякаш по е свята, / защото трябва да се разделим“ от стихотворението на Димчо Дебелянов „Аз искам да те помня все така...“.} \\
\noindent
\makebox[\textwidth]{\dotfill} \\
\makebox[\textwidth]{\dotfill} \\
\makebox[\textwidth]{\dotfill} \\
\makebox[\textwidth]{\dotfill} \\
\makebox[\textwidth]{\dotfill}

\vspace{0.6cm}

\noindent
\textbf{39. Разтълкувайте в рамките на 3 – 4 изречения символиката на луната в стихотворението „Посвещение“ на Петя Дубарова.} \\
\noindent
\makebox[\textwidth]{\dotfill} \\
\makebox[\textwidth]{\dotfill} \\
\makebox[\textwidth]{\dotfill} \\
\makebox[\textwidth]{\dotfill}

\vspace{0.6cm}

\noindent
\textbf{40. Напишете теза за есе върху стихотворението на Христо Фотев „Колко си хубава!“ на тема „Любовта – истина за смисъла на човешкото битие“.} \\
\noindent
\makebox[\textwidth]{\dotfill} \\
\makebox[\textwidth]{\dotfill} \\
\makebox[\textwidth]{\dotfill} \\
\makebox[\textwidth]{\dotfill} \\
\makebox[\textwidth]{\dotfill}
\newpage

\subsection{ОБОБЩИТЕЛЕН ТЕСТ КЪМ ДЯЛ 1 „ЛЮБОВТА“}
\noindent
\textit{Име и фамилия:} \dotfill \textit{№} \makebox[2cm]{\dotfill}

\vspace{0.6cm}

\begin{center}
{\large\bfseries ТЕСТ 4} \\
{\bfseries ОБОБЩИТЕЛЕН ТЕСТ КЪМ ДЯЛ 1 „ЛЮБОВТА“}
\end{center}

\vspace{0.5cm}

\noindent
\textbf{1. В кой от редовете \textbf{НЕ} е допусната правописна грешка?} \\
\textbf{А)} разложка, стоношка, нишка \\
\textbf{Б)} медийна, стадии, поразий \\
\textbf{В)} невинаги, навярно, незнаейки \\
\textbf{Г)} повърхностни, тягостни, безнадеждни

\vspace{0.6cm}

\noindent
\textbf{2. В кое от изреченията \textbf{НЕ} е допусната правописна грешка?} \\
\textbf{А)} Пожарникарът, който спасява света, е герой от роман, определен за бестселър № 1 за 2016 г. \\
\textbf{Б)} Повечето Троянски коне се стремят да откраднат с цел финансова печалба. \\
\textbf{В)} Познавам ги от дете и двете – майката е известна пианистка, а дъщеря й е виолончелистка. \\
\textbf{Г)} Тя е израснала в музикално семейство и затова още от малка свири виртуозно на цигулка.

\vspace{0.6cm}

\noindent
\textbf{3. В кое от изреченията \textbf{НЕ} е допусната граматична грешка при членуването?} \\
\textbf{А)} Каза ми, че през зимния сезон жълтите и червени ябълки са най-подходящи за диетата ми. \\
\textbf{Б)} Участвам с голям ентусиазъм в дарителската кампания на Университет по библиотекознание и информационни технологии. \\
\textbf{В)} Времето е най-безпристрастният съдник на събитията в човешката история, на делата и приноса на големите личности за благото на хората. \\
\textbf{Г)} Васил Иванов Кунчев, Апостолът на свободата, е един от идеолозите на българската национална революция.

\vspace{0.6cm}

\noindent
\textbf{4. В кое от изреченията \textbf{Е} допусната граматична грешка?} \\
\textbf{А)} Сали Яшар осъзнава силата на неговите харуци да носят радост и щастие в света, изпълнен с мъки и нещастия. \\
\textbf{Б)} Типичен пример за нравствен катарзис е разказът „Индже“, който поставя проблема за осъзнаване на жизнената цел на човека. \\
\textbf{В)} Очаквам и се надявам на вашето положително отношение към този проект, защото е свързан с развитието на общината. \\
\textbf{Г)} Би Би Си съобщават, че ще намалят разходите за уебсайта си с 25\% още другия месец, като въведат някои подобрения.

\vspace{0.6cm}

\noindent
\textbf{5. В кое от изреченията \textbf{НЕ} е допусната пунктуационна грешка?} \\
\textbf{А)} Няколко парчета диня, кората от 1 лимон, 3/4 ч. ч. лимонов сок, 1/2 ч. ч. светъл мед и един лимон на резени за украса, са достатъчни да подготвите освежаващата лимонада. \\
\textbf{Б)} Благодарение на усилията на по-заможни и просветени възрожденци, по време на Българското възраждане работят над 130 читалища. \\
\textbf{В)} В днешно време, често се говори за миграцията като личен избор, от една страна, и като спешна необходимост, от друга. \\
\textbf{Г)} Именно чрез събудената съвест човешкото съществуване придобива истински смисъл според В. Франкъл.

\vspace{0.6cm}

\noindent
\textbf{6. В кое от изреченията е допусната пунктуационна грешка?} \\
\textbf{А)} Мисълта с какво мога да допринеса на света, не ми даваше покой през всичките тези години. \\
\textbf{Б)} Тя е наясно, че, за да успее, трябва да я възприемат като уверена и безкомпромисна. \\
\textbf{В)} Припомних си мъдростта на Андерсен: „Можеш да видиш истинския облик на планината само когато се отдалечиш от нея“. \\
\textbf{Г)} Естествено, вярвам й, защото винаги е искрена с мен, а в трудни моменти е на първа линия.

\vspace{0.6cm}

\noindent
\textbf{7. Коя двойка фразеологични словосъчетания \textbf{НЕ} са синоними?} \\
\textbf{А)} от дън душа – от все сърце \\
\textbf{Б)} стоя в сянка – в глуха линия съм \\
\textbf{В)} между нас казано – между Сцила и Харибда \\
\textbf{Г)} хвърлям зад гърба си – махвам с ръка

\vspace{0.6cm}

\noindent
\textbf{8. От коя творба е посоченият откъс?} \\
Ела! Ще измиєм луната от грях! \\
Ще хвърлим трупа на умрелия страх, \\
ще пеем с тъмнежния корабен глас \\
на морската нощ във добрия Бургас. \\

\vspace{-0.1cm}

\noindent
\begin{tabular}{@{}p{0.48\textwidth} p{0.48\textwidth}@{}}
\textbf{А)} „Две хубави очи“ & \textbf{В)} „Колко си хубава!“ \\
\textbf{Б)} „Аз искам да те помня все така“ & \textbf{Г)} Посвещение
\end{tabular}

\vspace{0.6cm}

\noindent
\textbf{9. Кое от посочените твърдения \textbf{НЕ} е вярно за междутекстовия диалог на „Аз искам да те помня все така“, „Колко си хубава!“ и „Посвещение“?} \\
\textbf{А)} И трите стихотворения са създадени през годините, в които властва комунистическата идеология, а любовната тема не е сред престижните в литературата. \\
\textbf{Б)} „Аз искам да те помня все така“ на Димчо Дебелянов е различна от другите творби като настроение и в жанрово отношение. \\
\textbf{В)} И към трите лирически произведения има посвещения, но трите посвещения имат различно значение за разбирането на творбите. \\
\textbf{Г)} И в трите творби лирическите герои се обръщат към любимия човек и внушават съкровена интимност на споделените чувства.

\vspace{0.6cm}

\noindent
\textbf{10. Кое е характерно за съотнасянето на любовта с времето в трите творби? Посочете \textbf{НЕВЯРНОТО} твърдение.} \\
\textbf{А)} Само една от творбите разкрива моменти от миналото на обичания човек. \\
\textbf{Б)} В Дебеляновата творба любовта е представена като спомен в настоящето. \\
\textbf{В)} Според лирическата героиня на Петя Дубарова любовта ще стане реалност в бъдещето. \\
\textbf{Г)} Според героя на Фотев любовта е силна само в мига на настоящето, а бъдещето е несигурно.

\vspace{0.6cm}

\noindent
\textbf{11. В коя от творбите „Аз искам да те помня все така“, „Колко си хубава!“ и „Посвещение“ присъстват повелителни глаголни форми като израз на императивност на посланието или като молба?} \\

\vspace{-0.1cm}

\noindent
\begin{tabular}{@{}p{0.48\textwidth} p{0.48\textwidth}@{}}
\textbf{А)} „Аз искам да те помня все така“ & \textbf{В)} „Посвещение“ \\
\textbf{Б)} във всяка от изброените & \textbf{Г)} в нито една от изброените
\end{tabular}

\vspace{0.6cm}

\noindent
\textbf{12. В коя творба любовта е изцяло духовна и протича в призрачна символистична атмосфера?} \\

\vspace{-0.1cm}

\noindent
\begin{tabular}{@{}p{0.48\textwidth} p{0.48\textwidth}@{}}
\textbf{А)} „Посвещение“ & \textbf{В)} „Аз искам да те помня все така“ \\
\textbf{Б)} „Колко си хубава!“ & \textbf{Г)} в нито една от изброените
\end{tabular}

\vspace{0.6cm}

\noindent
\textbf{13. В коя от посочените двойки творби образът на нощта \textbf{НЕ} се свързва единствено със страха и самотата?} \\

\vspace{-0.1cm}

\noindent
\begin{tabular}{@{}p{0.58\textwidth} p{0.58\textwidth}@{}}
\textbf{А)} „Посвещение“ и „Колко си хубава!“ & \textbf{В)} „Посвещение“ и „В часа на синята мъгла“ \\
\textbf{Б)} „Аз искам да те помня все така“ и „Колко си хубава!“ & \textbf{Г)} във всички посочени двойки
\end{tabular}

\vspace{0.6cm}

\noindent
\textbf{14. Коя от творбите е написана в бели стихове?} \\

\vspace{-0.1cm}

\noindent
\begin{tabular}{@{}p{0.48\textwidth} p{0.48\textwidth}@{}}
\textbf{А)} „Аз искам да те помня все така“ & \textbf{В)} „Посвещение“ \\
\textbf{Б)} „Колко си хубава!“ & \textbf{Г)} всички посочени творби
\end{tabular}

\vspace{0.6cm}

\noindent
\textbf{15. Трима ученици разискват образа на мъжа в трите текста. Изберете това твърдение, което съвпада с вашата гледна точка, и посочете 2 аргумента, за да го защитите.}

\vspace{0.3cm}

\noindent
\begin{tabularx}{\linewidth}{|>{\centering\arraybackslash}X|>{\centering\arraybackslash}X|>{\centering\arraybackslash}X|}
\hline
Образът на мъжа е различен в трите творби. & В „Посвещение“ и „Колко си хубава“ образът на мъжа има общи черти. & Образът на мъжа има общи черти в трите творби. \\
\hline
\end{tabularx}

\vspace{0.4cm}

\noindent
\textbf{Аргументи:} \\
\noindent\rule{\linewidth}{0.4pt}
\noindent\rule{\linewidth}{0.4pt}
\noindent\rule{\linewidth}{0.4pt}
\noindent\rule{\linewidth}{0.4pt}
\noindent\rule{\linewidth}{0.4pt}
\noindent\rule{\linewidth}{0.4pt}
\noindent\rule{\linewidth}{0.4pt}
\noindent\rule{\linewidth}{0.4pt}
\noindent\rule{\linewidth}{0.4pt}
\noindent\rule{\linewidth}{0.4pt}
\noindent\rule{\linewidth}{0.4pt}
\noindent\rule{\linewidth}{0.4pt}

\vspace{0.6cm}

\noindent
\textbf{Прочетете двата текста и изпълнете задачи 16 – 23.}

\vspace{0.4cm}

\noindent
\textbf{Текст 1} \\
В Южно Мароко има вихрушка, гаджеж, от която фелахите се защитават с ножове. Вятърът африко е достигнал понякога до Рим. Алм е есенен, идващ от Югославия. А арифи, или ареф, или рифи, изпепелява всичко по пътя си с хиляди огнени езици. Това са постоянни ветрове, които живеят в сегашно време.

Има и други, непостоянни, дето сменят посоката си, събират конете си и се обръщат срещу движението на часовниковата стрелка. Бист роз върлува сто седемдесет и седем дни в Афганистан, поми-буря, която преоблича всичко в яркожълто и си отива, последвана от дъжд. Харматан духа, докато издуха и себе си в Атлантическия океан. Имбат е морски бриз от Северна Африка. Има ветрове, които само въздишат към небето. Има нощни пясъчни бури, които носят студа със себе си. [...] Дату идва откъм Гибралтар и донася ухания.

Познат е и тайният вятър на пустинята. Неговото име било завинаги заличено от владетеля, чийто син попаднал във вихъра му и там намерил гибелта си. Нафхат е силен вятър от Арабия. Бешабар е чер и сух, „черен вятър“, и идва от североизток, от Кавказ.

\vspace{0.2cm}
\hfill \textit{Из „Английският пациент“ от М. Ондатджи}

\vspace{0.4cm}

\noindent
\textbf{Текст 2} \\
Четири човешки събии насред ужасите на Втората световна война. Сплотени от болката и желанието за оцеляване. Събрани от съдбата и войната и спасени от любовта.

През 1992 г. Майкъл Ондатджи публикува романа „Английският пациент“. В творбата писателят успява да събере четири човешки съдби на хора, наранени от Втората световна война. Действието се развива в Италия – в изолирана тосканска вила в края на 1944 и началото на 1945 г. Романът е изграден основно върху две сюжетни линии и те се пресичат в „английския пациент“ – човек, който е обезорал до неузнаваемост. Първата разказва за любовта на английския пациент – унгарския граф Алмаси. Той изследва либийската пустиня и има връзка с Катрин Клифтън – съпругата на друг изследовател.

Втората сюжетна линия разказва за Хана – двайсетгодишна канадецка медицинска сестра, която се грижи за обегорелия от самолетна катастрофа и умиращ пациент. По-късно към тях се присъединяват Караваджо – крадец, шпионин, измъчван и осакатен от военнопрестъпници, приятел на баща ѝ, и Кип – млад сикх и салчър, който през деня разчиства района от немски бомби и мини, а през нощта е с Хана.

Алмаси е наречен „английският пациент“, тъй като документите за самоличността му са загубени, когато катастрофира със самолета си в северноафриканската пустиня. Обезорял много лошо, той е спасен от бедуини. Не може да си спомни кой е, но е лекуван като британски офицер. Причините са, че говори английски и очевидно е с високо социално положение.

Романът е носител на наградата „Букър“ за 1992 г. През 1996 г. по книгата е заснет и едноименен филм с участието на Ралф Файнс, Кристин Скот Томас и Жулиет Бинош. Филмът е отличен с две награди „Златен глобус“ и девет награди „Оскар“.

\vspace{0.2cm}
\hfill \textit{Из интернет}

\vspace{0.6cm}

\noindent
\textbf{16. Каква функция изпълнява Текст 1?} \\
\textbf{А)} Посочва конкретното мястото и времето на действие в романа. \\
\textbf{Б)} Пресъздава експресивно духа на пустинята чрез ветровете в нея. \\
\textbf{В)} Констатира последиците от ветровете в африканската пустиня. \\
\textbf{Г)} Представя научна информация за метеорологичните условия в Африка.

\vspace{0.6cm}

\noindent
\textbf{17. Използваният в Текст 1 израз „които живеят в сегашно време“ е употребен със значението:} \\
\textbf{А)} които не се употребяват с други глаголни времена \\
\textbf{Б)} които в бъдещето се изгубват без следа \\
\textbf{В)} които никога не престават \\
\textbf{Г)} които непрекъснато се променят

\vspace{0.6cm}

\noindent
\textbf{18. Кое от посочените твърдения е вярно според Текст 1?} \\
\textbf{А)} Представени са два вида ветрове – постоянни и непостоянни. \\
\textbf{Б)} Ветровете имат по едно име, а само един вятър няма име. \\
\textbf{В)} Ветровете създават само проблеми за хората. \\
\textbf{Г)} Ветровете не могат да прекосяват континентите.

\vspace{0.6cm}

\noindent
\textbf{19. В коя сфера на общуване се използва Текст 2?} \\

\vspace{-0.1cm}

\noindent
\begin{tabular}{@{}p{0.48\textwidth} p{0.48\textwidth}@{}}
\textbf{А)} в битовата сфера & \textbf{В)} в научната сфера \\
\textbf{Б)} в естетическата сфера & \textbf{Г)} в медийната сфера
\end{tabular}

\vspace{0.6cm}

\noindent
\textbf{20. Кое от посочените твърдения се отнася за Текст 2?} \\

\vspace{-0.1cm}

\noindent
\begin{tabular}{@{}p{0.48\textwidth} p{0.48\textwidth}@{}}
\textbf{А)} Текстът представлява реклама с апелативна функция. & \textbf{В)} Текстът е изчерпателен и емоционално въздействащ. \\
\textbf{Б)} Текстът включва кратко резюме на книгата. & \textbf{Г)} Текстът е оформен стандартизирано.
\end{tabular}

\vspace{0.6cm}

\noindent
\textbf{21. Запишете буквите САМО на онези 4 факта, които НЕ присъстват в Текст 1 и Текст 2.} \\
\textbf{А)} Суданската пясъчна буря е в жълти цветове и е последвана от дъжд. \\
\textbf{Б)} Ветровете са назовани със своите имена, които ги идентифицират. \\
\textbf{В)} Хабуб е горещ сух вятър от Тунис, който изостря нервите. \\
\textbf{Г)} Ветровете, които въздишат към небето, са много слаби. \\
\textbf{Д)} Основната сюжетна линия в романа на Ондатджи е за любовта на Алмаси и Хана. \\
\textbf{Е)} Хана – двайсетгодишната медицинска сестра, се грижи за умиращия пациент в Канада. \\
\textbf{Ж)} В Мароко има вихрушка, от която местните се отбраняват с оръжие. \\
\textbf{З)} Караваджо е шпионин, осакатен от военнопрестъпници, приятел на Алмаси.

\vspace{0.6cm}

\noindent
\textbf{22. Относно това дали Текст 1 е подходящ откъс, който може да се добави към анотацията на романа „Английският пациент“ (Текст 2), двама ученици изказват различно мнение.} \\
\textit{Първо мнение:} Напълно подходящ – интригуващо описание! \\
\textit{Второ мнение:} Не мисля, че е подходящ – по-скоро обърква читателя!

\vspace{0.2cm}
\noindent
\textbf{Изберете една от позициите и я защитете с 2 аргумента.}

\vspace{0.3cm}
\noindent\rule{\linewidth}{0.4pt}
\noindent\rule{\linewidth}{0.4pt}
\noindent\rule{\linewidth}{0.4pt}
\noindent\rule{\linewidth}{0.4pt}
\noindent\rule{\linewidth}{0.4pt}
\noindent\rule{\linewidth}{0.4pt}
\noindent\rule{\linewidth}{0.4pt}
\noindent\rule{\linewidth}{0.4pt}
\noindent\rule{\linewidth}{0.4pt}
\noindent\rule{\linewidth}{0.4pt}

\vspace{0.6cm}

\noindent
\textbf{23. В рамките на 3–4 изречения формулирайте СОБСТВЕНА теза на тема \textit{Името и неговата сила}.}

\vspace{0.3cm}
\noindent\rule{\linewidth}{0.4pt}
\noindent\rule{\linewidth}{0.4pt}
\noindent\rule{\linewidth}{0.4pt}
\noindent\rule{\linewidth}{0.4pt}
\noindent\rule{\linewidth}{0.4pt}
\noindent\rule{\linewidth}{0.4pt}
\noindent\rule{\linewidth}{0.4pt}
\noindent\rule{\linewidth}{0.4pt}
\noindent\rule{\linewidth}{0.4pt}
\noindent\rule{\linewidth}{0.4pt}

\vspace{0.6cm}

\noindent
\textbf{24. Запишете правилните форми на думите, поставени в скоби.} \\
\textbf{А) Членуване} \\
(Проект) \dotfill на Постановление на (Министерски) \dotfill съвет за приемане на \\
Устройствен правилник на Агенция „Митници“ бе дискутиран продължително. \\
С първите лъчи на слънцето той откри невероятна гледка в (пустош) \dotfill. \\
Ето го (учебник) \dotfill по който се готвих за изпита.

\vspace{0.4cm}

\noindent
\textbf{Б) Местоимения} \\
Няма да чакам (никой/никого) \dotfill повече, защото достатъчно се забавихме. \\
Той беше упорит състезател, (чийто/чиито) \dotfill победи ме възхищаваха. \\
Тя ме впечатли не толкова със зашеметяващата си фигура, колкото с магията на (нейните/своите) \dotfill очи.

\vspace{0.4cm}

\noindent
\textbf{В) Учтива форма} \\
Госпожо, ако сте (бил) \dotfill (щастлив) \dotfill тук, бихте ли се (върнал) \dotfill отново при нас?

\vspace{0.6cm}

\noindent
\textbf{25. Препишете текста, като редактирате допуснатите 6 грешки.} \\
Поетът с ватенката, певецът на новото време – нито едно от тези популярни определения не е точно за човека
Пеньо Пенев, нито едно от тях не разкрива дълбоката му душевна драма, житейските сътресения, трагизмът и обреченоста му. Той е сред най-известните ни поети от първото следвоенно поколение творци. Днес обаче малцина знаят
Пеньо Пеневата поема Дни на проверка само защото се отнася към едно друго време – към зората на тоталитарното
строителство.

\vspace{0.3cm}
\noindent\rule{\linewidth}{0.4pt}
\noindent\rule{\linewidth}{0.4pt}
\noindent\rule{\linewidth}{0.4pt}
\noindent\rule{\linewidth}{0.4pt}
\noindent\rule{\linewidth}{0.4pt}
\noindent\rule{\linewidth}{0.4pt}
\noindent\rule{\linewidth}{0.4pt}
\noindent\rule{\linewidth}{0.4pt}
\noindent\rule{\linewidth}{0.4pt}
\noindent\rule{\linewidth}{0.4pt}
\noindent\rule{\linewidth}{0.4pt}
\noindent\rule{\linewidth}{0.4pt}

\vspace{0.6cm}

\noindent
\textbf{Прочетете откъсите от творбите и изпълнете задачи 26 – 30.}

\vspace{0.3cm}

\noindent
\begin{tabularx}{\linewidth}{|>{\raggedright\arraybackslash}X|>{\raggedright\arraybackslash}X|}
\hline
\textbf{„Аз искам да те помня все така“} & \textbf{„Колко си хубава!“} \\
\hline
А все по-страшно пада нощ над нас, & Не ме изгубвай никога – откривай ме, \\
чертаят мрежа прилепите в мрака, & изпълвай ме с детинско изумление. \\
утеха сетна твоята немощ чака, & Отново да се уверя в ръцете ти, \\
а в своята вяра сам не вярвам аз. \textbf{(А)} & В нозете ти, в очите ти... Обичай ме! \\
И ти отпущаш пламнала ръка & Как искам да те задържа завинаги. \\
и тръгваш, поглед в тъмнината впила, & Да ме обичам завинаги – \\
изгубила дори за сълзи сила. – & завинаги. \\
Аз искам да те помня все така... & И колко ми е невъзможно... \textbf{(Б)} Колко си \\
& ти пясъчна... И моля те, не казвай ми, \\
& че искаш да ме задържиш завинаги, \\
& да ме обичаш завинаги, \\
& завинаги. \\
\hline
\end{tabularx}

\vspace{0.6cm}

\noindent
\textbf{26. Запишете 3 образа или мотива, които присъстват и в двата откъса.}\\
\noindent\rule{\linewidth}{0.4pt}
\noindent\rule{\linewidth}{0.4pt}
\noindent\rule{\linewidth}{0.4pt}
\noindent\rule{\linewidth}{0.4pt}
\noindent\rule{\linewidth}{0.4pt}
\noindent\rule{\linewidth}{0.4pt}
\noindent\rule{\linewidth}{0.4pt}
\noindent\rule{\linewidth}{0.4pt}
\noindent\rule{\linewidth}{0.4pt}

\vspace{0.6cm}

\noindent
\textbf{27. Със СВОИ ДУМИ посочете 2 примера, които подкрепят твърдението, че в двата откъса мотивът за раздялата е интерпретиран по различен начин.} \\
\noindent\rule{\linewidth}{0.4pt}
\noindent\rule{\linewidth}{0.4pt}
\noindent\rule{\linewidth}{0.4pt}
\noindent\rule{\linewidth}{0.4pt}
\noindent\rule{\linewidth}{0.4pt}
\noindent\rule{\linewidth}{0.4pt}
\noindent\rule{\linewidth}{0.4pt}
\noindent\rule{\linewidth}{0.4pt}
\noindent\rule{\linewidth}{0.4pt}
\noindent\rule{\linewidth}{0.4pt}

\vspace{0.6cm}

\noindent
\textbf{28. Каква е функцията на образа на ръката/ръцете в двата текста? Посочете 2 примера.} \\
\noindent\rule{\linewidth}{0.4pt}
\noindent\rule{\linewidth}{0.4pt}
\noindent\rule{\linewidth}{0.4pt}
\noindent\rule{\linewidth}{0.4pt}
\noindent\rule{\linewidth}{0.4pt}
\noindent\rule{\linewidth}{0.4pt}
\noindent\rule{\linewidth}{0.4pt}
\noindent\rule{\linewidth}{0.4pt}

\vspace{0.6cm}

\noindent
\textbf{29. Каква е функцията на образа на очите/погледа в двата текста? Посочете 2 примера.}\\
\noindent\rule{\linewidth}{0.4pt}
\noindent\rule{\linewidth}{0.4pt}
\noindent\rule{\linewidth}{0.4pt}
\noindent\rule{\linewidth}{0.4pt}
\noindent\rule{\linewidth}{0.4pt}
\noindent\rule{\linewidth}{0.4pt}
\noindent\rule{\linewidth}{0.4pt}
\noindent\rule{\linewidth}{0.4pt}

\vspace{0.6cm}

\noindent
\textbf{30. Определете какви са подчертаните изразни средства в откъсите, и разтълкувайте техния смисъл.} \\
\textbf{А)} \dotfill
\vspace{0.2cm}
\noindent\rule{\linewidth}{0.4pt}
\noindent\rule{\linewidth}{0.4pt}
\noindent\rule{\linewidth}{0.4pt}

\vspace{0.3cm}

\noindent
\textbf{Б)} \dotfill
\vspace{0.2cm}
\noindent\rule{\linewidth}{0.4pt}
\noindent\rule{\linewidth}{0.4pt}
\noindent\rule{\linewidth}{0.4pt}

\vspace{0.6cm}

\noindent
\textbf{31. Коментирайте в 3 – 4 изречения следния цитат: „Изтрива имената! Изтрива националностите! Това беше духът на пустинята“ (Майкъл Ондатджи), като имате предвид цитираните в теста и изучените художествени творби.}\\

\noindent\rule{\linewidth}{0.4pt}
\noindent\rule{\linewidth}{0.4pt}
\noindent\rule{\linewidth}{0.4pt}
\noindent\rule{\linewidth}{0.4pt}
\noindent\rule{\linewidth}{0.4pt}
\noindent\rule{\linewidth}{0.4pt}
\noindent\rule{\linewidth}{0.4pt}
\noindent\rule{\linewidth}{0.4pt}
\noindent\rule{\linewidth}{0.4pt}

\newpage

\subsection{ТЕСТ 5}

\noindent
\textit{Име и фамилия:} \dotfill \textit{№} \makebox[2cm]{\dotfill}

\vspace{0.6cm}

\begin{center}
{\large\bfseries ТЕСТ 5}
\end{center}

\vspace{0.3cm}

\noindent
\textbf{І МОДУЛ}

\vspace{0.2cm}

\noindent
\textbf{Прочетете двата текста и изпълнете задачите към тях (от 1. до 5. включително).}

\vspace{0.6cm}

\noindent
\textbf{ТЕКСТ 1} \\
Има дни, в които не мога да пиша. Има цели седмици, в които не мога да пиша. Това нямаше да е такъв проблем, ако можех предварително да определя кои ще са тези дни и седмици. Щях да намаля щетите и да поема на маратон по Темза. Или пък да хвана влака за обиколка по галериите... Но за съжаление, разбирам, че един ден е такъв, когато по обяд осъзнавам, че съм се взирал с часове в празния екран. Или съм драскал някаква тежка, неубедителна проза, която трябва да изтрия след това.

Мисля, че проблемът е в това, че не съм ужасно добър писател. Но съм много хладнокръвен и прецизен редактор (който по стечение на обстоятелствата е женен за дори по-добър редактор). Освен това съм безцеремонен, когато трябва да изтрия някакъв слаб текст. Всъщност трия поне три четвърти от това, което пиша. Прецеждам и прецеждам. И някъде между версии номер 15 и 25 нещо се случва. Стигам дотам, че когато четеш собствените си думи, вече имаш чувство, че някой друг ги е писал...

Понякога, когато не се чувствам много уверен, прекарвам сутрини, работейки в някое квартално кафене. Тълпата ме накара да се чувствам като част от света на заетите хора. В такива моменти препрочитам написаното в предишните дни, полирам, изразявам отново себе си.

Писането е като правене на операция или като пилотиране – трябва да работят всички тилиндри. Иначе ще се наложи да правиш нещо друго. Затова, когато усетя, че не става, се прибирам у дома, пускам пералнята, рисувам или просто правя крос по Темза.

Не мисля, че бих могъл да пиша, ако не чувах заглушено в главата си онзи глас, който ми казва пак и пак: „Това не е достатъчно добро... Пиши повече, пиши по-добре... Време лети“.

\vspace{0.2cm}
\hfill \textit{Марк Хадън, автор на „Странна куче с куче през нощта“, „Петното“ и „Червената къща“}

\vspace{0.6cm}

\noindent
\textbf{ТЕКСТ 2} \\
Какво друго е истинският автор на фрагменти освен човек, способен да обобщи частица от световния порядък в една-едничка мисъл. Колкото по-близо стига до едничката мисъл, толкова е по-достоен за похвала. И би бил съвършен, ако попада винаги в центъра на завършения и концентриран изказ. Гледна точка, от която собствената ми работа върху тази книга с фрагменти ме разочарова.

Аз губя много време, много място там, където би трябвало да оставя само категоричен знак; думите ми не са способни да обхванат енергията, изстреляна от съзнанието в залп, да изградят от нея ясна и съвсем прибрана звукова и смислова постройка. Достатъчно е да оня, който познава истинската сила на думите – една безкрайна тайна, – да си послужи с пет или шест от тях: слънцето изгрява, защото великолепието му е обяснено, а е възможно и да ни затиснат тъмни грамади, пороища и страх. Вярно е, десеторно повече думи биха могли да кажат същото, но и с десет пъти по-малка сила.

В странни посоки може да те отведе такова чувство на вина. Откриеш ли идеята си (върху която, без съмнение, ще се разпростре ненужно), отхрувисто и блестящо изложена у друг, преминаваща високо над главите като звуци от сребърен тромпет, ти си готов да останеш с вдигнати очи и да благодариш. Понесла те е люлката на разума, изплетена така отдавна за човечеството; собствената ти мисъл е поднесена вече съвършено и ти се чувстваш свободен от голяма грижа. Чувстваш се закрилян, ученик, дете...

Днешният човек, уж пестелив към времето си, уж галопиращ през живота, е забравил (именно защото галопира) що е истински размисъл и концентрация. Точно затова е многословен.

\vspace{0.2cm}
\hfill \textit{Димитър Коруджиев, „Болката, когато е прекрасна“}

\vspace{0.6cm}

\noindent
\textbf{1. Кое от посочените твърдения НЕ се отнася до Текст 1?} \\
\textbf{А)} За да пише добре, един писател се нуждае от разговор със своята съвест. \\
\textbf{Б)} Писателят често е твърде критичен към текстовете, които смята за слаби. \\
\textbf{В)} Когато чете свои произведения, творецът усеща, че у него се пробужда силно вдъхновение. \\
\textbf{Г)} Ако можеш да предвидиш ефективността на творческата работа, писателят би имал време и за други занимания.

\vspace{0.6cm}

\noindent
\textbf{2. Защо според Текст 2 съвременният човек трудно може да оформи мисълта си в синтезиран вид?} \\
\textbf{А)} Защото изпитва огорчение, че други преди него вече са развили неговите идеи. \\
\textbf{Б)} Защото динамиката на живота е променила способността му за концентрация. \\
\textbf{В)} Защото е загубил детската си непосредственост и спонтанност на изказа. \\
\textbf{Г)} Защото е изгубил способността си да познава истинската сила на думите.

\vspace{0.6cm}

\noindent
\textbf{3. Обяснете защо според Текст 1 писането е като правене на операция или пилотиране.}
\noindent\rule{\linewidth}{0.4pt}
\noindent\rule{\linewidth}{0.4pt}
\noindent\rule{\linewidth}{0.4pt}
\noindent\rule{\linewidth}{0.4pt}
\noindent\rule{\linewidth}{0.4pt}
\noindent\rule{\linewidth}{0.4pt}
\noindent\rule{\linewidth}{0.4pt}
\noindent\rule{\linewidth}{0.4pt}

\vspace{0.6cm}

\noindent
\textbf{4. Като имате предвид информацията в Текст 1 и Текст 2, запишете ДВА аргумента в подкрепа на твърдението, че процесът на писане изисква от твореца да бъде критичен и взискателен към себе си.}
\noindent\rule{\linewidth}{0.4pt}
\noindent\rule{\linewidth}{0.4pt}
\noindent\rule{\linewidth}{0.4pt}
\noindent\rule{\linewidth}{0.4pt}
\noindent\rule{\linewidth}{0.4pt}
\noindent\rule{\linewidth}{0.4pt}
\noindent\rule{\linewidth}{0.4pt}
\noindent\rule{\linewidth}{0.4pt}
\noindent\rule{\linewidth}{0.4pt}
\noindent\rule{\linewidth}{0.4pt}
\noindent\rule{\linewidth}{0.4pt}
\noindent\rule{\linewidth}{0.4pt}

\vspace{0.6cm}

\noindent
\textbf{5. Един от важните индикатори за ценностните нагласи на младите хора е свързан с четенето на книги. Графиката онагледява данни от Годишния доклад за младежта за 2019 г., изготвен по поръчка на Министерството на младежта и спорта, свързани с този индикатор.}

\vspace{0.3cm}

\begin{center}
\textbf{Четете ли или не четете някакви книги в момента?} \\
Спонтанен отговор. База: 1400 души

\vspace{0.3cm}

\usetikzlibrary{patterns}

\begin{tikzpicture}[scale=1.5]
    \fill[lightgray] (0,0) -- (3,0) arc (0:244.8:3) -- cycle;
    \fill[pattern=north east lines, pattern color=gray!60] (0,0) -- (244.8:3) arc (244.8:252:3) -- cycle;
    \fill[darkgray] (0,0) -- (252:3) arc (252:360:3) -- cycle;
    
    \begin{scope}[shift={(3.8, 1.2)}]
        \fill[darkgray] (0,0.4) rectangle (0.3,0.7);
        \node[right] at (0.4,0.55) {Чета (30\%)};
        
        \fill[lightgray] (0,-0.1) rectangle (0.3,0.2);
        \node[right] at (0.4,0.05) {Не чета (68\%)};
        
        \fill[pattern=north east lines, pattern color=gray!60] (0,-0.6) rectangle (0.3,-0.3);
        \draw (0,-0.6) rectangle (0.3,-0.3);
        \node[right] at (0.4,-0.45) {Без отговор (2\%)};
    \end{scope}
\end{tikzpicture}
\end{center}

\vspace{0.3cm}

\noindent
\textbf{Като имате предвид данните от графиката, коментирайте в кратък текст от 4–5 изречения следното:}
\begin{itemize}[noitemsep, topsep=0pt, parsep=0pt, itemsep=0.4em]
    \item Кой е преобладаващият отговор? На какво според вас се дължи той?
    \item Обяснете към коя група бихте се причислили и защо.
    \item Смятате ли, че отразената в графиката тенденция е тревожна?
\end{itemize}

\vspace{0.3cm}
\noindent\rule{\linewidth}{0.4pt}
\noindent\rule{\linewidth}{0.4pt}
\noindent\rule{\linewidth}{0.4pt}
\noindent\rule{\linewidth}{0.4pt}
\noindent\rule{\linewidth}{0.4pt}
\noindent\rule{\linewidth}{0.4pt}
\noindent\rule{\linewidth}{0.4pt}
\noindent\rule{\linewidth}{0.4pt}
\noindent\rule{\linewidth}{0.4pt}
\noindent\rule{\linewidth}{0.4pt}
\noindent\rule{\linewidth}{0.4pt}
\noindent\rule{\linewidth}{0.4pt}
\noindent\rule{\linewidth}{0.4pt}
\noindent\rule{\linewidth}{0.4pt}

\vspace{0.6cm}

\noindent
\textbf{6. Къде е допусната правописна грешка?} \\
\begin{tabular}{@{}p{0.48\textwidth} p{0.48\textwidth}@{}}
\textbf{А)} изсвирвам & \textbf{В)} смръчава се \\
\textbf{Б)} подтиквам & \textbf{Г)} унесен
\end{tabular}

\vspace{0.6cm}

\noindent
\textbf{7. Къде НЕ е допусната правописна грешка?} \\
\begin{tabular}{@{}p{0.48\textwidth} p{0.48\textwidth}@{}}
\textbf{А)} вестникар, сблъсак, подплатен & \textbf{В)} непоколебим, безнадежност, сфера \\
\textbf{Б)} месност, надпис, сбръсква се & \textbf{Г)} съкращение, звездно, установявам
\end{tabular}

\vspace{0.6cm}

\noindent
\textbf{8. В кое изречение НЕ е допусната правописна грешка?} \\
\textbf{А)} Жълто-червен флаг на ивици, развян по време на рали, предупреждава за хлъзгав участък, наличие на отломки или масло на пътя вследствие на инцидент. \\
\textbf{Б)} През лятото ще работя като помощник редактор в едно дигитално списание за музика. \\
\textbf{В)} За рождения си ден Иван получи страхотен три измерен пъзел от своите родители. \\
\textbf{Г)} Един от компонентите на този комплекс от науки е проучвателната гео-физика, която изучава структурата на Земята.

\vspace{0.6cm}

\noindent
\textbf{9. В кое от изреченията НЕ е допусната граматична грешка?} \\
\textbf{А)} От магазина купих килограм праскови, няколко пъпеши и само два големи ананаса. \\
\textbf{Б)} Наградиха и коня, с който спечелих международното състезание миналия месец. \\
\textbf{В)} След като се прибрах, установих, че автомобилът ми го няма. \\
\textbf{Г)} Господине, бихте ли ми услужил с една монета от 2 лева?

\vspace{0.6cm}

\noindent
\textbf{10. В коя от подчертаните думи НЕ е допусната граматична грешка?} \\
Срещнах \underline{учуденият} \textbf{(А)} ѝ поглед, след като успяхме да \underline{обясниме} \textbf{(Б)} пред комисията как само за няколко кратки \underline{месеци} \textbf{(В)} ще предложим съвършено нова организация на \underline{работния} \textbf{(Г)} процес.

\vspace{0.6cm}

\noindent
\textbf{11. В кое от изреченията е допусната пунктуационна грешка?} \\
\textbf{А)} Учителят по биология изглеждаше като най-обикновен и безобиден французин на Земята. \\
\textbf{Б)} По набраздените бузи на мъжа се стичаха големи кръгли сълзи и капаха върху дланите му, но той не ги забелязваше. \\
\textbf{В)} Всеки ден се разхождахме, а после се разделяхме до малката църква където лятно време търговците на плодове изнасяха сергии. \\
\textbf{Г)} Стояхме на снега на стълбите на университета и държахме шарени табели и плакати, подскачайки от студ.

\vspace{0.6cm}

\noindent
\textbf{12. В кое изречение НЕ е допусната пунктуационна грешка?} \\
\textbf{А)} В новия интернет сайт за образованието, в ЕС може да се намери всякаква информация, свързана с изграждането на европейското образователно пространство. \\
\textbf{Б)} Иновацията се превръща в нещо, което носи полза, когато е свързана с даване на конкретни отговори, свързани с области като: климат, енергетика, земеделие. \\
\textbf{В)} Пред комисията ще изложиш своето виждане, едва когато председателят запознае със становището на експертите. \\
\textbf{Г)} Момичето беше разбълнувано и ръкомахаше нервно без да може да контролира движенията си.

\vspace{0.6cm}

\noindent
\textbf{13. Срещу съответната буква запишете онзи от вариантите, поставени в скоби, който е правилен.} \\
\textbf{А)} (Парфюмерийно-козметичната / Парфюмерийнокозметичната) промишленост е свързана с използването на естествени суровини за производство на голяма гама продукти. \\

\vspace{-0.3cm}

\noindent
\textbf{Б)} Български екипажи грабнаха победата в двете най-големи и важни категории на едноседмичния (рали марафон / рали-маратон).

\vspace{0.3cm}
\noindent\textbf{А)} \dotfill \\
\noindent\textbf{Б)} \dotfill

\vspace{0.6cm}

\noindent
\textbf{14. Напишете изречението, като редактирате допуснатата правописна грешка.} \\
Метеоролозите съобщават, че през Ноември ни очакват по-високи от обичайните температури, но през третата десетдневка ще настъпи истинска зима.

\vspace{0.2cm}
\noindent\rule{\linewidth}{0.4pt}
\noindent\rule{\linewidth}{0.4pt}
\noindent\rule{\linewidth}{0.4pt}
\noindent\rule{\linewidth}{0.4pt}


\vspace{0.6cm}

\noindent
\textbf{15. Запишете формата, която е употребена неправилно в изречението, и формата, с която трябва да се замени, за да бъде поправена допуснатата граматична грешка.} \\

\vspace{-0.2cm}
\noindent
Днес трябва да съобщя на некому за възникналата повреда в инсталацията.
\vspace{0.2cm}

\vspace{0.2cm}
\noindent
\begin{tabularx}{\linewidth}{|>{\centering\arraybackslash}X|>{\centering\arraybackslash}X|}
\hline
\textbf{Неправилна форма} & \textbf{Правилна форма} \\
\hline
 &  \\
\hline
\end{tabularx}

\vspace{0.6cm}

\noindent
\textbf{16. Запишете правилните форми на думите, поставени в скоби.} \\

\vspace{-0.2cm}
\noindent
На борда на двата космически (кораба / кораби) са трима руски (космонавти / космонавта) и американски астронавт, които ще извършат подготвителни дейности, за да се освободи (лабораторният / лабораторния) кораб „Наука“.

\vspace{0.2cm}
\noindent\rule{\linewidth}{0.4pt}
\noindent\rule{\linewidth}{0.4pt}
\noindent\rule{\linewidth}{0.4pt}
\noindent\rule{\linewidth}{0.4pt}
\noindent\rule{\linewidth}{0.4pt}
\noindent\rule{\linewidth}{0.4pt}
\noindent\rule{\linewidth}{0.4pt}
\noindent\rule{\linewidth}{0.4pt}

\vspace{0.6cm}

\noindent
\textbf{17. Напишете текста, като поставите пропуснатите ЧЕТИРИ препинателни знака.} \\
Изискванията към качеството на медийното позициониране нарастват с всеки изминал ден. И брандовете знаят добре кои нагласи на потребителите могат да им бъдат от полза и какво трябва да променят в своята стратегия. Затова те разполагат със специален екип от професионалисти посветени на ярките и креативни решения. \\
Според Стоян Стоянов, директор Продажби в телевизионната компания размерът на екрана и факторът като високотехнологично обезпечаване на продукциите продължителност и сила на присъствието, вплитане на изразни средства и въздействие върху сетивата имат пряко отношение към успеха на дадена реклама.

\vspace{0.2cm}
\noindent\rule{\linewidth}{0.4pt}
\noindent\rule{\linewidth}{0.4pt}
\noindent\rule{\linewidth}{0.4pt}
\noindent\rule{\linewidth}{0.4pt}
\noindent\rule{\linewidth}{0.4pt}
\noindent\rule{\linewidth}{0.4pt}
\noindent\rule{\linewidth}{0.4pt}
\noindent\rule{\linewidth}{0.4pt}
\noindent\rule{\linewidth}{0.4pt}
\noindent\rule{\linewidth}{0.4pt}
\noindent\rule{\linewidth}{0.4pt}
\noindent\rule{\linewidth}{0.4pt}
\noindent\rule{\linewidth}{0.4pt}
\noindent\rule{\linewidth}{0.4pt}
\noindent\rule{\linewidth}{0.4pt}
\noindent\rule{\linewidth}{0.4pt}
\noindent\rule{\linewidth}{0.4pt}
\noindent\rule{\linewidth}{0.4pt}

\newpage

\subsection{Езиков тест XII клас 18.11.2024 г.}

\noindent
\textit{Име и фамилия:} \dotfill \textit{№} \makebox[2cm]{\dotfill}

\vspace{0.6cm}

\begin{center}
\noindent
\textbf{Езиков тест XII клас 18.11.2024 г.}
\end{center}

\vspace{0.6cm}

\noindent
\textbf{1. В кой ред всички думи НЕ са написани вярно?} \\
\begin{tabular}{@{}p{0.48\textwidth} p{0.48\textwidth}@{}}
\textbf{А)} компютри, мюзикъли, катетри & \textbf{В)} лакейство, героичен, крайща \\
\textbf{Б)} кървясал, млъквам, кълвач & \textbf{Г)} ограждам, подражател, упокой
\end{tabular}

\vspace{0.6cm}

\noindent
\textbf{2. В кое изречение е допусната правописна грешка?} \\
\textbf{А)} Художникът Иван Димитров откри самостоятелна изложба в Градската художествена галерия във Варна. \\
\textbf{Б)} Този път актьорът влиза в ролята на вдъхновител, който разпалва представата ни за това, което е възможно да постигнем. \\
\textbf{В)} Пикът на Лунното затъмнение ще бъде в 11:05 часа местно време и ще бъде видимо от Северна и Южна Америка. \\
\textbf{Г)} Затъмненията в зодиакалните земни знаци обикновено предвещават недостиг на пшеница, царевица и много селскостопански продукти.

\vspace{0.6cm}

\noindent
\textbf{3. В кое изречение НЕ е допусната правописна грешка?} \\
\textbf{А)} Днес за небостъргач обикновенно се считат сгради с височина над 150 метра, но през 19. век няма официална дефиниция. \\
\textbf{Б)} В края на 1800 г. индустриялизацията увеличава населението в града и цените на земята, което прави високите сгради все по-рентабилни. \\
\textbf{В)} Строителството на високи сгради в Азия може да се дължи от части на по-ниските строителни разходи в тези градове. \\
\textbf{Г)} Намаляването на въглеродните емисии е новото предизвикателство, пред което са изправени архитектите днес.

\vspace{0.6cm}

\noindent
\textbf{4. Препишете от изречението думите, които са с правописна грешка, като ги редактирате.} \\
Поведението на слоновете се свързва с уникална животинска интелигентност, която показва мъка, алтруизъм, състрадание, самосъзнание, игра, изкуство и музика. \\
\vspace{0.2cm}
\noindent\rule{\linewidth}{0.4pt}
\noindent\rule{\linewidth}{0.4pt}

\vspace{0.6cm}

\noindent
\textbf{5. В кое изречение е допусната пунктуационна грешка?} \\
\textbf{А)} Африканските слонове имат големи уши, оформени като континента Африка. \\
\textbf{Б)} Прасето или домашната свиня е бозайник от разред чифтокопитни. \\
\textbf{В)} Младите лъвове, прокудени от групата, често изминават големи разстояния в търсене на други прокудени мъжки. \\
\textbf{Г)} Разбира се ловът се извършва предимно от здрач до зори и през по-хладните часове на деня.

\vspace{0.6cm}

\noindent
\textbf{6. В кое изречение НЕ е допусната пунктуационна грешка?} \\
\textbf{А)} Известно е, че зайчетата изострят зъбите си върху електрически проводници, стайни растения и мебели. \\
\textbf{Б)} Трябва да къпете зайчето си само, когато има причина и ви е препоръчано от ветеринарен лекар. \\
\textbf{В)} Учените са доказали, че кучетата имат усещане за време, и знаят разликата между час и пет минути. \\
\textbf{Г)} Проучване твърди, че кучето ви може истински да ревнува когато вижда, че проявявате обич към друго същество.

\vspace{0.6cm}

\noindent
\textbf{7. В кое изречение е допусната пунктуационна грешка?} \\
\textbf{А)} От една страна, обичам предизвикателствата, но от друга страна те ме плашат. \\
\textbf{Б)} Не можехме да разберем кой беше първият, повдигнал въпроса за парите. \\
\textbf{В)} Виж сега, Ирина, аз много добре разбирам проблема на групата. \\
\textbf{Г)} Ей, колко се ядосах заради пропуснатия мач на брат ми.

\vspace{0.6cm}

\noindent
\textbf{8. Препишете изречението, като поставите пропуснатите пунктуационни знаци.} \\
Благородник далечен роднина на автора на романа Война и мир Лев Толстой Алексей Толстой редува различни жанрове по време на кариерата си но се специализира в научната фантастика и историческите романи.

\vspace{0.2cm}
\noindent\rule{\linewidth}{0.4pt}
\noindent\rule{\linewidth}{0.4pt}
\noindent\rule{\linewidth}{0.4pt}
\noindent\rule{\linewidth}{0.4pt}
\noindent\rule{\linewidth}{0.4pt}
\noindent\rule{\linewidth}{0.4pt}
\noindent\rule{\linewidth}{0.4pt}
\noindent\rule{\linewidth}{0.4pt}

\vspace{0.6cm}

\noindent
\textbf{9. В кое изречение НЕ е допусната граматична грешка?} \\
\textbf{А)} Рисуването от натура е единственият начин да се опознаят формата, обемът, перспективата и цветовата композиция. \\
\textbf{Б)} Пресъздаването на действителни образи посредством методиките на изобразително изкуство sa перфектния начин да се тренира ръката. \\
\textbf{В)} Изграждайки обем и светлосянка, трябва да се изобразят няколко основни елемента, като колонна сянка и светлина. \\
\textbf{Г)} Проучването и претворяването на форми от природата и изкуството се основават на знания за формите, пространството и светлосянката.

\vspace{0.6cm}

\noindent
\textbf{10. В кое изречение е допусната граматична грешка?} \\
\textbf{А)} Предлагаме днес да говоримe за добродетелите на българина през вековете. \\
\textbf{Б)} Домакинята предложи студен чай на гостите, предполагайки, че са ожаднели от пътуването. \\
\textbf{В)} Днес се видях и се запознах с всички новопостъпили във фирмата през последния месец. \\
\textbf{Г)} Говорихме за пътя на четата и нейната неподготвеност за сблъсъка в добре въоръжения противник.

\vspace{0.6cm}

\noindent
\textbf{11. Запишете правилните бройни форми на думите, поставени в скоби.} \\
За летуването в планината взех чифт (кец), два топли (панталон) и няколко (пуловер).

\vspace{0.2cm}
\noindent\rule{\linewidth}{0.4pt}

\vspace{0.6cm}

\noindent
\textbf{12. Обяснете със свои думи значението на фразеологичното словосъчетание \textit{ходя по нервите}.}

\vspace{0.2cm}
\noindent\rule{\linewidth}{0.4pt}
\noindent\rule{\linewidth}{0.4pt}

\vspace{0.4cm}

\noindent
\textbf{13. В кое изречение НЕ е допусната лексикална грешка?} \\
\textbf{А)} На празника плаках и леех сълзи от умиление. \\
\textbf{Б)} Проявете малко повече емоции и чувства към това дете. \\
\textbf{В)} Развълнувано разказа какво видя на площада през галерията. \\
\textbf{Г)} Влязохме забързано у сградата с високите вити стълби.

\vspace{0.6cm}

\noindent
\textbf{14. Коя от думите е излишна?} \\
\begin{tabular}{@{}p{0.24\textwidth} p{0.24\textwidth} p{0.24\textwidth} p{0.24\textwidth}@{}}
\textbf{А)} смерч & \textbf{Б)} ураган & \textbf{В)} пожар & \textbf{Г)} буря
\end{tabular}

\vspace{0.6cm}

\noindent
\textbf{Прочетете текста и отговорете на въпроси от 15. до 24. включително.}

\vspace{0.3cm}

\noindent
Никола Фичев се счита за една от най-ярките ренесансови фигури на българското възраждане. От своите съвременници е бил ценен като майстор (уста), а от наследниците си – като признат първомайстор на националната архитектура. Много са мостовете, къщите, църквите, камбанариите и чешмите, които е сътворил с ръцете си и вложил частица от себе си в тях. Ето част от неговите архитектурни шедьоври:

Храм „Св. Троица“ – съграден през 1867 г., е един от най-интересните архитектурни паметници в Свищов. Иконите са рисувани от свищовлията Николай Павлович. Храмът е пострадал много при земетресението от 1977 г., но през 1992 г. реставрацията му е завършена.

Мостът над река Росица до Севлиево. По време на Кримската война (1853–1856) движението на войска, пътници и стока се засилило. На местната власт било възложено със свои сили да организира построяването на каменен мост. За главен майстор на строежа на моста бил ангажиран майстор Никола Фичев от Дряново, прочут български строител и архитект. Приемането на сравнително голямата поръчка било изпитание за възможностите на майстора.

Строеният по-късно мост при Бяла е много по-внушителен, но той бил изграден с предоставени от централната власт средства след специален закон от 1864 г. В Севлиево такива средства нямало, но през 1856 г. строежът започнал. Работата вървяла сравнително бавно заради липсата на пари, но следващите една-две години грубият строеж бил приключен. Мостът бил дълъг 110 метра, поддържали го седем свода, като централният бил най-висок и оформятъл „гърбица“.

Съоръжението е от ранните хидротехнически постройки на майстора, в което развива своите идеи за традиционния възрожденски тип „гърбат“ мост. Най-висок е вторият свод, над който мостовото платно е леко изгърбено. През XX век гърбицата е премахната и мостът е удължен с насипи от двата му края.

Каменният мост при Бяла е забележителна творба от Възрожденската епоха. Носещите стълбове са украсени със символни скулптури: грифон, лъв, нимфа, лебед. Мостът е истинско художествено произведение. При него се наблюдава уникално съчетание между архитектура и скулптура. До този момент подобни орнаменти са използвани само за украса на сгради. Неравната зидария на моста е облицована с пълтен, добре обработен варовик, устоите имат вълнообразен завършек.

Всеки устой наподобява фасада на сграда, пропускателните отвори са оформени като отворени и затворени прозорци. За изграждането на моста помагали ковачи от Габрово, дърводелци и каменари от Трявна и Дряново, селяни от Бяла. Дължината му е 276 м, а ширината 9,5 м. Има 13 подпорни стълба и 14 полукръгли свода. В средата се издига колона, на която е поставена мраморна плоча с арабски надпис. \\
\textit{(https://opoznal.bg/articles/view/356)}

\vspace{0.6cm}

\noindent
\textbf{15. Кое твърдение е вярно за текста?} \\
\textbf{А)} Предоставя информация за майстор Кольо Фичето, предназначена за тесни специалисти. \\
\textbf{Б)} Дава информация за някои от строежите на майстор Кольо Фичето, достъпна за всеки. \\
\textbf{В)} Художествено изобразява заслугите на Кольо Фичето за българското строителство. \\
\textbf{Г)} Показва напредничавото мислене и ограниченията пред строителството през Възраждането.

\vspace{0.6cm}

\noindent
\textbf{16. Кое твърдение НЕ е вярно според текста?} \\
\textbf{А)} Мостът над река Росица е изграден със средства, предоставени от централната власт. \\
\textbf{Б)} Храмът в Свищов е изграден от Кольо Фичето, но иконите в него са дело на друг. \\
\textbf{В)} Заради Кримската война местната власт в Севлиево трябва да построи мост над реката. \\
\textbf{Г)} Мостът при Бяла е от ранните хидротехнически постройки на Кольо Фичето.

\vspace{0.6cm}

\noindent
\textbf{17. Със свои думи обяснете какво означава „гърбат“ мост според текста.}

\vspace{0.2cm}
\noindent\rule{\linewidth}{0.4pt}

\vspace{0.6cm}

\noindent
\textbf{18. Защо майстор Кольо Фичето е определян като ренесансова фигура?}

\vspace{0.2cm}
\noindent\rule{\linewidth}{0.4pt}

\vspace{0.6cm}

\noindent
\textbf{19. Посочете две особености, които превръщат моста при Бяла в художествено произведение.}
\vspace{0.2cm}

\noindent\rule{\linewidth}{0.4pt}
\noindent\rule{\linewidth}{0.4pt}

\vspace{0.6cm}

\noindent
\textbf{20. Обяснете в 1 – 2 изречения с какво мостът при Бяла се отличава от моста при Севлиево.}

\vspace{0.2cm}
\noindent\rule{\linewidth}{0.4pt}
\noindent\rule{\linewidth}{0.4pt}
\noindent\rule{\linewidth}{0.4pt}
\noindent\rule{\linewidth}{0.4pt}

\vspace{0.6cm}

\noindent
\textbf{21. Напишете кратка теза (4 – 5 изречения) на тема: \textit{Българските поводи за гордост}.}

\vspace{0.2cm}
\noindent\rule{\linewidth}{0.4pt}
\noindent\rule{\linewidth}{0.4pt}
\noindent\rule{\linewidth}{0.4pt}
\noindent\rule{\linewidth}{0.4pt}
\noindent\rule{\linewidth}{0.4pt}
\noindent\rule{\linewidth}{0.4pt}
\noindent\rule{\linewidth}{0.4pt}
\noindent\rule{\linewidth}{0.4pt}
\noindent\rule{\linewidth}{0.4pt}
\noindent\rule{\linewidth}{0.4pt}

\vspace{0.6cm}

\noindent
\textbf{22. Кой израз от текста е термин?} \\
\begin{tabular}{@{}p{0.48\textwidth} p{0.48\textwidth}@{}}
\textbf{А)} обработен варовик & \textbf{В)} хидротехнически постройки \\
\textbf{Б)} всепризнат майстор & \textbf{Г)} главен майстор на строежа
\end{tabular}

\vspace{0.6cm}

\noindent
\textbf{23. Каква е връзката между думите майстор и уста?} \\
\begin{tabular}{@{}p{0.48\textwidth} p{0.48\textwidth}@{}}
\textbf{А)} И двете са разговорна лексика. & \textbf{В)} Синоними са. \\
\textbf{Б)} Имат еднакъв произход. & \textbf{Г)} Сродни думи са.
\end{tabular}

\vspace{0.6cm}

\noindent
\textbf{24. Какво е значението на думата свищовлия?}

\vspace{0.2cm}
\noindent\rule{\linewidth}{0.4pt}

\vspace{0.6cm}

\noindent
\textbf{25. Прочетете изреченията и открийте думата, в която е допусната лексикална грешка. Запишете САМО редактирания вариант.} \\

\vspace{-0.2cm}
\noindent
БНБ в навечерието на 3-ти март пусна в обръщение златна монета с изображение на връх Шипка.

\vspace{0.2cm}
\noindent\rule{\linewidth}{0.4pt}

\vspace{0.2cm}
\noindent
Тя отправи към него безпристрастен поглед, който го закова на място.

\vspace{0.2cm}
\noindent\rule{\linewidth}{0.4pt}

\vspace{0.2cm}
\noindent
Не знам защо на този човек му преписват всевъзможни грехове.

\vspace{0.2cm}
\noindent\rule{\linewidth}{0.4pt}

\vspace{0.6cm}

\noindent
\textbf{26. Прочетете изреченията и открийте фразеологичните словосъчетания в тях. Запишете САМО тяхното значение.} \\

\vspace{-0.2cm}
\noindent
Създаде се конфликтна ситуация с колегите, от която излязох сух от водата.

\vspace{0.2cm}
\noindent\rule{\linewidth}{0.4pt}

\vspace{0.2cm}
\noindent
В края на учебната година ми предстои да премина Рубикон.

\vspace{0.2cm}
\noindent\rule{\linewidth}{0.4pt}

\vspace{0.2cm}
\noindent
Съзнавах, че съм малодушен, свивах се в черупката си и чаках по-добри дни.

\vspace{0.2cm}
\noindent\rule{\linewidth}{0.4pt}

\vspace{0.6cm}

\noindent
\textbf{27. Прочетете текста и поставете пропуснатите пунктуационни знаци.} \\
Заедно с френската войска в Египет навлезли и много цивилни учени историци, филолози, архитекти, математици ботаници художници. Бонапарт съвестно се грижел за тях почти колкото за своя обоз и преди битка никога не забравял да каже Магаретата и учените – в средата. И мъжете които войниците започнали да наричат магарета и учени по-късно се нахвърляли на египетските старини опаковали артефактите в сандъци, измервали, копирали правели гипсови отливки, рисували и чертаели.

\vspace{0.2cm}
\noindent\rule{\linewidth}{0.4pt}
\noindent\rule{\linewidth}{0.4pt}
\noindent\rule{\linewidth}{0.4pt}
\noindent\rule{\linewidth}{0.4pt}
\noindent\rule{\linewidth}{0.4pt}
\noindent\rule{\linewidth}{0.4pt}
\noindent\rule{\linewidth}{0.4pt}
\noindent\rule{\linewidth}{0.4pt}
\noindent\rule{\linewidth}{0.4pt}
\noindent\rule{\linewidth}{0.4pt}
\noindent\rule{\linewidth}{0.4pt}

\vspace{0.6cm}

\noindent
\textbf{28. Запишете правилната форма на думата, поставена в скоби.} \\
\noindent
\textbf{А) членуване} \\
Не очакваха, че той ще бъде (първи) \dotfill , който ще се опълчи на началника за преувеличените данни в доклада. \\
В реалити шоу, излъчено по \\(Българска национална телевизия) \dotfill  изгря звезда на нов музикален талант.

\vspace{0.2cm}
\noindent
\textbf{Б) местоимения} \\
Препрочетох книгата, в (чиито/чийто) \dotfill  страници бе скрит кодът към криптата. \\
Непознатите се скриха в тълпата, без да подозират,\\ че (тя/те) \dotfill  ги понасят точно към патрула. \\
В гаража рано сутринта ръмжеше моторът на момичето.\\ Съседите не можеха да (я/го) \dotfill  търпят повече.

\vspace{0.2cm}
\noindent
\textbf{В) учтива форма} \\
Господин Любомиров, бихте ли (споделил) \dotfill  наблюденията\\ си върху поведението на пчелите през зимата, за което вече\\ сте (бил) \dotfill  (канен) \dotfill да говорите пред сдружението на пчеларите у нас.

\newpage
\subsection{ТЕСТ № 2}

\noindent
\textit{Име и фамилия:} \dotfill \textit{№} \makebox[2cm]{\dotfill}

\vspace{0.6cm}

\noindent
\begin{center}
  \textbf{ТЕСТ № 2}   
\end{center}

\vspace{0.2cm}
\noindent
\textbf{ПЪРВА ЧАСТ (60 минути)}

\vspace{0.4cm}

\noindent
\textbf{1. В кой ред думата е изписана правилно?} \\
\begin{tabular}{@{}p{0.24\textwidth} p{0.24\textwidth} p{0.24\textwidth} p{0.24\textwidth}@{}}
\textbf{А)} стояли & \textbf{Б)} свитaк & \textbf{В)} предверие & \textbf{Г)} смегчавам
\end{tabular}

\vspace{0.6cm}

\noindent
\textbf{2. В кое от изреченията има дума с правописна грешка?} \\
\textbf{А)} Най-добре запазеният архитектурен паметник от римската епоха в София е ротондата „Св. Георги“. \\
\textbf{Б)} Възнамерявам още утре да се обърна към председатела на съда за съдействие по заведеното дело. \\
\textbf{В)} Когато имате запушен слъзен канал, сълзите не могат да се оттичат нормално от раздразненото око. \\
\textbf{Г)} До края на седмицата ще получите поръчаните два термометъра – един живачен и един дигитален.

\vspace{0.6cm}

\noindent
\textbf{3. В кое изречение НЕ е допусната правописна грешка?} \\
\textbf{А)} Гамалъчите са форма на електромагнитно излъчване с много малка дължина на вълната. \\
\textbf{Б)} В Иван-Вазовата повест Странджата и Македонски са представени като герои и мъченици. \\
\textbf{В)} Даначната политика е сърцевината на икономическата и фискалната политика на държавата. \\
\textbf{Г)} Нараства търсенето на по-малки жилища в големите градове с цел оддаването им под наем.

\vspace{0.6cm}

\noindent
\textbf{4. В кое изречение е допусната граматична грешка?} \\
\textbf{А)} Младите хора вече приемат образованието като изключително важно за тяхното бъдеще. \\
\textbf{Б)} Можете да намерите новите колекции дамски, мъжки и детски дрехи в онлайн магазина. \\
\textbf{В)} Иванов, учителят ни по математика, този път бе много доволен от резултатите на теста. \\
\textbf{Г)} Ще получите информация по отношение на частната и публичната държавна собственост.

\vspace{0.6cm}

\noindent
\textbf{5. В кое изречение НЕ е допусната граматична грешка?} \\
\textbf{А)} В кутията има три вида ръчно изработени бонбона от черен, млечен и бял шоколад. \\
\textbf{Б)} Един от уроците, който научих в живота си, е да проявявам толерантност към хората. \\
\textbf{В)} Българската национална телевизия ще излъчи всички срещи от тазгодишния шампионат. \\
\textbf{Г)} Договорът за двустранно сътрудничество участниците подписаха по време на конгреса.

\vspace{0.6cm}

\noindent
\textbf{6. В кое изречение е допусната пунктуационна грешка?} \\
\textbf{А)} Бащата на Светла беше тих, замислен и дълбоко вглъбен в себе си човек. \\
\textbf{Б)} Детското отделение разполага с 47 нови легла, разпределени в два сектора. \\
\textbf{В)} Независимо от многото заявки, мили колеги надявам се, че ще се справим. \\
\textbf{Г)} Цяла чиния с бисквитки, разбира се, е за предпочитане пред една бисквитка.

\vspace{0.6cm}

\noindent
\textbf{7. В коя от позициите е допусната пунктуационна грешка?} \\
Построената от Хрельо Драголов средновековна църква, \textbf{(А)} все пак е оцеляла до 1834 г., \textbf{(Б)} когато по решение на манастирското братство е разрушена, \textbf{(В)} като върху част от нейните основи е построена днешната съборна църква, \textbf{(Г)} но въпреки това тя остава в паметта на местните като Хрельовата църква.

\vspace{0.6cm}

\noindent
\textbf{8. Запишете правилната форма на думата, поставена в скоби.} \\
Ще имате възможност да се запознаете с историите на създаването на едни от най-красивите (ария).

\vspace{0.2cm}
\noindent\rule{\linewidth}{0.4pt}

\vspace{0.4cm}

\noindent
\textbf{9. Запишете правилната форма САМО на думата, в която е допусната правописна грешка.} \\
Клубът по алпинизъм възнамерява да организира няколко изкачвания на върхове в и извън България.

\vspace{0.2cm}
\noindent\rule{\linewidth}{0.4pt}

\vspace{0.6cm}

\noindent
\textbf{10. Запишете правилната граматична форма на думата, поставена в скоби.} \\
Геологическият институт на (Българска) академия на науките ще отбележи с тържество своята 75-годишнина.

\vspace{0.2cm}
\noindent\rule{\linewidth}{0.4pt}

\vspace{0.6cm}

\noindent
\textbf{11. Запишете правилната форма САМО на думата, в която е допусната граматична грешка.} \\
Основният проблем, който стои пред нас и който трябва в най-скоро време да решим, е как възнамеряваме да се справиме с предизвикателствата в дигиталната среда.

\vspace{0.2cm}
\noindent\rule{\linewidth}{0.4pt}

\vspace{0.6cm}

\noindent
\textbf{12. Запишете правилната форма САМО на думата, в която е допусната граматична грешка.} \\
Господине, Вие сте одобрен за участие в конкурса, но бихте ли попълнил още два формуляра, за да ги добавим към документите?

\vspace{0.2cm}
\noindent\rule{\linewidth}{0.4pt}

\vspace{0.6cm}

\noindent
\textbf{13. В текста са пропуснати САМО ПЕТ препинателни знака. Препишете ЦЕЛИЯ текст в свитка за свободните отговори, като поставите пропуснатите знаци.} \\
Да бъдат независими е доминиращият стремеж на тийнейджърите. Важно е да запомните думите на един известен поет Децата ви не са ваши деца. Те са синове и дъщери на копнежа за живота сам по себе си. Те идват чрез вас но не от вас, и въпреки че са с вас те не ви принадлежат.

\vspace{0.6cm}

\noindent
\textbf{\textit{Прочетете двата текста и изпълнете задачите към тях (от 14. до 21. включително).}}

\vspace{0.3cm}
\noindent\textbf{Текст 1} \\
На 12 януари 1915 г. в елитното софийско кино „Модерен театър“ се е състояла премиерата на първия български филм „Българан е галант“. „Сензация в „Модерен театър“!“ – няколко дни тръбят всекидневниците „Пряпорец“, „Дневник“ и „Мир“. Краттките реклами информират, че „комедията е заснета в София и е играна от българските артисти при Народния театър т-жа Мара Липина и т-н Васил Гендов“. 

Още докато учи във Виена (1908 – 1909), Гендов решава да създаде кинотворба. Но за целта са му нужни доста пари. Затова по-късно, в София, режисьорът се обръща за помощ към няколко търговци, индустриалци и богаташи, но получава отказ. Вероятно през 1914 г. напористият младеж се среща с предприемчивия унгарец Аладар Оттай – съсобственик на „Модерен театър“, който веднага се запалва по идеята и предоставя на Гендов снимачна техника, пари и своя оператор.

Историята във филма наподобява тогавашните чуждестранни градски комедии. Елегантен франт, наречен Българан, среща на улицата млада дама и започва да я ухажва. Тя решава да даде урок на натрапника и му предлага да я придружи по магазините, където пазарува с размах. По някое време даматя открива, че е забравила парите си вкъщи, и моли ухажора да й даде назаем. На път за дома и срещат нейния съпруг, който освобождава „носача“ на покупкитя и му се отблагодарява щедро с... 50 стотинки. Сюжетът предполага да се снима по елегантните и луксозни места в тогавашната София – пред къщата на богатия търговец Димитър Съселов на улица „Славянска“ № 19, пред красивите витрини на скъпите магазини на централната „Леге“ и край шикозните заведения, а героите да са облечени по последна мода.

Голяма трудност създава подборът на актьори. Театралните звезди отказват да се снимат заради лошото мнение за киното по онова време. Докато театърът вече е станал част от елитарните забавления на бомонда, „живата фотография“ – представяна от скитащи „палаткаджии“ – се смята за простонародно зрелище. Бюджетът на „Българан“ е нищожен. Поради тази причина проблем има и с реквизита – парите едва покриват наема за дрехи „от частни лица за временно ползване“. Заради лошия имидж на кинаджиите никой не желае да заеме дрехите за подобни „маймунджилъци“. Любопитен е случай с собственика на шапка. Режисьорът убеждава столичанката, че шапката е нужна за... благотворително представление за Червения кръст. След началото на снимките обаче даматя идва и си я иска обратно. Разбрала за какво са я ползвали, в пристъп на гняв тя се разкрещява: „Как ще сложа аз утре тази шапка, когато така я омаскарихте?!“ Кинографът е налице – героинята влиза в луксозен магазин с една шапка, а на излизане вече е с друга. Е, може да си я е купила вътре...

Външните снимки на „Българан“ разгневяват и собствениците на къщи и магазини. Един от тях заплашва екипа с полиция, ако табелата на дюкяна му е попадала в кадър – това би било лоша реклама за търговията. А днес компаниите хвърлят огромни пари за продуктово позициониране, и то за няколко секунди на екран... 

Васил Гендов споделя в спомените си, че „Модерен театър“ не отделя нито стотинка от приходите за екипа, създал филма. Режисьорът продава своя часовник, за да услужи с пари на актрисата от Пловдивския театър Мара Липина, изпълнила женска роля, когато престоят й в София се удължава и тя остава без средства. Въпреки това на режисьора е заплатно... „не възторзите на публиката, което за начало все пак е най-голямата награда“. После „Българан“ потъва в тайнственост и до днес не се знае къде се намира първият български нерален филм.

\vspace{0.6cm}

\noindent
\textbf{Текст 2} \\
\textbf{Закон за филмовата индустрия} \\
\textbf{Глава първа}

\vspace{0.2cm}
\noindent
\textbf{ОБЩИ ПОЛОЖЕНИЯ} \\
\textbf{Чл. 1.} Този закон урежда отношенията, свързани с производството, разпространението, промоцията и показа на филмите в Република България и държавното подпомагане на българската филмова индустрия, като създава условия за нейното развитие.

\noindent
\textbf{Чл. 2. (1)} Филмова индустрия по смисъла на закона е производството, разпространението, промоцията, показът и съхранението на филми. \\
\textbf{(2)} Дейността по съхранението на филми се осъществява от Българската национална филмотека като държавен културен институт с национално значение.

\noindent
\textbf{Чл. 3. (1)} Държавното подпомагане на филмовата индустрия цели: \\
1. утвърждаване на филмовото творчество като важна област на националната култура; \\
2. стимулиране на производството, разпространението, промоцията и показа на националната филмова продукция; \\
3. поощряване създаването и разпространението на високохудожествени произведения на филмовото творчество. \\
\textbf{(2)} Приоритети на държавната политика във филмовата индустрия са: \\
1. правото на обществен достъп до разнообразни форми на филмово творчество; \\
2. защита на правата и интересите на зрителите; \\
3. подкрепа на нови таланти и млади автори; \\
4. представяне на българското кино в страната и в чужбина.

\vspace{0.4cm}
\noindent
{\footnotesize{\textsuperscript{1}Бомонд – висше, отбрано общество.} \hfill}

\vspace{0.6cm}

\noindent
\textbf{14. Кое от следните твърдения е вярно според Текст 1?} \\
\textbf{А)} Появата на първия български филм е приета нееднозначно от медиите. \\
\textbf{Б)} На предприемача Аладар Оттай се налага да търси финансиране за филма. \\
\textbf{В)} Като декор във филма са използвани луксозни сгради в центъра на София. \\
\textbf{Г)} Собствениците на къщи и магазини са доволни от безплатната реклама.

\vspace{0.6cm}

\noindent
\textbf{15. Защо киното е наречено \textit{живата фотография}, а снимачният екип – \textit{палаткаджии}, в подчертаното изречение от Текст 1?} \\
\textbf{А)} За да се оразличат киното и снимачният екип от фотографията и фотографите. \\
\textbf{Б)} За да бъдат защитени киното и снимачният екип от несправедливите обвинения. \\
\textbf{В)} За да бъдат изтъкнати предимствата на несправедливо отричаното киноискуство. \\
\textbf{Г)} За да се подчертае непрестижността на киноизкуството в съпоставка с театъра.

\vspace{0.6cm}

\noindent
\textbf{16. Кой от проблемите, с които се сблъскват създателите на първия български филм, НЕ е засегнат в разпоредбите от Текст 2?} \\
\textbf{А)} подкрепата на кинопроизводството \\
\textbf{Б)} съобразяването с нагласите на зрителите \\
\textbf{В)} съхраняването на заснетата материал \\
\textbf{Г)} поощряването на младите творци

\vspace{0.6cm}

\noindent
\textbf{17. С какво значение е употребена думата \textit{промоцията} в Текст 2?} \\
\textbf{А)} даване, присъждане на определена научна степен на някого \\
\textbf{Б)} раздаване на дипломи на завършили студенти или докторанти \\
\textbf{В)} публично представяне на специална среща или чрез медиите \\
\textbf{Г)} предлагане на по-ниска цена или безплатно с рекламна цел

\vspace{0.6cm}

\noindent
\textbf{18. Като използвате информацията от Текст 2, запишете кой е най-важният критерий един български филм да бъде подпомогнат от държавата.}

\vspace{0.2cm}
\noindent\rule{\linewidth}{0.4pt}
\noindent\rule{\linewidth}{0.4pt}
\noindent\rule{\linewidth}{0.4pt}
\noindent\rule{\linewidth}{0.4pt}

\vspace{0.6cm}

\noindent
\textbf{19. Като използвате информацията от Текст 1, запишете два аргумента в подкрепа на твърдението, че мнението за киното в началото на 20. век е негативно.}

\vspace{0.2cm}
\noindent\rule{\linewidth}{0.4pt}
\noindent\rule{\linewidth}{0.4pt}

\vspace{0.4cm}

\noindent
\textbf{20. Като имате предвид информацията от Текст 1, коментирайте ролята на многоточието след подчертаните изречения:} \\
\textit{Кинографът е налице – героинята влиза в луксозен магазин с една шапка, а на излизане вече е с друга. Е, може да си я е купила вътре...}

\vspace{0.2cm}
\noindent\rule{\linewidth}{0.4pt}
\noindent\rule{\linewidth}{0.4pt}

\vspace{0.4cm}

\noindent
\textbf{21. В текст от 4 – 5 изречения коментирайте защо българското общество в началото на миналия век се отнася с такова недоверие към киноизкуството.}

\vspace{0.2cm}
\noindent\rule{\linewidth}{0.4pt}
\noindent\rule{\linewidth}{0.4pt}
\noindent\rule{\linewidth}{0.4pt}
\noindent\rule{\linewidth}{0.4pt}
\noindent\rule{\linewidth}{0.4pt}

\newpage

\vspace{0.2cm}
\noindent
\textbf{ВТОРА ЧАСТ (60 минути)}

\vspace{0.4cm}

\noindent
\textbf{22. Запишете САМО паронима, с който да поправите лексикалната грешка в изречението.}\\
В настоящото проучване се споделя мнението, че военният конфликт вероятно ще доведе до тежка икономическа криза в Европа.

\vspace{0.2cm}
\noindent\rule{\linewidth}{0.4pt}

\vspace{0.6cm}

\noindent
\textbf{23. За всяко празно място изберете УМЕСТНАТА дума и я запишете срещу съответната буква в листа за отговори.}\\
С развитието на технологиите се \underline{\hspace{3cm}} \textbf{(А)} и рисковете за сигурността, свързани с тяхното използване, както и потенциалът за злоупотреба с тях. Хакерите могат да използват силата на изкуствения интелект, за да разработват \underline{\hspace{3cm}} \textbf{(Б)} кибератаки, да заобикалят мерките за сигурност и да откриват \underline{\hspace{3cm}} \textbf{(В))} в системите.

\vspace{0.2cm}

\noindent
\textbf{А)} усилват, увеличават, усложняват\\
\textbf{Б)} по-усъвършенствани, по-несъвършени, по-осъществени\\
\textbf{В)} пробоини, иновации, подобрения\\
\vspace{0.2cm}

\noindent
\textbf{24. Кой конфликт НЕ е заложен в интерпретацията на темата за живота и смъртта в стихотворението „До моето първо либе“?}\\
\textbf{А)} робство -- свобода\\
\textbf{Б)} любов -- омраза\\
\textbf{В)} страх -- надежда\\
\textbf{Г)} минало -- настояще

\vspace{0.6cm}

\noindent
\textbf{25. Кой от проблемите е интерпретиран и в двата посочени откъса?}\\

\vspace{-0.3cm}

\noindent
\textit{Аз умирам и светло се раждам --\\
разнолика, нестройка душа,\\
през деня неуморно изграждам,\\
през нощта без пощада руша.
}
\vspace{0.3cm}
\\
\noindent
\textit{Призова ли дни светло-смирени,\\
гръмват бури над тъмно море,\\
а подиря ли буря -- край мене\\
всеки вопъл и ропот замре.} \\
\quad (Димчо Дебелянов, „Черна песен“)

\vspace{0.5cm}

\noindent
\textit{Аз не живея: аз горя. Непримирими\\
в гърдите ми се борят две души:\\
душата на ангел и демон. В гърди ми\\
те пламъци дишат и плам ме суши.} \\
\quad (Пейо Яворов, „Две души“)

\vspace{0.4cm}

\noindent
\textbf{А)} за непреодолимото раздвоение в човешката душа\\
\noindent
\textbf{Б)} за несправедливостите в човешкото битие\\
\noindent
\textbf{В)} за неизбежния край на човешкото щастие\\
\noindent
\textbf{Г)} за безвъзвратно загубения национален идеал

\vspace{0.6cm}

\noindent
\textbf{26. Какво увековечава стихотворението „Новото гробище над Сливница“?}
\\
\noindent
\textbf{А)} Смъртта, която доказва, че националният идеал е жив.
\\
\noindent
\textbf{Б)} Смъртта, която преосмисля съдбата на враговете.
\\
\noindent
\textbf{В)} Смъртта, която е следствие от безсмислена саможертва.
\\
\noindent
\textbf{Г)} Смъртта, която поражда любов към другия етнос.

\vspace{0.6cm}

\noindent
\textbf{27. Коя композиционна особеност НЕ се открива в „Балада за Георг Хених“?}
\\
\noindent
\textbf{А)} емоционални обръщения на разказвача към майстора лютиер
\\
\noindent
\textbf{Б)} композиционна рамка, която представя настоящето на героя разказвач
\\
\noindent
\textbf{В)} повествовател, който е и един от героите в творбата
\\
\noindent
\textbf{Г)} подреждане на епизодите в естествената им времева последователност

\vspace{0.6cm}

\noindent
\textbf{28. Кое твърдение за думите на героинята от стихотворението „Посвещение“ е вярно?}

\vspace{0.2cm}

\noindent
\textit{Ела! Ще измием луната от грях!\\
Ще хвърлим трупа на умрелия страх...}

\vspace{0.3cm}

\noindent
\textbf{А)} Поставят акцент върху невъзможността да живее сама в трудните дни.
\\
\noindent
\textbf{Б)} Разкриват очакваното духовно съзряване и жадуваното себепостигане.
\\
\noindent
\textbf{В)} Подсказват наивността и желанието й да победи неизбежната смърт.
\\
\noindent
\textbf{Г)} Загатват нейния нелек избор да поеме по пътя на осъзнатия грях.

\vspace{0.6cm}

\noindent
\textbf{29. Кое от твърденията за цитирания откъс от „История“ е НЕВЯРНО?}

\vspace{0.2cm}

\noindent
\textit{Ще хванеш контурите само,\\
а вътре, знам, ще бъде празно\\
и няма никой да разказва\\
за простата човешка драма.}

\vspace{0.3cm}

\noindent
\textbf{А)} Насочва към размисъл за съдбата на онези, които водят борба за оцеляване.
\\
\noindent
\textbf{Б)} Разкрива незаслужената забрава на онеправданите, които воюват с живота.
\\
\noindent
\textbf{В)} Представя страданието от настъпилия разрив между историята и човека.
\\
\noindent
\textbf{Г)} Внушава идеята, че историята е неспособна да разкаже болката на човека.

\vspace{0.6cm}

\noindent
\textbf{30. Коя тема е интерпретирана в „Хлътни от мен!“ от Яна Язова?}

\vspace{0.2cm}

\noindent
\textit{Хлътни от мен, сладка песен,\\
при него и запей му ти!\\
Ако по онзи край е есен,\\
омай го, да се не смути!\\
\\
Прегръща ли го буйна пролет --\\
звъни на неговия пир!\\
На мен не трябва твоя полет --\\
аз имам тих самотен мир.\\
\\
Идете, слънце, извор, рози!\\
Правете му живота скъп!\\
Попътен вятър да го вози,\\
кога по мен го пари скръб.\\
\\
Ти, месец, дай му упоение,\\
блаженство, в ласки да замре.\\
Но вие, сълзи, стойте си при мене! --\\
Той няма да ви разбере.}

\vspace{0.3cm}

\noindent
\begin{tabular}{@{}p{0.48\textwidth} p{0.48\textwidth}@{}}
\textbf{А)} за природата & \textbf{В)} за труда \\
\textbf{Б)} за любовта & \textbf{Г)} за творчеството
\end{tabular}

\vspace{0.6cm}

\noindent
\textbf{31. Кое твърдение е вярно?}

\vspace{0.2cm}

\noindent
\textbf{А)} И в „Аз искам да те помня все така“, и в „Колко си хубава“ любовта е изобразена в своята преходност и нетряйност.
\\
\noindent
\textbf{Б)} За разлика от „Аз искам да те помня все така“, в „Колко си хубава“ любовта е споделено чувство между двама влюбени.
\\
\noindent
\textbf{В)} За разлика от „Колко си хубава“, в „Аз искам да те помня все така“ любовта е осмислена като блажено тайнство в света.
\\
\noindent
\textbf{Г)} И в „Аз искам да те помня все така“, и в „Колко си хубава“ е разкрита любовта като мост между минало, настояще и бъдеще.

\vspace{0.6cm}

\noindent
\textbf{32. Кое от тълкуванията съответства на смисъла на „Балкански синдром“?}

\vspace{0.2cm}

\noindent
\textbf{А)} Произведението внушава невъзможността на човека да победи пороците в бита.
\\
\noindent
\textbf{Б)} Творбата представя парадоксална тъжно-смешна картина на българските нрави.
\\
\noindent
\textbf{В)} Пиесата разкрива самотата на героите като единствена алтернатива на съществуването.
\\
\noindent
\textbf{Г)} Текстът представя липсата на любов като основна причина за страданието на героите.

\vspace{0.6cm}

\noindent
\textbf{33. В коя творба НЕ присъства мотивът за духовното пробуждане на народа?}

\vspace{0.2cm}

\noindent
\begin{tabular}{@{}p{0.48\textwidth} p{0.48\textwidth}@{}}
\textbf{А)} „Железният светилник“ & \textbf{В)} „Паисий“ \\
\textbf{Б)} „Борба“ & \textbf{Г)} „История“
\end{tabular}

\vspace{0.6cm}

\noindent
\textbf{34. Кои две творби утвърждават избора на свободолюбивата личност?}

\vspace{0.2cm}

\noindent
\textbf{А)} „Две души“ и „Ветрената мелница“
\\
\noindent
\textbf{Б)} „Песента на колелетата“ и „Вяра“
\\
\noindent
\textbf{В)} „Честен кръст“ и „Потомка“
\\
\noindent
\textbf{Г)} „Молитва“ и „Посвещение“

\vspace{0.6cm}

\noindent
\textbf{35. Съпоставете интерпретациите на темата за родното в „Бай Ганьо журналист“ и в „Железният светилник“ и запишете ЕДНА РАЗЛИКА между тях.}

\vspace{0.2cm}
\noindent\rule{\linewidth}{0.4pt}
\noindent\rule{\linewidth}{0.4pt}
\noindent\rule{\linewidth}{0.4pt}
\noindent\rule{\linewidth}{0.4pt}

\vspace{0.6cm}

\noindent
\textbf{36. Разтълкувайте заглавието „Ноев ковчег“ в контекста на творбата.}

\vspace{0.2cm}
\noindent\rule{\linewidth}{0.4pt}
\noindent\rule{\linewidth}{0.4pt}
\noindent\rule{\linewidth}{0.4pt}
\noindent\rule{\linewidth}{0.4pt}

\vspace{0.6cm}

\noindent
\textbf{37. Запишете в 2 -- 3 изречения каква е ролята на извънземното в комедията „Балкански синдром“.}

\vspace{0.2cm}
\noindent\rule{\linewidth}{0.4pt}
\noindent\rule{\linewidth}{0.4pt}
\noindent\rule{\linewidth}{0.4pt}
\noindent\rule{\linewidth}{0.4pt}

\vspace{0.6cm}

\noindent
\textbf{38. Запишете ЕДНО ПРЕНОСНО ЗНАЧЕНИЕ на блатото от цитирания откъс в контекста на темата за обществото и властта.}

\vspace{0.2cm}

\noindent
\textit{
-- Какво е това? -- запита съдията.\\
-- Блато, господин съдия... Пътят върви през него. То е плитко, хич не се бой. Тук-таме само има дупки... Колко пъти съм го минавал аз, и с кола, и пеши, и... Дий, дий, господа! Дръжте се здраво, господин съдия!\\
Конете навлязоха в студената вода, в която се оглеждаше небето, и предпазливо зацамбуркаха напред, като почакаха до потъват все по-дълбоко и по-дълбоко. Мъртвата и блестяща, зелено бисерна вода на блатото се размърда и оживя.\\
-- Стой, говедо! -- извика по едно време съдията и се изправи изплашено в кожуха си. -- Ще ме удавиш, шопе! Не видиш ли, че каруцата се напълни с вода? Спри! Спри!\\
Съдията, разгневен, почана да псува...\\
Андрешко спря конете, Каруцата, потънала до пода, застана сред блатото, краят на което се губеше в непроницаемата тъмнина.\\
}
\vspace{-0.2cm}

\hspace{11.5cm} (Елин Пелин, „Андрешко“)

\vspace{0.6cm}

\noindent
\textbf{39. Свържете заглавието на всяка от творбите с нейния автор, като срещу съответната буква запишете номера, съответстващ на името на автора.}

\vspace{0.3cm}

\noindent
\begin{tabular}{@{}p{0.55\textwidth} p{0.45\textwidth}@{}}
\textbf{А)} „При Рилския манастир“ & 1. Емилиян Станев \\
\textbf{Б)} „Крадецът на праскови“ & 2. Христо Смирненски \\
\textbf{В)} „Приказка за стълбата“ & 3. Иван Вазов
\end{tabular}

\vspace{0.6cm}

\noindent
\textbf{40. В текст до 5 изречения обяснете ролята на библейския символ в посочения цитат за смисъла на творбата „Честен кръст“ от Борис Христов.}

\vspace{0.3cm}

\noindent
\textit{... и бързам да напусна на поезията храма,\\
преди да запълзят змии по младото ми рамо.}

\vspace{0.2cm}
\noindent\rule{\linewidth}{0.4pt}
\noindent\rule{\linewidth}{0.4pt}
\noindent\rule{\linewidth}{0.4pt}
\noindent\rule{\linewidth}{0.4pt}
\noindent\rule{\linewidth}{0.4pt}
\noindent\rule{\linewidth}{0.4pt}
\noindent\rule{\linewidth}{0.4pt}
\noindent\rule{\linewidth}{0.4pt}

\newpage

\textbf{ЧАСТ 3 (Време за работа: 120 минути)}
\vspace{0.2cm}

\noindent
{\textit{\underline{Аргументативния текст напишете в листа за отговори на предвиденото място.}}}
\vspace{0.4cm}

\textbf{41.} Напишете аргументативен текст в обем до 4 страници по ЕДНА от следните две теми:
\vspace{0.2cm}
\textbf{Тема за есе:}
\vspace{0.1cm}
\begin{center}
\textbf{„Род и чест“}
\\
\textit{(Есе)}
\end{center}

\vspace{0.2cm}
\noindent
\textbf{Тема за интерпретативно съчинение:}
\vspace{0.1cm}
\begin{center}
\textbf{„Свободният избор и родовата памет“}
\\
\textit{(Интерпретативно съчинение върху „Потомка“ на Елисавета Багряна)}
\end{center}
\begin{flushleft}
Няма прародителски портрети,
\\
ни фамилна книга в моя род
\\
и не знам аз техните завети,
\\
техните лица, души, живот.
\vspace{0.5cm}
\\
Но усещам, в мене бие древна,
\\
скитническа, непокорна кръв.
\\
Тя от сън ме буди нощем гневно,
\\
тя ме води към греха ни пръв.
\vspace{0.5cm}
\\
Може би прабаба тъмноока,
\\
в свилени шалвари и тюрбан,
\\
е избягала в среднощ дълбока
\\
с някой чуждестранен, светъл хан.
\vspace{0.5cm}
\\
Конски тропот може би кънтял е
\\
из крайдунавските равнини
\\
и спасил е двама от кинжала
\\
вятърът, следите изравнил.
\vspace{0.5cm}
\\
Затова аз може би обичам
\\
необхватните с око поля,
\\
конски бяг под плясъка на бича,
\\
волен глас, по вятъра разлян.
\vspace{0.5cm}
\\
Може би съм грешна и коварна,
\\
може би средпът ще се сломя —
\\
аз съм само щерка твоя вярна,
\\
моя кръвна майчице-земя.
\end{flushleft}

\end{document}
